\documentclass[output=paper,colorlinks,citecolor=brown]{langscibook}
\ChapterDOI{10.5281/zenodo.22078577}

\author{Berit Gehrke\orcid{0000-0002-6315-9532}\affiliation{Humboldt-Universität zu Berlin}}
% replace the above with you and your coauthors
% rules for affiliation: If there's an official English version, use that (find out on the official website of the university); if not, use the original
% orcid doesn't appear printed; it's metainformation used for later indexing

%%% uncomment the following line if you are a single author or all authors have the same affiliation
\SetupAffiliations{mark style=none}

%% in case the running head with authors exceeds one line (which is the case in this example document), use one of the following methods to turn it into a single line; otherwise comment the line below out with % and ignore it
% \lehead{anonymous%Gehrke}
% \lehead{Radek Šimík et al.}

\title{Cross-Slavic aspect, passives, and temporal definiteness}
% replace the above with your paper title
%%% provide a shorter version of your title in case it doesn't fit a single line in the running head
% in this form: \title[short title]{full title}
\abstract{The chapter discusses cross-Slavic variation in aspect use in four different contexts (sequences of single events, habituality, historical present, factual imperfectives), with a particular focus on Czech vs. Russian. A new context is added to the mix, namely passives (participial and reflexive), and it is shown that Czech and Russian differ in that Czech, but not Russian, regularly derives both types of passives from both aspects, with predictable aspectual meanings. I summarise the few existing formal-semantic accounts of cross-Slavic aspectual variation, which all have in common that they apply some concept of definiteness to the verbal and/or temporal domain. Finally, I outline a general research programme to exploit parallels between individuals, events, and times regarding definiteness. %goes here and should not have more than 150 words.

\keywords{aspect, definiteness, passive, participle, reflexive}
}

\IfFileExists{../localcommands.tex}{
   \addbibresource{../localbibliography.bib}
   \usepackage{tabularx,multicol}
\usepackage{url}
\urlstyle{same}

 
\usepackage{langsci-optional}
\usepackage{langsci-lgr}
\usepackage{langsci-gb4e}

% ADDED BY VOLUME EDITORS

% \usepackage{langsci-osl} % macros specific to Open Slavic Linguistics; PLACED DIRECTLY IN localcommands.tex
\usepackage{siunitx}
\usepackage[linguistics]{forest}
% \usepackage{caption} %borik (?)
\usepackage{pifont} %geist

% 08_muellerreichau

% \usepackage{caption}
\usepackage{subcaption}
\usepackage{cleveref}

\usepackage{drs}
\usepackage{diagbox}

% 09_simik

\usepackage{multirow}
\usepackage{tabto}

% 10_stegovec

\usetikzlibrary{arrows,patterns}
\usepackage{listings}
% \usepackage{multirow}

% 01_gehrke-simik

% \usepackage{comment}

% 13_wagiel

% \usepackage{pifont} % for checkmarks (\ding{51}, \ding{55})


\usepackage{langsci-branding}

   % \newcommand*{\orcid}{}


\makeatletter
\let\thetitle\@title
\let\theauthor\@author
\makeatother

\newcommand{\togglepaper}[1][0]{
   \bibliography{../localbibliography}
   \papernote{\scriptsize\normalfont
     \theauthor.
     \titleTemp.
     To appear in:
     E. Di Tor \& Herr Rausgeberin (ed.).
     Booktitle in localcommands.tex.
     Berlin: Language Science Press. [preliminary page numbering]
   }
   \pagenumbering{roman}
   \setcounter{chapter}{#1}
   \addtocounter{chapter}{-1}
}

\newbool{bookcompile}
\booltrue{bookcompile}
\newcommand{\bookorchapter}[2]{\ifbool{bookcompile}{#1}{#2}}

\AtBeginDocument{\nogltOffset} %removes extra spacing before trans line in examples - SPECIFIC TO OPEN SLAVIC LINGUISTICS, PREVIOUSLY AGREED UPON


%%%%%%%%%%%%%%%%%%%%%%%%%%%%%%%%%%%%%%%%%%%%%%%%%%%%%%%%%%%%%%%%%%%%%
%%      File: langsci-osl.sty
%%    Author: Radek Simik and Language Science Press (http://langsci-press.org)
%%      Date: 2019-02-18
%%   Purpose: This file contains macros for the Open Slavic Linguistics at Language Science Press series and is included 
%%            into langscibook.cls
%%  Language: LaTeX
%%   Licence: 
%%%%%%%%%%%%%%%%%%%%%%%%%%%%%%%%%%%%%%%%%%%%%%%%%%%%%%%%%%%%%%%%%%%%%

% FIXING GRAY

\definecolor{lsDOIGray}{cmyk}{0,0,0,0.45}

% PACKAGES REQUIRED

\RequirePackage{relsize}
\RequirePackage{stmaryrd}
\RequirePackage{array}
\RequirePackage{tabularx}

% % KEYWORDS

% \newcommand{\keywords}[1]{\noindent\textbf{Keywords:} {#1}}

% SEMANTICS MACROS - GENERAL FOR ALL PAPERS

\newcommand{\sib}[1]{$\llbracket${#1}$\rrbracket$} % = semantic interpretation brackets; puts text into double square brackets
\newcommand{\stb}[1]{\langle{#1}\rangle} % = semantic type brackets; puts stuff into semantic type brackets; necessary to use in the form $\st{}$
\newcommand{\cnst}[1]{\textsf{\relsize{-1}\MakeUppercase{#1}}} %creates sans-serif small-caps, used for typesetting logical constants which have no other appropriate symbols (such as Greek letters)

% MINUS HSPACE MACRO

\newcommand{\minsp}[1]{#1\hspace{-2.5pt}} % use for non-glossed elements in glossed examples (typically stars, brackets, slashes, ...)

% CITATION MACROS

\newcommand{\citeposst}[1]{\citeauthor{#1}'s (\citeyear{#1})} % produces Chomsky's (1995)
\newcommand{\citeposstpg}[2]{\citeauthor{#1}'s (\citeyear[#2]{#1})} % produces Chomsky's (1995: page)
\newcommand{\citepossalt}[1]{\citeauthor{#1}'s \citeyear{#1}} % produces Chomsky's 1995
\newcommand{\citepossaltpg}[2]{\citeauthor{#1}'s \citeyear[#2]{#1}} % produces Chomsky's 1995: page

% BARPLOT

\usepackage{pgfplots}
\pgfplotsset{compat=1.18} 
\newcommand{\barplot}[5]{%
  \begin{tikzpicture}
    \begin{axis}[
	xlabel={#1},  
	ylabel={#2}, 
	axis lines*=left, 
        width  = {#3\textwidth},
	height = .3\textheight,
    	nodes near coords, 
	xtick=data,
	x tick label style={},  
	ymin=0,
	symbolic x coords={#4},
	]
	\addplot+[ybar,lsRichGreen!80!black,fill=lsRichGreen] plot coordinates {
	    #5
	}; 
    \end{axis} 
  \end{tikzpicture} 
}



%%%%%%%%%%%%%%%%%%%%%%%%%%%
%%% 04_filip %%%

% \newcommand{\mC}[1]{\multicolumn{1}{c}{#1}} % multicolumn with a single column; ALREADY DEFINED

\usetikzlibrary{arrows, arrows.meta}
\usetikzlibrary{positioning, decorations, 
                decorations.pathreplacing, calligraphy}

%%%%%%%%%%%%%%%%%%%%%%%%%
%%% 08_muellerreichau %%%

\captionsetup[subfigure]{subrefformat=simple,labelformat=simple}
\renewcommand\thesubfigure{(\alph{subfigure})}

%%%%%%%%%%%%%%%%%%%%%%%%%%%
%%% 09_simik

\newcommand{\citepossts}[1]{\citeauthor{#1}' (\citeyear{#1})} % produces Rawlins' (2008)
\newcommand{\citepossalts}[1]{\citeauthor{#1}' \citeyear{#1}} % produces Rawlins' 2008


%%%%%%%%%%%%%%%%%%%%%%%%%%%%%%%
%%% 10_stegovec

\newcommand{\poscite}[1]{\citeauthor{#1}'s (\citeyear{#1})}

%extra commands
\newcommand{\fst}{\textsc{1p}}
\newcommand{\snd}{\textsc{2p}}
\newcommand{\trd}{\textsc{3p}}
\newcommand{\refl}{\textsc{refl}}
\newcommand{\excl}{\textsc{excl}}
\newcommand{\incl}{\textsc{incl}}
\newcommand{\sg}{\textsc{sg}}
\newcommand{\pl}{\textsc{pl}}
\newcommand{\dl}{\textsc{dl}}
\newcommand{\fem}{\textsc{f}}
\newcommand{\masc}{\textsc{m}}
\newcommand{\neut}{\textsc{n}}
\newcommand{\ImpCl}{IMPERATIVE}


\newcommand{\tikzmark}[1]{\tikz[overlay,remember picture] \node (#1) {};}
\forestset{M/.style={
		for tree={s sep=5mm, inner sep=0.5pt, l=5mm,
			calign=fixed edge angles,
			calign primary angle=-65,
			calign secondary angle=65},
	},
	Mid/.style={
		for tree={s sep=0mm, inner sep=0pt, l=0mm,
			calign=fixed edge angles,
			calign primary angle=-62.5,
			calign secondary angle=62.5},
	},
	word/.style={before typesetting nodes={content=\emph{##1}},
		no edge,l=0,for parent={l sep=0}},
	feat/.style={before typesetting nodes={content={{[}##1{]}}},
		no edge,l=0,for parent={l sep=0}},
	feat2/.style={before typesetting nodes={content={##1}},
		no edge,l=0,for parent={l sep=0}},
	llap/.style={
		tikz+={%
			\edef\forest@temp{\noexpand\node[\option{node options},
				anchor=base east,at=(.base east)]}%
			\forest@temp{#1\phantom{\option{environment}}};
		}
	},
	rlap/.style={
		tikz+={%
			\edef\forest@temp{\noexpand\node[\option{node options},
				anchor=base west,at=(.base west)]}%
			\forest@temp{\phantom{\option{environment}}#1};
		}
	}
}

\makeatletter
\newcommand{\coloruline}[2]{%
	\newcommand\temp@reduline{\bgroup\markoverwith
		{\textcolor{#1}{\rule[-0.5ex]{2pt}{0.4pt}}}\ULon}%
	\temp@reduline{#2}%
}

\newcommand{\coloruuline}[2]{%
	\UL@protected\def\temp@uuline{\leavevmode \bgroup
		\UL@setULdepth
		\ifx\UL@on\UL@onin \advance\ULdepth2.8\p@\fi
		\markoverwith{\textcolor{#1}{\lower\ULdepth\hbox
				{\kern-.03em\vbox{\hrule width.2em\kern1\p@\hrule}\kern-.03em}}}%
		\ULon}%
	\temp@uuline{#2}%
}

\newcommand{\coloruwave}[2]{%
	\UL@protected\def\temp@uwave{\leavevmode \bgroup 
		\ifdim \ULdepth=\maxdimen \ULdepth 3.5\p@
		\else \advance\ULdepth2\p@ 
		\fi \markoverwith{\textcolor{#1}{\lower\ULdepth\hbox{\sixly \char58}}}\ULon}
	\font\sixly=lasy6 % does not re-load if already loaded, so no memory drain.
	\temp@uwave{#2}%
}
\makeatother

%%%%%%%%%%%%%%%%%
%%% 13_wagiel %%%

\newcommand\irregularcircle[2]{% radius, irregularity
  \pgfextra {\pgfmathsetmacro\len{(#1)+rand*(#2)}}
  +(0:\len pt)
  \foreach \a in {10,20,...,350}{
    \pgfextra {\pgfmathsetmacro\len{(#1)+rand*(#2)}}
    -- +(\a:\len pt)
  } -- cycle
}


% IF POSSIBLE, CAN WE USE THIS SETTING JUST FOR 13_wagiel?

% \forestset{
% default preamble={for tree={s sep=10mm, inner sep=0, l=0}},
% fairly nice empty nodes/.style={
%         delay={where content={}{shape=coordinate,
%               for siblings={anchor=north}}{}},
% for tree={s sep=4mm}},
% }

   %% hyphenation points for line breaks
%% Normally, automatic hyphenation in LaTeX is very good
%% If a word is mis-hyphenated, add it to this file
%%
%% add information to TeX file before \begin{document} with:
%% %% hyphenation points for line breaks
%% Normally, automatic hyphenation in LaTeX is very good
%% If a word is mis-hyphenated, add it to this file
%%
%% add information to TeX file before \begin{document} with:
%% %% hyphenation points for line breaks
%% Normally, automatic hyphenation in LaTeX is very good
%% If a word is mis-hyphenated, add it to this file
%%
%% add information to TeX file before \begin{document} with:
%% \include{localhyphenation}
\hyphenation{
    par-a-digm
    Cart-wright %filip
    Truswell %docekal
    Shcher-bi-na %simik
    Mün-chen %gehrke
    Mo-ni-ka %gehrke
    spi-sov-né %gehrke
    li-te-ra-tur-no-go %gehrke
    russ-koj %gehrke
    so-ver-šen-no-go %gehrke
    russ-kom %gehrke
    sla-vjan-sko-mu %gehrke
    Zeit-schrift %gehrke
    Na-ta-sha %gehrke
    Zu-stands-pas-siv %gehrke
    Um-gangs-spra-che %gehrke
    Bra-so-vea-nu %geist
}

\hyphenation{
    par-a-digm
    Cart-wright %filip
    Truswell %docekal
    Shcher-bi-na %simik
    Mün-chen %gehrke
    Mo-ni-ka %gehrke
    spi-sov-né %gehrke
    li-te-ra-tur-no-go %gehrke
    russ-koj %gehrke
    so-ver-šen-no-go %gehrke
    russ-kom %gehrke
    sla-vjan-sko-mu %gehrke
    Zeit-schrift %gehrke
    Na-ta-sha %gehrke
    Zu-stands-pas-siv %gehrke
    Um-gangs-spra-che %gehrke
    Bra-so-vea-nu %geist
}

\hyphenation{
    par-a-digm
    Cart-wright %filip
    Truswell %docekal
    Shcher-bi-na %simik
    Mün-chen %gehrke
    Mo-ni-ka %gehrke
    spi-sov-né %gehrke
    li-te-ra-tur-no-go %gehrke
    russ-koj %gehrke
    so-ver-šen-no-go %gehrke
    russ-kom %gehrke
    sla-vjan-sko-mu %gehrke
    Zeit-schrift %gehrke
    Na-ta-sha %gehrke
    Zu-stands-pas-siv %gehrke
    Um-gangs-spra-che %gehrke
    Bra-so-vea-nu %geist
}

   \boolfalse{bookcompile}
   \togglepaper[5]%%chapternumber
}{}


\begin{document}
\maketitle

\section{Introduction}
\label{sec:intro}

All Slavic languages have a grammatical category of aspect: each verb form is either imperfective (IPFV) or perfective (PFV) (in a given contexts), and in many environments one or the other aspect is obligatory. It is therefore commonly assumed that the lexical meaning of a given verb can be described using either aspect, so that for many lexical verb meanings we get an \textsc{aspectual pair} of IPFV and PFV forms.\footnote{I am abstracting away from so-called biaspectual verbs, which are interpreted as IPFV or PFV in a given context, such as Czech \textit{akceptovat} `accept' or Russian \textit{kaznit'} `execute' (see, e.g., \citealt{janda07} for discussion of Russian).} All Slavic languages use primarily prefixes and suffixes on the verb to make aspectual distinctions% \citep[see also][]{Mueller-Reichau2026, tatevosovslavsem}
.\footnote{\label{fn:aorimpf}In addition to the aspectual system described here (illustrated with examples from Russian and Czech), Bulgarian, Macedonian, Sorbian, and some Serbian dialects have aorist and imperfect forms, which make aspectual contributions that are orthogonal to the IPFV-PFV-distinction and similar to what we find in, e.g., Spanish and French \citep[see, e.g.,][]{deswart98}. To my knowledge, there is no formal-semantic account of the interaction between these aspectual tense forms and (I)PFV in these languages. The other Slavic languages lost these distinctions and only have one past tense form, which is historically related to the perfect participle (the so-called \textit{l}-participle) and combines with the auxiliary `be' in some languages (e.g. Czech) or is a stand-alone synthetic past tense form in others (e.g. Russian).} In \REF{ex:asppairs} we see typical examples of aspectual pairs from Czech and Russian, which are either derived by prefixing a morphologically \textsc{simple} IPFV to derive a PFV partner, or by suffixing a prefixed or a morphologically simple PFV to derive an IPFV partner (the latter forms are descriptively labeled \textsc{secondary imperfectives}, SIs).\footnote{I set aside less common morphological means to derive an aspectual partner, as well as suppletive forms. For similar data from some other Slavic languages, see, for instance, \citet{arsenijevic06,zaucerdiss,kwapidiss}.}

\ea\label{ex:asppairs}
\ea \textsc{ipfv} \textit{psát} $>$ \textsc{pfv} \textit{na-psat} `write'\\
    \textsc{pfv} \textit{pode-psat} $>$ \textsc{ipfv} \textit{pode-pis-ova-t} `sign'\\
    \textsc{pfv} \textit{dát} $>$ \textsc{ipfv} \textit{dá-va-t} `give' \hfill (Czech)
\ex \textsc{ipfv} \textit{pisat'} $>$ \textsc{pfv} \textit{na-pisat'} `write'\\
    \textsc{pfv} \textit{pod-pisat'} $>$ \textsc{ipfv} \textit{pod-pis-yva-t'} `sign'\\
    \textsc{pfv} \textit{dat'} $>$ \textsc{ipfv} \textit{da-va-t'} `give'\hfill (Russian)
\z
\z

\noindent Thus, there is no uniform morphological marking of either aspect: there are PFVs that contain a prefix or not (e.g. PFV `write' vs. PFV `give' in \REF{ex:asppairs}), as well as IPFVs that are morphologically simple (e.g. IPFV `write') or complex, as in the case of SIs (e.g. IPFV `sign'). Nevertheless, native speaker intuitions clearly group such verb forms into either the IPFV or the PFV group, and there are diagnostics to show which group a verb form belongs to. For example, in Czech and Russian, only IPFVs can form periphrastic future forms and can combine with phase verbs, as illustrated for Russian in \REF{ex:diag} (see, e.g., \citealt{Schoorlemmer1995,Filip1999,Borik2006,Gehrke2008} for discussion and further diagnostics).

\ea\label{ex:diag}
\gll    Ona \minsp{\{} budet / načala\} \minsp{\{} pisat' / podpisyvat' /\hspace{1.3cm} \minsp{*} napisat' / \minsp{*} podpisat'\} pis'mo.\\
    she.\textsc{nom} {} will {} began.\textsc{pfv} {} write.\textsc{ipfv.inf} {} under.write.\textsc{si.inf} {} {} on.write.\textsc{pfv.inf} {} {} under.write.\textsc{pfv.inf} letter.\textsc{acc}\\
\glt (Intended:) `She \{will / started to\} \{write / sign\} a/the letter.' \hfill (Russian)
%\ex
%\gll    Ona načala \minsp{\{} pisat' / podpisyvat' / *napisat' / *podpisat'\} pis'mo.\\
%    she.\textsc{nom} began.\textsc{pfv} {} write.\textsc{ipfv} {} sign.\textsc{si} {} write.\textsc{pfv} {} sign.\textsc{} letter.\textsc{acc}\\
%\glt (Intended:) `She started to \{write / sign\} a/the letter.' \hfill (Russian)
%\z 
\z 

\noindent Despite the overall morphological commonalities in the aspect systems of all Slavic languages, descriptive Slavicists early on noted cross-Slavic variation in aspect use, particularly between Czech and Russian \citep[][]{dokulil48%,dokulil53
,krizkova55,%krizkova58,krizkova61,
bares56,%bondarko58,
bondarko59,ivancev59/60,%sirokova63,
sirokova71}, and these two languages have been the focus of subsequent descriptive work (e.g. \citealt{eckert84,stunova93,petruxina00}, with the latter also taking into account Bulgarian, Polish, and Slovak). An investigation into the differences between ten Slavic languages in various contexts is provided by \citet{Dickey2000}, who proposes a cognitive semantic account with two contrasting aspect types, a Western type with its prototype Czech, and an Eastern type, with Russian as its prototype. Subsequently, additional corpus studies broadened the empirical picture \citep[][]{gehrke02,waldenfels12,waldenfels14,alvestad13,duebbersdiss,klimek22}. 

The descriptive observations about cross-Slavic differences in the use of (I)PFV are by now widely acknowledged, but only recently a few formal-semantic accounts have been proposed aiming at capturing the differences. These almost exclusively deal with differences in so-called general-factual contexts (see \sectref{sec:formsemasp}) \citep[][]{alvestad13,Mueller-Reichau2018a,klimek22}, with the exception of \citet{mueller23}, who includes descriptions of ongoing and habitual events. A general theory to account for the differences, be it semantic or pragmatic or a combination of both, is still absent.\footnote{Since this chapter is part of a volume on the semantics of Slavic languages, I mostly set aside potential (morpho-)syntactic considerations, although capturing the full picture of the variation should ultimately also take syntax into account.}  

The aim of this chapter is three-fold. First, after an outline of formal-semantic background assumptions in \sectref{sec:formsemasp}, \sectref{sec:crossslav} provides an overview of empirical investigations into cross-Slavic differences in the use of (I)PFVs in various contexts, with focus on Czech vs. Russian. Second, \sectref{sec:passives} adds an empirical domain that is missing almost entirely from the discussion, namely the domain of passives, for which we observe further aspectual differences between Czech and Russian.\footnote{I say ``almost'' because a recent corpus study \citep{wiemer+23} investigates the synchrony and diachrony of Polish and Russian past passive participles from the point of view of aspectual forms and readings.} Third, \sectref{sec:formal} summarises existing formal-semantic accounts of cross-Slavic aspectual variation, which all have in common that they employ some concept of temporal definiteness. \sectref{sec:definiteness} takes stock and outlines a general research programme to explore parallels between the nominal and verbal domain with respect to definiteness. 

\section{Formal-semantic background assumptions}
\label{sec:formsemasp}
\largerpage

This section briefly outlines formal-semantic background assumptions, relevant for later sections. These include the semantics of aspect (more detailed discussion can be found in \citetv{chapters/08_muellerreichau} and \citetv{chapters/11_tatevosov}; see also \citetv{chapters/07_gronn} on Slavic tense), from a cross-linguistic perspective and including cross-linguistic variation, the so-called general-factual use of Slavic IPFVs, as well as the notions of definiteness and specificity in the nominal domain (more detailed discussion can be found in \citetv{chapters/02_borik}). The latter will serve as a point of departure for the discussion of proposals that use the notion of temporal definiteness to account for cross-Slavic variation in aspect use. 

\subsection{The semantics of aspect and cross-linguistics variation}

Throughout the chapter, I assume that there is a distinction between different event types at the level of lexical or inner aspect (related to telicity, resultativity, stativity, change-of-state), on the one hand, and grammatical or outer aspect, which operates on top of event structures of various complexity (e.g. \citealt{Smith1991,Filip1999}; for further discussion see \citealt{Gehrke2008}).\footnote{I use the term \textsc{event} as a cover term for both events and states (or for what \citealt{Vendler1957} called states, activities, accomplishments, achievements)%, similar to the use of the term \textsc{eventuality} in \citet{bach81,bach86}
.} From a descriptive point of view, it has been observed that Slavic PFV forms appear in descriptions of ``completed'' events (however one might want to formally characterise the notion of completion), whereas IPFV forms appear in the description of ongoing or durative events, regular circumstances,\footnote{A special case of regular circumstances are generic statements; I will not address generic sentences in this chapter, but see \textcitetv{chapters/04_filip} for discussion.} and (at least in some Slavic languages) iterative and habitual events (e.g. \citealt{isacenko62,Filip1999,%borik02, 
Borik2006}, for Czech or Russian). Cross-linguistically, these are common meanings associated with grammatical aspect (see, e.g., \citealt{deo15} for more recent discussion), which is why it is standardly assumed that Slavic (I)PFVs are instances of grammatical aspect. On the other hand, there are (morphosyntactic) proposals in the literature that relate this distinction to inner aspect \citep[e.g.][]{arsenijevic06,lazorczyk10}, but at least for Czech and Russian it has been shown convincingly that semantically, (I)PFV is distinct from (a)telicity \citep[see, e.g.,][]{Filip1999, Borik2006, Gehrke2008}.\footnote{\label{fn:innerouter}A theoretical option that to my knowledge has not been explored formally could be that Slavic languages vary in this respect in that in some Slavic languages (I)PFV is an inner-aspectual distinction, whereas in others it is an outer-aspectual one (see \citealt{stunova93} for descriptive generalisations along this line, concerning Czech vs. Russian aspect; see also \citealt{breu00,scholze08}, who argue, informally, that the colloquial Upper Sorbian aspect system is based on ``terminativity''). For example, \citet{lazorczyk10} as a proponent of the inner-aspectual analysis brings forward data from Bulgarian to argue that the aorist-imperfect distinction (recall footnote \ref{fn:aorimpf}) belongs to the realm of outer aspect% \citep[see also][]{todorovicdiss}
, so that (I)PFV must contribute to the level of inner aspect. It could very well be that in the presence of aorist/imperfect (and possibly other contributing factors), this language uses (I)PFV differently than Slavic languages that lack the aorist/imperfect distinction. In the remainder of the chapter, I will mostly focus on Czech vs. Russian, for which I view the arguments for treating (I)PFV as related to grammatical aspect convincing. A separate issue, which I return to below, is whether (I)PFV forms directly or only indirectly correlate with grammatical aspect meaning (or syntax).} 

Cross-linguistically, it is common to analyse the semantics of (I)PFV as expressing a temporal relation between the event time (E) (or situation time; see \citealt{Klein1994,Klein1995}) and a reference time (R) (or topic/assertion time; see \citealt{Klein1994,Klein1995}), building on \citet{Reichenbach47}. More precisely, (I)PFV have been treated as aspectual operators that take as input a predicate of events and turn it into a predicate of times, by introducing a reference time that gets temporally related to the run time of the event (represented as $\tau(e)$, following \citealt{Krifka1998}) \REF{ex:graspsem}.

\ea\label{ex:graspsem}
\ea\label{ex:graspsema} 
\sib{PFV} = $\lambda P \lambda t \exists e[\tau(e) \subseteq t \wedge P(e)]$
\ex\label{ex:graspsemb} 
\sib{IPFV} = $\lambda P \lambda t \exists e[t \subseteq \tau(e) \wedge P(e)]$
\z
\z

\largerpage
\noindent With the PFV aspect \REF{ex:graspsema}, the run time of the event is included in the reference time $t$, leading to an external perspective on the event which can therefore be viewed in its entirety or totality; with the IPFV \REF{ex:graspsemb}, the reference time is included in the event time, leading to an internal perspective on the described event.\footnote{See, for example, \citet{kratzer98,ferreira2016} for further discussion. In \REF{ex:graspsem}, I use a simplified extensional representation; to fully grasp the semantics especially of IPFVs it might be necessary to take into account intensions and possible worlds, but I will abstract away from this. There are also modal accounts, in particular for the IPFV, going back to at least \citet{Dowty1979}. Furthermore, there are accounts in particular of Slavic (I)PFV in terms of partial vs. total events \citep[e.g.][]{Filip1999,altshuler14} or in terms of linking the PFV to maximality \citep[e.g.][]{Filip2008}; I view these accounts as special instances of the more general account in \REF{ex:graspsem}.}

\largerpage
Broadly speaking and over-simplifying to a certain extent, there are at least two ways in which the semantic account of a given aspect can vary, which have been used in explaining cross-linguistic variation in grammatical aspect.\footnote{A third way is proposed in \citet{Arregui.etal2014}, who argue that the variation in the interpretation of imperfectivity in Slavic, Romance, and Jê additionally derives from a difference in the modal bases employed by the IPFV operators in these languages. An interesting and sophisticated theoretical proposal of the cross-linguistic variation as to whether an IPFV operator in a language can express the (narrower) progressive reading (e.g. English) or (broader) imperfective reading (e.g. Gujarati) is proposed in \citet{deo15}.

There are also possible (morpho-)syntactic accounts that I will not address in detail here, since I will be mostly concerned with semantic approaches. For example, \citet[][]{biskup23adv}, following \citet{Tatevosov2011,Tatevosov2015b} (see below), assumes that grammatical aspect operators are silent operators not directly linked to verbal aspectual morphology, but their presence is determined by the syntactic operation Agree; variation in the ordering or merging of particular morphemes could then lead to differences in aspect behaviour, both across and within Slavic languages.} One point of variation concerns proper inclusion ($\subset$) vs. improper inclusion ($\subseteq$). For example, \citet{altshuler14} argues that the English progressive, an instance of IPFV, involves proper inclusion of the reference time in the event time, whereas the Russian IPFV involves improper inclusion, and that this accounts for empirical differences between the two languages in the use of IPFVs.

Another point of variation has to do with markedness: In the opposition between IPFV and PFV one aspect in a given language might be the marked, the other the unmarked member in the opposition. This could involve morphological markedness, or markedness could be seen as more abstract, concerning the (wider or narrower) meaning spectrum of a given form. For example, the English  progressive (IPFV) is taken to be the morphologically and also semantically marked aspect, whereas there is no morphological marking of the PFV; this leads to a wideheld assumption that, e.g., the simple past expresses a PFV semantics \citep[see, e.g.,][]{Smith1991,Klein1994}, albeit unmarked.\footnote{A question that arises in this context is why the same does not hold for the simple present.} Alternatively, one could view non-progressive forms as expressing no grammatical aspect (see, e.g., \citealt{deswart98} for such a view, as well as \citealt{minor+23} for recent psycholinguistic evidence that points in this direction), and a PFV semantics to come about simply due to the overall context as well as due to the input, e.g. whether the underlying predicate is eventive/stative, telic/atelic, etc.; the absence of IPFV marking would then only play a contributing factor (e.g. it could be an anti-presupposition effect), but it would not be the decisive factor. In this chapter, I side with the latter view. For Slavic languages, on the other hand, it is common to assume that the PFV is the (semantically) marked member of the opposition, so that IPFVs either involve the negation of the semantics of the PFV (see, e.g., \citealt{%borik02,
Borik2006} for such an account of Russian), or can express either PFV or IPFV semantics (e.g. \citealt{gronn15} for Russian, whose account is discussed in more detail in \sectref{sec:gronn}).\footnote{Morphologically, it is clear from the examples in \REF{ex:asppairs} that one cannot take either aspect as marked/unmarked: in some cases the PFV is morphologically more complex than its PFV partner, in others it is the other way around \citep[see also discussion in][]{%jakobson32,jakobson56,
jakobson66}.}

This latter point also relates to the general question whether a given form (e.g. Slavic (I)PFV forms) always comes with the same semantics, or whether there can be form-meaning mismatches. Given the markedness assumptions outlined above, it is common to assume (at least implicitly) that for the English progressive and the Slavic PFV the forms involved correlate respectively with an IPFV and PFV semantics (but see \citealt{paslawskastechukr} for a general account of the Ukrainian tense-aspect system that fully dissociates form and meaning). For Slavic IPFVs, in contrast, there are authors that explicitly argue that IPFV forms do not necessarily correspond to an IPFV semantics \citep[for Russian, see][]{%borik02,
Borik2006,%borik18,
gronndiss,gronn15}, and others that (implicitly or explicitly) assume a form-meaning-correspondence (see e.g. \citealt{Filip1999,gehrke22,gehrketrueipf} for Czech and Russian, and \citealt{altshuler14} for Russian). % see \citep[for Russian, see][]{%Altshuler2012,altshuler13,
%altshuler14,gehrke22,gehrketrueipf}.
%Furthermore, there is the question of whether a given morphological form corresponds directly to a given aspect semantics. 
\citeauthor{Tatevosov2011} (\citeyear{Tatevosov2011}, \citeyear{Tatevosov2015b}, \citeyear{chapters/11_tatevosov} [this volume]), in turn, argues that there is no direct correlation between form and meaning, building on insights from \citet{Klein1994,Klein1995}: What we call (I)PFV forms (in Russian, but presumably in Slavic more generally) does not directly correspond to (I)PFV semantic operators, given that aspectual morphology appears directly on the verb, whereas grammatical aspect is located above the VP. Nothing I say here is in principle incompatible with this view, to which I am fully sympathetic; once we reach the level of grammatical aspect, a given IPFV form could still uniformly lead to an IPFV semantics.\footnote{For example, building on Tatevosov's work, \textcitetv{chapters/08_muellerreichau} proposes a system that derive Russian aspect semantics from the morphological input and the overall context.} Throughout the chapter, I use (I)PFV for the forms that appear in contexts that convey a typical (I)PFV semantics. 

There is a particular use of Slavic IPFVs that at first sight seems problematic for accounts that directly relate (I)PFV forms with an (I)PFV semantics (in a given sentence). This use has first been described for Bulgarian and Russian by \citet{maslov59}, who termed it \textsc{general-factual}. It is peculiar from a cross-linguistic perspective because factual IPFVs are found for the description of seemingly ``completed'' events (among others), in which one would expect PFV forms but for some reason the IPFV is still used. Subsequently there has been a lot of work on in particular Russian factual IPFVs, both from a descriptive and a formal perspective \citep[e.g.][]{glovinskaja81,glovinskaja89,paduceva96,mehlig01,mehlig13,gronndiss,gronn15,muellerhab,gehrketrueipf}, but to this date there is no consensus as to the precise analysis. Since this IPFV use has played a pivotal role in semantic accounts of (in particular Russian) aspect and is a prominent point of variation within the Slavic languages, the following section is devoted to it.

\subsection{The (general-)factual IPFV}
\label{sec:OF}

%\citeposst{gronndiss} account, which I will outline in \sectref{sec:gronn}, is taken as a point of departure for two proposals dealing with differences between Czech and Russian, among others \citep[][]{alvestad13,Mueller-Reichau2018a,mueller23}, to which I will turn afterwards. First, however, I will discuss \citeposst{Dickey2000} cognitive semantic proposal that Russian PFVs require temporal definiteness whereas Czech does not, an idea that is further explored more formally in \citet{Mueller-Reichau2018a,mueller23} and \citet{klimek22}.  
%As noted in \sectref{sec:formsemasp}, there is an IPFV use that is rather peculiar from a cross-linguistic perspective, and it is possible that it is not attested outside Slavic. This is the so-called \textsc{(general-)factual} use, first observed by \citet{maslov59} for Bulgarian and Russian (Russian \textit{obščefaktičeskoe}). 
In factual contexts, IPFVs can appear to describe bounded, ``completed'' events.\footnote{Again, the term ``event completion'' remains at an intuitive level. The traditional literature also discerns subtypes of the general-factual with intuitively non-completed events \citep[e.g.][]{glovinskaja81,glovinskaja89,paduceva96}. In formal accounts of factual IPFVs \citep[e.g.][]{gronndiss}, these subtypes are usually not addressed, leading to the somewhat distorted impression that general-factual contexts always involve ``completed'' events (see \citealt{gehrketrueipf} for further discussion).} In such contexts, it is often assumed that the IPFV is in \textsc{aspectual competition} with the PFV \citep[a term that goes back to at least][]{mathesius38}, because both can (often) be used interchangeably with only subtle differences that are hard to pin down. This section first describes two commonly assumed subtypes of factuals (existential, presuppositional); since most of the literature on factual IPFVs deals with Russian, I use Russian examples to illustrate. I then outline a semantic proposal for Russian factual IPFVs (\citealt{gronndiss} et seq.), which has been used in accounts of cross-Slavic aspect variation, to which I return in \sectref{sec:formal}. 

\subsubsection{Existential and presuppositional factuals}

The literature on Russian aspect distinguishes at least two subtypes of factual IPFVs, the existential type \citep[][]{paduceva96, gronndiss} and what \citeauthor{gronndiss} calls the presuppositional type \citep[the ``actional'' type in][]{paduceva96}. The \textsc{existential IPFV} is illustrated in \REF{ex:existential} \citep[corpus example from][180; my glosses]{gronndiss}.

\ea\label{ex:existential} 
\gll Ne bylo somnenij, čto ja prežde vstrečal ee.\\
not was.\textsc{ipfv.3sg.n} doubts.\textsc{gen} that I before met.\textsc{si.sg.m} her\\
\glt `There was no doubt that I had met her before.' \hfill (Russian)
\z

\noindent In this example, the (male)\footnote{Czech and Russian past tense forms display gender and number agreement with the subject (e.g. in the singular \textit{-$\varnothing$} for masculine, \textit{-a} for feminine, \textit{-o} for neuter), because they originate from a participle. I add this to the glosses only where relevant.} speaker states that he had a meeting with a female person in the past. From the context it is clear that a meeting was ``completed'' and happened in the past; nevertheless the IPFV is used. Existential IPFVs usually involve stress on the verb form and can be paraphrased as `There has been/is/will be (at least) one event of this type.' \citep[following the idea that existential IPFVs involve event types or kinds; see][]{mehlig01, mehlig13, %muellerkrat, muellerPI, 
gehrkemueller}. The paraphrase for \REF{ex:existential} would therefore be `There has been at least one event of the type ``meet her'''. The exact time when this event happened, and also whether it happened more than once, remains unspecified. In some existential contexts, Russian has to use the IPFV and cannot switch to the PFV, for example in combination with the temporal adverb \textit{kogda-nibud'} `ever', as illustrated in \REF{ex:kogda-nibud'} \citep[adapted from][104]{Dickey2000}.

\ea\label{ex:kogda-nibud'}
\gll Ty kogda-nibud' \minsp{\{*} prygnul / prygal\} s tramplina?\\
    you.\textsc{nom} ever {} jumped.\textsc{pfv} {} jumped.\textsc{ipfv} off diving.board\\
\glt `Have you ever jumped off a diving board?' \hfill (Russian) 
\z 

\noindent In \sectref{sec:defspec}, we will see that \textit{nibud'}-marked indefinite pronouns express scopal non-specificity, and I will get back to such adverbs in \sectref{sec:OFCZRU} and \sectref{sec:definiteness}. What examples like these show is that the general impression about the interchangeability of IPFV and PFV in factual (in this case existential) contexts that is sometimes created by the literature should be taken with a grain of salt \citep[see also][]{duebbersdiss}.

The \textsc{presuppositional IPFV} is illustrated in \REF{ex:presuppositional} \citep[from][108; my glosses and translation]{glovinskaja81}.

\ea\label{ex:presuppositional}
\gll Zimnij Dvorec stroil Rastrelli.\\
winter.\textsc{adj.acc} palace.\textsc{acc} built.\textsc{ipfv} Rastrelli.\textsc{nom}\\
\glt `It was Rastrelli who built the Winter Palace.' \hfill (Russian)
\z

\noindent In this example, we have a ``completed'' event in the past, the building of the Winter Palace (which now houses the Hermitage Museum in Saint Petersburg). It is known that this event happened exactly once and also when it happened, so there is no temporal non-specificity, unlike what we have with the existential IPFV. \citet{gronndiss} labels this IPFV use presuppositional because it arises when the existence of the event in question is given in or derivable from the context, hence it is presupposed (and in \citeauthor{gronndiss}'s account backgrounded, following \cite{geurtssandt97}). An utterance with a presuppositional IPFV adds further information about this presupposed event, and this new information is in focus. A suitable paraphrase is `The (already mentioned or contextually retrievable) event was/is/will be such and such.' In \REF{ex:presuppositional}, context or world knowledge presupposes the existence of the event `build (the) Winter Palace' (as objects of verbs of creation tend to do, but it could also have been talked about in the previous context); the new information is that the architect of the building was Rastrelli, which is where the focus lies. This IPFV use is usually accompanied by a particular information structure (see also \citealt{Borik.Gehrke2018} for further discussion); in our examples the presupposed (backgrounded) material appears sentence-initially and unstressed (the building of the Winter Palace) and the new information in focus is Rastrelli, in sentence-final position, resulting in a non-canonical OVS order.\footnote{Russian, like other Slavic languages, is a ``free word order'' language with canonical SVO order; deviations from this canonical order are primarily information-structurally motivated (see \citealt{jasinsimik} for discussion).} 

There is a common assumption in the literature that factual IPFVs do not just occur in Russian but also in other Slavic languages, and this is independent of whether the authors in question assume a single factual meaning or several subtypes; nevertheless, its use in Russian is reported to be more frequent than in, e.g., Czech \citep[e.g.][]{Dickey2000,alvestad13%,alvestad14
,duebbersdiss,Mueller-Reichau2018a,klimek22}%, as can be seen from \citeposst{dickey15} tables \ref{tab:xslavEASTWEST} and \ref{tab:xslavTRANS}
, as we will see in \sectref{sec:OFCZRU}. In contrast, it is argued in \citet{gehrke22} that Czech only makes use of presuppositional, but not of existential IPFVs; I will come back to this in \sectref{sec:OFgehrke}. Let us turn to \citeauthor{gronndiss}'s (\citeyear{gronndiss} et seq.) account% of Russian factual IPFVs and the semantics of (I)PFVs more generally
. 

\subsubsection{``Fake'' IPFVs and the Aspect Neutralisation Rule in Russian}
\label{sec:gronn}

Following the traditional view that the Russian IPFV is semantically unmarked, \citet{gronndiss} employs a weak IPFV semantics, which merely requires overlap between the event time and the reference time ($e \bigcirc t$) \citep[building on][]{Klein1995}. This rather weak semantics gets pragmatically strengthened to a ``true'' IPFV (the reference time is part of the event time), or to an actual PFV semantics (the event time is part of the reference time), which, he argues, we find with factual IPFVs. \citeauthor{gronndiss} takes into account the role of information structure to characterise the contexts in which strengthening happens in one or the other direction. This account is a precursor to \citet{gronn15}, in which it is proposed that IPFV forms can express both IPFV and PFV semantics, as in \REF{ex:gronn15PFIPF}.

\ea\label{ex:gronn15PFIPF} 
\ea \sib{PFV} = $\lambda t \lambda e. e \subseteq t$
\ex \sib{IPFV\textsubscript{ongoing}} = $\lambda t \lambda e. t \subseteq e$
\ex\label{ex:fakeipf}
\sib{IPFV\textsubscript{factual}} = $\lambda t \lambda e. e \subseteq t$ \hspace{1em} ``Fake'' IPFV
\z 
\z 

\noindent \citeauthor{gronn15} calls the IPFV that has the same semantics as the PFV in \REF{ex:fakeipf} a \textsc{``fake'' IPFV}. He proposes that the existence of IPFV\textsubscript{factual} alongside the PFV leads to aspectual competition, and also here he follows the Slavistic tradition. In the default case the PFV is argued to appear but in certain contexts the IPFV\textsubscript{factual} wins the competition. This is proposed to give rise to the presuppositional IPFV in cases where narrative progression is to be avoided (under the assumption that the PFV always leads to narrative progression; see below), and to the existential IPFV when the reference time is too large for the PFV semantics to be informative. 

Starting from the well-known observation that tenses can be both quantificational and referential, \citet{gronnstechow16} spell out a research programme that draws parallels between events and times, on the one hand, and the semantics of bare nominals in articleless languages (e.g. Russian), on the other. In particular, they treat tenses (R-S relations), aspects (E-R relations), and verbal predicates as relational predicates. Further information about times and events, given for example by adverbials, is added via predicate modification. Covert definite and indefinite operators turn events and times into dynamic generalised quantifiers, which, respectively, are anaphoric to a previous referent, maximally presupposing given information, or introduce a new referent. Following \citet{gronndiss,gronn15}, Russian IPFV forms have either an IPFV or a PFV (``fake'' IPFV) semantics. Covertly, on top, we get a definite or an indefinite event. Indefinite, complete events, on the other hand, are regularly referred to by PFVs, to ensure narrative progression, which is analysed as a pragmatic effect (``be orderly'').

Employing this system, \citet{gronn15} proposes that presuppositional IPFVs involve a definite event and a definite reference time, whereas existential IPFVs come with an indefinite event and an indefinite reference time. In order to account for why the IPFV is used in presuppositional contexts, \citet{gronnstechow16} propose the rule in \REF{ex:aspneugst}, which builds on the general idea that the IPFV is the unmarked member of the aspectual opposition.

\eanoraggedright \textsc{Aspect Neutralisation Rule} \citep[see][]{gronnstechow16}\label{ex:aspneugst}\\
When a semantically PFV aspect is definite/anaphoric, it is morphologically neutralised to IPFV. 
\z 

\noindent This system is the point of departure for \citeposst{alvestad13} account of cross-Slavic differences in aspect use in imperatives, which I will outline in \sectref{sec:alvestad}. The proposal to treat factual IPFVs as ``fake" and not ``true'' IPFVs has been challenged by \citet{gehrketrueipf}, and I return to this alternative account in \sectref{sec:OFgehrke}. Finally, \sectref{sec:definiteness} draws more general parallels between the nominal and verbal domain with respect to definiteness. In the following, let us move to formal-semantic assumptions about definiteness and specificity in the nominal domain. 

\subsection{Definiteness and specificity}
\label{sec:defspec}

This section briefly summarises semantic approaches to definiteness in the nominal domain (for more detailed discussion see \citetv{chapters/02_borik}), in as far as it is relevant for the notion of ``temporal definiteness'' that will play a role in later sections (see also \sectref{sec:gronn}). (In)definites appear both in argument as well as predicative position, and it is common to take the argument use as basic and to derive the predicative use from it. A widespread approach to singular definites takes them to presuppose uniqueness and existence \citep[following][]{frege92,strawson50}, while plural definites presuppose maximality \citep[following][]{sharvy80}. For example, \citet{heim11} spells out the semantics of (in)definiteness (not necessarily of \textit{the/a}) in \REF{ex:heimsemindef}, generalising maximality to both singular and plural definites. 

\ea\label{ex:heimsemindef}
\ea \sib{$+\textsc{def}$} = $\lambda P:\exists x \forall y [\textsc{max}(P)(y) \leftrightarrow x = y].\iota x.\textsc{max}(P)(x)$
\ex \sib{$-\textsc{def}$} = $\lambda P.\lambda Q.\exists x[P(x) \wedge Q(x)]$
\z 
\z 

\noindent \citeauthor{heim11} proposes that indefinites and definites form a scale, with indefinites being logically weaker. She argues that in languages that have a definite article, such as English, the indefinite cannot be used in definite contexts, due to an anti-uniqueness implication with indefinites, which arises in competition with definites, assuming the principle of Maximise Presupposition. 

\citet{coppockbeaver15}, in turn, take the predicative use as basic and propose that (singular) definites presuppose uniqueness but not existence; existence only comes in through a covert type shift of definites in argument position. Yet in other approaches definiteness is understood in terms of familiarity: (in)definites are associated with discourse referents, and with definites the referent is anaphorically linked to a previously introduced discourse referent (\citealt{kamp81,Heim1982}; see also \citealt{gronnstechow16}, as outlined in \sectref{sec:gronn}). Finally, there is also a distinction between strong definites, which satisfy uniqueness, and weak definites (e.g. \textit{go to the supermarket}), which do not and which, e.g., \citet{aguilarzwarts11} treat as kind-referring.

In addition, there is another notion that is relevant, namely specificity (see also \citetv{chapters/06_geist}). Also this notion has different uses in the literature. For example, \citet{geist10a} distinguishes between epistemic specificity (whether or not the speaker can identify the referent) and scopal specificity (whether an indefinite has wide or narrow scope with respect to other quantifiers or operators). She shows that Russian, which lacks articles, has several series of indefinite pronouns that are specialised to signal epistemic specificity (the \textit{koe-}series, as opposed to the epistemically non-specific \textit{-to}-series) or scopal non-specificity (the \textit{-nibud'}-series). The particular morphemes are added to \textit{wh}-items (e.g. \textit{kto} `who') to derive an indefinite pronoun (\textit{koe-kto, kto-to, kto-nibud'} `(different types of) some/any\-one'), and they are not limited to the nominal domain, as we already saw with \textit{kogda-nibud'} `ever' in \REF{ex:kogda-nibud'}, which is based on the \textit{wh}-word `when'.

All of the notions discussed in this section, apart from weak definites, have been transferred from the nominal to the verbal or sentential domain, and they will come up again in one or the other account of cross-Slavic variation in aspect use, summarised in \sectref{sec:formal}; some have already been addressed in \sectref{sec:gronn}. With these theoretical background assumptions in hand, let us turn to the empirical description of cross-Slavic variation in aspect use in selected contexts.

\section{Cross-Slavic variation in aspect use}
\label{sec:crossslav}

This section describes empirical generalisations about cross-Slavic variation in the use of aspect, focusing on Czech and Russian and on particular contexts for illustration: sequences of single events, habitual events, the historical present, and factual contexts. %In the second part (\sectref{sec:formal}), I will address \citeposst{Dickey2000} proposal of a Slavic East-West-isogloss in more detail to then turn to the few existing formal accounts addressing some of the differences. 
%
%\subsection{Empirical generalisations}
%\label{sec:empirical}
%
First, though, let me start with a quantitative picture that emerges from the literature on differences in aspect use between Czech and Russian. Table \ref{unterschiede} summarises the findings for some contexts.\footnote{Sequences of single events are discussed in \citet{ivancev59/60,eckert84,monnesland1984,%stunova88,stunova92,
stunova93,Dickey2000,petruxina00,gehrke02,gehrke22,barentsen08,fortuin2015typology}, iterativity and habituality in \citet{%sirokova63,petruxina78,petruxina00,
eckert84,%eckert85,stunova86,
stunova93,Dickey2000,kresin00,gehrke02,gehrke22,duebbersdiss,fortuin2015typology}, the historical present in \citet{krizkova55,bondarko58,bondarko59,petruxina83,stunova93,%stunova94,
Dickey2000,fortuin2015typology}, running instructions and commentaries in \citet{Dickey2000}. 

Further differences between Czech and Russian (and sometimes also other Slavic languages) have been described for imperatives \citep[][]{dokulil48,eckert84,benacchio10,waldenfels12,waldenfels14,alvestad13%,alvestad14
}, motion verbs \citep[][]{eckert91,gehrke02,gehrke22}, various prefixes and suffixes \citep[][]{nuebler92,petruxina00,dickey01,dickey05,dickey11,dickeyhutch}, contexts involving negation \citep[][]{dickeykresin09,waldenfels14,duebbersdiss}, nominalisations \citep[][]{Dickey2000,biskup23adv,gehrkeadv23}, as well as factual contexts \citep[][]{Dickey2000,gehrke02,gehrke22,alvestad13,duebbersdiss,fortuin2015typology,Mueller-Reichau2018a,klimek22}; I will address the latter in more detail in \sectref{sec:OF}.}

\begin{table}
	\caption{Some aspectual differences between Czech and Russian}
	\label{unterschiede}
%\setlength{\tabcolsep}{8.2mm}
		\begin{tabular}{lll}
	  \lsptoprule
			&\textsc{Czech} 	&\textsc{Russian}  \\ \midrule[0.25ex] 
				Chains of single events (past)	&IPFV, PFV &(almost excl.) PFV\\[0.25ex]
				Iterativity, habituality (past, present)	&IPFV, PFV&(almost excl.) IPFV 	
				\\[0.25ex]
				Historical present	&IPFV, PFV&(almost excl.) IPFV	 \\[0.25ex]
 			Running instructions \& commentaries	&IPFV, PFV&(almost excl.) IPFV	 \\
      \lspbottomrule
	\end{tabular}
\end{table}

The first, more coarse-grained picture that emerges from this table is that there are contexts in which Czech allows for the use of both aspects, whereas Russian uses just one aspect, and except for the first context it is the IPFV that is quasi-obligatory in Russian.\footnote{The motivation for adding ``almost excl(usively)'' in Table \ref{unterschiede} will be addressed in \sectref{sec:chains}--\sectref{sec:HP}.} Such differences have therefore also been framed in terms of the obligatory use of a particular aspect in Russian, vs. the optional use of a particular aspect in Czech, in a particular context \citep[e.g.][]{bondarko59,krizkova61,sirokova71,petruxina00}. However, it is not clear what it means for a grammatical aspect to be optional, as this suggests some kind of arbitrariness, or at least that aspect use in Czech is just a matter of choice, which is usually found with lexical, but not with grammatical categories. In other words, at this coarse-grained level, where we simply count (I)PFV forms in particular contexts (which is also something one could do statistically, see, e.g., \citealt{waldenfels12, waldenfels14, duebbersdiss,klimek22}), noting these differences is of course important, but a quantitative analysis cannot be the endpoint to understanding the differences; we have to take it as the starting point for a detailed qualitative analysis. In the following, I focus on differences in four particular contexts, each time trying to explore the circumstances which motivate the use of one or the other aspect, and thus to give a positive characterisation of the reasons for the occurrence of a given aspect form. I will pay special attention to the verb type employed, the importance of which has also been stressed in some of the previous empirical research \citep[e.g.][]{eckert84,stunova93,gehrke02}. 

\subsection{Chains of single events}
\label{sec:chains}

In the descriptions of chains of single events, Russian almost obligatorily uses the PFV to give rise to a sequence of event (SOE) interpretation, whereas we find both aspects in Czech \citep[][]{ivancev59/60,eckert84,monnesland1984,%stunova88,stunova92,
stunova93,Dickey2000,petruxina00,gehrke02,gehrke22,barentsen08,fortuin2015typology}. For example, if the last event in a chain has a clear temporal onset but then evolves further, we have a case that is commonly labeled \textsc{ingressivity}. In such cases, it is plausible that there is tension between using the PFV to mark the temporal onset and to explicitly mark SOE, on the one hand, and the IPFV, on the other, to mark the process of the evolving event. It turns out that Russian consistently employs the first strategy (PFV), leaving the evolution of the process implicit, while Czech regularly goes for the second strategy (IPFV), so that in this language ingressivity is contextually derived, but not marked on the verb form. This systematic difference in aspect use was first noted by \citet{ivancev59/60}, and one of his examples is given in \REF{ex:ivancevingr} 
\citep[from][36; Czech original by Božena Němcová; my own glosses and translations; relevant verb forms (here and in the following longer examples) italicised]{ivancev59/60}.

\ea\label{ex:ivancevingr}
\ea 
\gll {...} zvolna si \textit{sedl} vedle mne a Josefa, \textit{položil} hlavu do dlaně a \textit{díval} na mne.\\
    {} slowly \textsc{refl} down.sat.\textsc{pfv} next.to me and Josef put.\textsc{pfv} head.\textsc{acc} to palm and looked.\textsc{ipfv} on me\\
\glt    `He slowly sat down next to me and Josef, put his head in his palm and looked at me.' \hfill (Czech) 
\ex 
\gll {...} on tixo \textit{sel} vozle menja i Iozefa, \textit{sklonil} golovu na ruki i  \textit{stal} \textit{smotret'} na menja.\\
    {} he quietly down.sat.\textsc{pfv} near me and Josef tilted.\textsc{pfv} head.\textsc{acc} on hands and began.\textsc{pfv} watch.\textsc{ipfv.inf} on me\\
\glt    `He quietly sat down near me and Josef, put his head on his hands and started watching me.' \hfill (Russian)\label{ex:stal}
\z
\z 

\largerpage
\noindent 
In this example the last event in the chain of events is described by an activity predicate, `look/watch'. In Czech, a simple IPFV is used and the circumstance that the watching follows the preceding events (which, in turn, are described by PFVs in both languages) is understood only contextually. In Russian, in contrast, the translator uses the PFV phase verb \textit{stat'} `begin' in combination with the (IPFV) infinitive of `watch', and thus explicitly signals the SOE reading. 

In other examples of this sort, one can find the so-called ingressive prefix \textit{za-},\linebreak which marks the onset of an otherwise unbounded event \citep[see][]{isacenko62}. This prefix (in this use) is productively used in Russian but arguably non-existent in Czech \citep[see also][]{ivancev59/60,Dickey2000,petruxina00}.\footnote{As pointed out by Petr Biskup (p.c.), some examples are mentioned in the literature that could be used to argue for the existence of ingressive \textit{za-} in Czech, such as \textit{zelenat se} -- \textit{za-zelenat se} `be green -- suddenly become green' (\citealt{nadenicek11}; translated from German). There are several reasons for arguing against the view that these are truly ingressive prefixes, rather than taking them to be homophonous. First, the meaning expressed with ingressives is one of a process evolving, whereas Czech \textit{za}-verbs of the type above express inchoativity (a sudden onset) (see also \citealt{isacenko62} for a discussion of ingressivity vs. inchoativity). Second, truly ingressive \textit{za-}verbs in Russian do not derive secondary IPFVs. However, there are SIs of Czech inchoative \textit{za-}verbs (e.g. \textit{za-zelená-va-t} `to start to become green'); similar inchoative aspectual pairs are found in Russian (e.g. \textit{za-bole(-va)-t'} `become sick'). Third, in parallel corpora there is usually no Czech \textit{za-}correspondence to Russian ingressive \textit{za-} (see, e.g. \citealt{petruxina00,gehrke02}). I therefore maintain that there is no ingressive \textit{za-} in Czech (see also \citealt{gehrke22} for further discussion).} One such case is illustrated with the Russian original from Mixail Bulgakov's \textit{Master i Margarita} and its Czech translation in \REF{ex:po} \citep[][21]{gehrke22} \citep[see also][]{gehrke02}.

\ea \label{ex:po}
\ea \label{ex:poRU} 
\gll On \textit{pomolčal}          nekotoroe vremja v smjatenii, \textit{vsmatrivajas’} v lunu, plyvuščuju za rešetkoj, i    {\textit{zagovoril}: […]}\\
    he   \textsc{po}.was.silent.\textsc{pfv}  some time	     in confusion   watch.\textsc{ipfv.ap}      in moon.\textsc{acc} swimming.\textsc{acc} behind bars and \textsc{za}.spoke.\textsc{pfv}\\
    \glt `He was silent for a while, confused, watching the moon swimming behind the bars, and said: ...'\hfill (Russian)
\ex 
\gll Chvíli     zaraženě \textit{mlčel},           \textit{sledoval}      plující měsíc               za mříží, a    pak    se    {\textit{zeptal}: […]}\\
 while.\textsc{acc} confused.\textsc{adv}    was.silent.\textsc{ipfv} followed.\textsc{ipfv} swimming.\textsc{acc} moon.\textsc{acc} behind bars and then     \textsc{refl} inquired.\textsc{pfv}\\
    \glt`He was silent for a while, confused, followed the moon swimming behind the bars, and then inquired: ...' \hfill (Czech)\label{ex:poCZ}
\z 
\z 

\noindent In the Russian example, the last event in the chain is described by a PFV verb of saying, containing ingressive \textit{za-}. The Czech translator did not opt for a direct translation but chose a different verbal predicate, `inquire'. The example illustrates further systematic differences between Czech and Russian in the context of chains of events, as argued by \citet{gehrke22}: Whereas Russian uses finite PFV forms for a SOE reading, Czech can also employ IPFV forms, in particular in descriptions of states and activities, or of accomplishments of a longer duration. Russian uses PFVs even with activity and state predicates. For example, PFV `be silent' in \REF{ex:poRU} contains the delimitative prefix \textit{po-}, which temporally bounds an otherwise unbounded activity \citep[see][]{isacenko62}. Whereas Russian delimitative \textit{po-} is quite productive, Czech \textit{po-} retains its spatial meaning `a bit' and is usually not found in the same contexts \citep[see, e.g.,][]{gehrke02,gehrke22,dickeyhutch,dickey11}.\footnote{Some examples of purely temporally delimiting \textit{po-} are also discussed in the literature on Czech (e.g. \textit{po-přemýšlel} `thought.\textsc{pfv} (for a while)' and \textit{po-hovořil} `chatted.\textsc{pfv} (for a while)' in \citealt{souckovapo}; see also \citealt{biskup24}) but they are far less frequent than in Russian. In the corpus data discussed in \citet{gehrke02}, for instance, there was not a single instance of delimitative \textit{po-} in Czech, but an abundance of Russian delimitative \textit{po-}verbs.} The flip side to this is that Russian uses non-finite verb forms, in this case the adverbial participle (AP) \textit{vsmatrivajas’} `(lit.) in-watching', to avoid a SOE reading, whereas Czech uses finite verb forms throughout. Thus, SOE or the absence thereof in this language is, again, only contextually induced and tied to neither (I)PFV nor (non-)finite verb forms.\footnote{Throughout the chapter, I call \textit{l}-participles, the only past tense form in both languages, finite verb forms, contrasting them with non-finite verb forms of the adverbial (or other) participle type, discussed in this section. Strictly speaking, Czech \textit{l}-participles are non-finite forms in a periphrastic verb form in combination with the auxiliary `be'. Finiteness and its role for the Russian, as opposed to the Czech aspect system, is further discussed in \citet{gehrkeadv23}.} 

As a final illustration of systematic differences between Czech and Russian in the context of the description of chains of single events, let us look at a Czech original from Bohumil Hrabal's short story \textit{Jetel růžák} and its Russian translation in \REF{sommer} \citep[discussed in][21f.]{gehrke22} \citep[see also][]{gehrke02}%; the relevant verb forms I want to focus on are italicised
. 

\ea\label{sommer}
\ea\label{sommera}	
\gll Když \textit{přišlo} pozdní jaro, když \textit{bylo} léto, když se \textit{setmělo} a \textit{byla} sobota, \textit{přešel} jsem osvětlený most, pak \textit{zahnul} k mlýnu a podle Staré rybárny jsem \textit{kráčel} kolem plotu farní zahrady.\\
when came.\textsc{pfv} late.\textsc{nom} spring.\textsc{nom} when was.\textsc{ipfv} summer.\textsc{nom} when \textsc{refl} got.dark.\textsc{pfv} and was.\textsc{ipfv} Saturday.\textsc{nom} across.went.\textsc{pfv} \textsc{aux.1sg} illuminated.\textsc{acc} bridge.\textsc{acc} then off.bent.\textsc{pfv} to mill.\textsc{dat} and past old.\textsc{acc} fish.restaurant.\textsc{adj.acc} \textsc{aux.1sg} strolled.\textsc{ipfv} around fence parish.\textsc{gen} yard.\textsc{gen}\\
\glt `When late spring arrived, when it was summer, when it got dark and it was Saturday, I crossed the illuminated bridge, then turned to the mill and past the Old Fisherman and strolled along the fence of the churchyard.' \hfill (Czech) %(CZ JR 109)
\ex	
\gll Kogda vesnja \textit{približalas'} k koncu, kogda \textit{bylo} uže počti leto, odnaždy v subbotnie sumerki ja \textit{perešel} osveščennyj most, a potom \textit{svernul} k mel'nice i \textit{zašagal} mimo starogo \minsp{`} Rybnogo podvor'ja' vdol' ogrady cerkovnogo sada.\\
	when spring approached.\textsc{si} to end.\textsc{dat}, when was.\textsc{ipfv} already almost summer.\textsc{nom} once in Saturday.\textsc{adj.pl} twilights I across.went.\textsc{pfv} illuminated.\textsc{acc} bridge.\textsc{acc} and then off.bent.\textsc{pfv} to mill.\textsc{dat} and \textsc{za}.strolled.\textsc{pfv} past old {} fish.\textsc{adj} inn along fence church.\textsc{adj.gen} yard.\textsc{gen}\\
\glt `When spring was coming to its end, when it was already almost summer, one Saturday evening I crossed the illuminated bridge, and then turned to the mill and started strolling past the old Fisherman's Inn along the fence of the churchyard.' \hfill (Russian)
\z 
\z 

\noindent The first four finite verb forms in the Czech original example describe backgrounded events that set the scene for the following passage, which contains a chain of three events that temporally follow one another (`cross the bridge', `turn to the mill', `stroll past the inn along the fence'). The scene setting is done in Czech alternating between PFV achievement predicates (`arrive', `get dark') and IPFV states (`be'). In the Russian translation, on the other hand, only two finite verb forms are used, for the first two events setting the scene, and both appear in the IPFV (`approach.\textsc{si}', `be.\textsc{ipfv}'). The other two scene-setting events (`get dark', `be Saturday') are translated non-verbally (lit. `in Saturday twilight'), and they are backgrounded to the chain of the three following events, the start of which is explicitly marked in Russian (by \textit{{odnaždy}} `once'), but not in Czech. Both languages use PFV accomplishment predicates for the description of the first two events in this chain, but the initial temporal bound of the third event, which temporally follows the second one and is described by the activity predicate `stroll', is marked explicitly only in Russian, by ingressive \textit{za-}, but remains to be deduced from the context in Czech.   

In sum, aspect use in Czech chains of single events is largely governed by the types of predicates employed (IPFVs for states and activities, primarily PFVs for accomplishments and achievements) and also whether the narrator wants to focus on the duration of a process, in which case the IPFV also appears with accomplishments. In contrast, in Russian a SOE interpretation (with single events) requires a PFV finite verb form, independently of the verb class. Let us then turn to iterative and habitual contexts.

\subsection{Iterativity, habituality}
\label{sec:hab}

The literature on Russian aspect opposes the description of single events (\textit{edi\-nič\-nost'}) with \textit{kratnost'}, which I translate as \textsc{repeatability}.\footnote{The Russian terms are abstract nominalisations related to the adjectives for `single' and `multiple', respectively.} The consensus in the literature on Russian is that repeatability requires the IPFV \citep[e.g.][]{paduceva96,zaliznjaksmelev97}, whereas this is not the case in Czech \citep[e.g.][]{%sirokova63,petruxina78,petruxina00,
eckert84,%eckert85,stunova86,
stunova93,Dickey2000,kresin00,gehrke02,gehrke22,duebbersdiss,fortuin2015typology}. From a theoretical point of view, the use of IPFVs in such contexts has been captured by treating the Russian IPFV in these cases as quantifying over sets of plural events, as in \citet{kagandiss,Altshuler2012,altshuler14} \citep[the latter following][]{ferreira2005}.

Russian can use the PFV in iterative contexts when the number of repetitions is known and the stretch of repetitions is presented as one whole. This PFV use has been labeled the \textsc{summation} (\textit{summarnoe}) reading and is illustrated in \REF{ex:sum} \citep[from][19; my glosses and translation]{zaliznjaksmelev97}.

\ea \label{ex:sum}
\gll Ona tri raza postučala v dver'.\\
    she three times knocked.\textsc{pfv} in door.\textsc{acc}\\
\glt    `She knocked on the door three times.' \hfill (Russian)
\z

\noindent Furthermore, in iterative and habitual contexts, Czech productively uses the suffix \textit{-va} to derive \textsc{frequentatives} from all kinds of simple IPFV verbs \REF{ex:freqa}, as well as from SIs \REF{ex:freqb} \citep[example adapted from][]{biskup23adv} (see also \citealt{kopecny62}, \citealt{petr86}, \citealt{Filip.Carlson1997}, \citetv{chapters/04_filip}).\footnote{I use \citeposst{petr86} term for frequentatives, but there is no established term; see \textcitetv{chapters/04_filip} for recent discussion.} 

\ea\label{ex:freq}
\ea IPFV \textit{mít} `have' $>$ IPFV \textit{mí-va-t} `have.\textsc{freq}' \hfill (Czech)\label{ex:freqa}
\ex SI \textit{vy-pis-ova-t} `excerpt' $>$ \textit{vy-pis-ová-va-t} `excerpt.\textsc{freq}'\label{ex:freqb}
\z 
\z

\noindent The derivation of frequentatives is not productive anymore in Russian, apart from a few remaining lexical items (e.g. \textit{byvat'} `be.\textsc{freq}'), and it is not possible from SIs. Some Czech researchers even consider frequentatives to represent a third category of aspect, in addition to IPFV and PFV \citep[see discussion in][]{kopecny62}. Furthermore, the frequentative suffix \textit{-va} is commonly treated as homophonous to one of the imperfectivising suffixes Czech employs in SIs, and the fact that it attaches to a verb form that is already IPFV is taken as an argument in favour of the homophony analysis. I will stay agnostic as to whether synchronically we are dealing with homophony or with the same suffix. Diachronically, both the Czech and the Russian imperfectivising suffix derives from the frequentative suffix, and this is the only remaining productive SI suffix in Russian.  

Let me leave strictly iterative contexts aside and concentrate on habitual ones. The descriptive literature sometimes distinguishes a micro-level/event (for each repetition) from a macro-level/event \citep[see, e.g.,][]{eckert84,stunova93}, and in the following I employ these terms descriptively. For example, while at the micro-level an event can be bounded (e.g. because it appears in a habitual chain of events), we can have non-boundedness at the macro-level, because we are dealing with a habitual discourse. This tension, again, potentially leads to variation in aspect use, namely by using the PFV to signal boundedness at the micro-level, but the IPFV to signal non-boundedness at the macro-level. In habitual contexts, we find both aspects in Czech, arguably conditioned by the same considerations as with single events, whereas Russian almost exclusively uses the IPFV; the Russian IPFV is quasi-obligatory in both present and past tense habitual contexts, with the only exception of the stylistically marked so-called vivid-exemplifying use of the PFV, which I will come back to in the discussion of \REF{ex:naglprim}. 

For illustration, let us look at the Russian example in \REF{28}, from Mixail Bulgakov's \textit{Rokovye jajca}, and its Czech translation \citep[adapted from][87]{gehrke02}. 

\ea\label{28}	
\ea	
\gll Mnogie iz 30 tysjač mexaničeskix ėkipažej, \textit{begavšie} v 28-m godu po Moskve, \textit{proskakivali} po ulice Gercena, \textit{šurša} po gladkim torcam, i čerez každuju minutu s gulom i skrežetom \textit{skatyvalsja} s Gercena k Moxovoj tramvaj {16, 22, 48} ili 53-go maršruta. \\
 	many.\textsc{nom} of 30 thousand mechanical wagons run.\textsc{ipfv.indet.pap} in 28th year along Moscow through.jumped.\textsc{si} along street  Gercen.\textsc{gen} rustle.\textsc{ipfv.ap} over smooth pavement and through every minute with roar and crunch down.rolled.\textsc{si.refl} from Gercen to Moxovaja tram.\textsc{nom} {16, 22, 48} or 53th.\textsc{gen} line.\textsc{gen}\\
\glt `Many of the 30.000 mechanical wagons, running in Moscow in 1928, sped through Gercen street, rustling over the smooth pavement, and every minute Tram lines 16, 22, 48 or 53 rolled down from Gercen street to Moxovaja street, roaring and crunching.' \hfill (Russian)%(RU RJ389/10)
\ex\label{ex:jezdily}	
\gll Mnohé z třiceti tisíc drožek, které v osmadvacátém \textit{jezdily} po Moskvě, \textit{proklouzly} Gercenovou ulicí a \textit{zasvištěly} na hladkém dřevěném dláždění; každou minutu se s řinkotem a skřípěním \textit{přehnala} od Gercenovy ulice k Mechové tramvaj číslo {16, 22, 48} nebo 53.\\
 	many.\textsc{nom} out.of thirty thousand carriages which in 28th drove.\textsc{indet.ipfv} along Moscow through.slid.\textsc{pfv} Gercen.\textsc{adj} street  and swished.\textsc{pfv} on smooth wooden pavement every.\textsc{acc} minute.\textsc{acc} \textsc{refl} with rattle and crunching past.chased.\textsc{pfv} from Gercen.\textsc{adj} street to Mechová  tram.\textsc{nom} number {16, 22, 48} or 53 \\
	\glt `Many of the 30.000 carriages that drove in Moscow in 1928 slid through Gercen street and swished on the smooth pavement. Every minute Tram no. 16, 22, 48 or 53 chased by from Gercen street to Mechová street, rattling and crunching.' \hfill (Czech)
\z 
\z 

\noindent This passage describes the regular public transport in Moscow in the late 1920s%, with tram lines quickly passing by. The verb forms I discuss in more detail are italicised
: wagons quickly move along Gercen Street, ass they rustle over the smooth pavement, and every minute a tramline passes through. Since the passage is habitual, Russian has to use the IPFV, whereas we find the PFV for the Czech descriptions of all foregrounded events; even the backgrounded event that is described by a non-finite adverbial participle in Russian (`rustling') is referred to with a Czech finite PFV form in a main clause (`swished'). Furthermore note that all Russian IPFVs are SIs. Thus, we get the impression, that the morphology is used to both render the quickness and ``completion'' of each micro-event (by a prefix) as well as the habituality and thus unboundedness of the macro-event (by the imperfectivising suffix). In Czech, on the other hand, habituality is deducible almost exclusively from the context, and the only verb form that makes it obvious that we are dealing with a habitual passage is the indeterminate IPFV motion verb \textit{jezdily} in the beginning of this passage (I will come back to this below). 

\iffalse
The example in \REF{ex:sext} from \citet[][85]{gehrke02}, whose Czech original is from Hrabal's short story \textit{Sextánka}, further illustrates \citep[discussed in][]{gehrke02}.

\ea\label{ex:sext}
\ea	\label{ex:sextcz}
\gll A nejen matka, ale to už \textit{přišli} ostatní herci, všechny jsem je znal, protože jsem jim \textit{mazal} chleba se sádlem a \textit{podával} láhve piva, taky \textit{přišli} v pěkných šatech, každý \textit{měl} pod paží srolovanou \'{u}lohu, když se \textit{stmívalo}, tak náměstím \textit{chodilo} sem a tam deset bílých srolovaných divadelních knížek [...]\\
	and \textsc{neg}.only mother.\textsc{nom} but there already to.went.\textsc{pfv} other.\textsc{pl.nom} actors.\textsc{nom} all.\textsc{pl.acc} \textsc{aux.1sg} I knew.\textsc{ipfv} because \textsc{aux.1sg} them.\textsc{dat} smeared.\textsc{ipfv} bread.\textsc{acc} with lard.\textsc{instr} and passed.\textsc{si} bottles.\textsc{acc} beer.\textsc{gen} also to.went.\textsc{pf} in pretty.\textsc{pl.prep} dresses.\textsc{prep} every.\textsc{nom} had.\textsc{ipfv} under arm.\textsc{instr} rolled.up.\textsc{acc} roll.\textsc{acc} when \textsc{refl} darkened.\textsc{n.SI} then square.\textsc{instr} went.\textsc{n.indet.ipfv} here and there ten.\textsc{nom} white rolled.up theatre.\textsc{adj.gen} booklets.\textsc{gen}\\
\glt 	`And not only (the) mother, but quickly also other actors came/arrived, I knew them all, because I made them lard breads and passed them beer bottles, they also came in pretty clothes/costumes, everyone had a rolled-up text under their arm, and when it turned dark, ten white rolled-up theatre texts flew here and there across the square.' \hfill (Czech) %(CZ, Sext 147)
\ex
\gll I tak \textit{vela} sebja ne tol'ko matuška, no i drugie aktery; ja vsex ix znal, {potomu čto} \textit{namazyval} im smalec na xleb i \textit{podaval} butylki s pivom, oni tože \textit{prixodili} prinarjažennye, i u každogo iz nix pod myškoj \textit{torčala} svernutaja v trubku rol', tak čto v sumerkax po ploščadi \textit{rasxaživalo} tuda-sjuda desjat' belyx truboček [...]\\
and so led.\textsc{ind.ipfv} \textsc{refl} not only mother.\textsc{nom} but also other.\textsc{pl.nom} actors.\textsc{nom} I all.\textsc{pl.acc} them.\textsc{acc} knew.\textsc{ipfv} because smeared.\textsc{si} them.\textsc{dat} lard.\textsc{acc} on bread.\textsc{acc} and passed.\textsc{si} bottles.\textsc{acc} with beer.\textsc{instr} they also to-went.\textsc{si} dressed.up.\textsc{nom} and at every.\textsc{gen} out-of them.\textsc{gen} under armpit.\textsc{instr} stuck.out.\textsc{ipfv} rolled.\textsc{nom} in little-tube.\textsc{acc} roll/part.\textsc{nom} so that in dusk.\textsc{prep} across square.\textsc{dat} apart.went.\textsc{si} here-there ten white little.tubes.\textsc{gen} \\
\glt `And not only (the) mother behaved like that, but also other actors; I knew them all because I smeared them lard on bread and passed them bottles with beer, they also came dressed up, and everyone of them had their part rolled up in a little tube sticking out out from under their armpits, so that when it turned dark ten white little tubes went here and there across the square.' \hfill (Russian)
\z
\z

\noindent This passage describes the memories of the protagonist of a regular theatre rehearsal, which involves his mother, who is one of the actors (relevant verb forms in italics): there are actors regularly coming to the rehearsal, they come dressed up, they have their texts rolled up, and when it gets dark, the text rolls are left flying across the square. In between this passage, the protagonist provides background knowledge about why he knows all the actors, namely he regularly smears them sandwiches and passes out beer bottles to them. Apart from this background knowledge, it is not clear from the Czech passage that we are actually dealing with a habitual passage, because the verb forms used would be the same in a description of single events: The arrival of the actors (dressed up) is described by a PFV, having texts rolled up is described by a simple IPFV, because we are dealing with a state, the getting dark is described by a SI, because focus is on the duration or process of getting dark, and the flying around of text rolls is described by a simple IPFV indeterminate motion verb, which could arguably seen as a verbal marker for distributivity, or also as a marker of habituality (I will come back to this below). The only Czech verb form that is arguably explicitly marked for habituality is the SI \textit{podával} `passed', which contains the same suffix as frequentative verbs in Czech (recall \REF{ex:freq}). The situation is quite different in the Russian translation, which makes clear with every verb form that we are in a habitual context: all verb forms are IPFV, and many of them are again SIs, which are doubly marked for event completion at the micro-level (prefix) and habituality/unboundedness at the macro-level (SI suffix).
\fi

There is one exception to the rule that Russian requires IPFVs in habitual contexts, namely the \textsc{vivid-exemplifying} use of PFV present tense forms. Such present tense forms occur independently of whether the habitual passage is in the past or present, so these present tense forms do not necessarily express a present tense meaning. For illustration let us look at %In the data discussed in \citet{gehrke02,gehrke22} this use in past tense contexts occurred just once, namely in the Russian translation of a Czech original that does not contain such present tense forms; see 
the Czech example in \REF{ex:naglprim}, from Milan Kundera's \textit{Žert}, and its Russian translation, discussed in \citet[][29]{gehrke22}% \citep[see also][]{gehrke02}
. 

\ea\label{ex:naglprim}
\ea\label{ex:cekavala} 
\gll [...] v poledne jsme \textit{neměli} čas ani poobědvat, \textit{snědli} jsme na sekretariátě ČSM dvě suché housky a pak jsme se zase třeba celý den \textit{neviděli}, \textit{čekávala} jsem na Pavla kolem půlnoci [...]\\
{} in noon \textsc{aux.1pl} \textsc{neg}.had.\textsc{ipfv} time even have.lunch.\textsc{pfv.inf} ate.\textsc{pfv} \textsc{aux.1pl} on secretariat ČSM two dry rolls and then \textsc{aux.1pl} \textsc{refl} again maybe whole.\textsc{acc} day.\textsc{acc} \textsc{neg}.saw.\textsc{ipfv} waited.\textsc{ipfv.freq.f} \textsc{aux.1sg} on Pavel around midnight \\
\glt `At noon, we did not even have time to have lunch, we ate two dry rolls at the secretariat of the ČSM [Czechoslovak Youth Union] and then again maybe did not see each other the whole day, I used to wait for Pavel around midnight.' \hfill (Czech)%(CZ Kun 21/23)
\ex 
\gll [...] v polden' nam ne \textit{xvatalo} vremeni daže poobedat', \textit{s''edim}, \textit{byvalo}, na sekretariate dve suxie bulki, a potom snova počti celyj den' ne \textit{vidimsja}, \textit{ždala} ja Pavla \textit{obyčno} k polunoči [...]\\
{} in noon us.\textsc{dat} not sufficed.\textsc{ipfv} time.\textsc{gen} even have.lunch.\textsc{pfv.inf} eat.\textsc{pfv.1pl.prs} was.\textsc{ipfv.freq.3sg.n} on secretariat two dry rolls and then again almost whole.\textsc{acc} day.\textsc{acc} not see.\textsc{ipfv.1pl.prs.refl} waited.\textsc{ipfv.f} I Pavel.\textsc{acc} usually to midnight\\
\glt `At noon we did not even have time to have lunch, we used to eat two dry rolls at the secretariat and then not see each other almost the whole day, I usually waited for Pavel until midnight.' \\\hfill (Russian)
\z 
\z 

\noindent This passage describes the daily sequential routine that the female narrator had with her partner Pavel. It starts with a negated stative expression (`not have time') that explains the shortness of the first event in the sequence (a quick lunch), followed by not seeing each other during the day, and then her waiting for Pavel around/until midnight. In the Czech original, all four verb forms are finite past tense forms; the one that describes the reason (`not have time') is in the IPFV because it describes a state. The Russian translator chose a different lexical item, but also an IPFV. In Czech, the first event of the habitual chain (`eat two dry rolls') is described by a PFV verb form, followed by an IPFV negated stative description (`not see each other'). The PFV verb form appears because at the micro-level two rolls were finished, rather quickly, %eaten %(each day two different ones) 
and %given the context it most likely involves a quick lunch so that 
contextually it does not make sense to focus on the duration% the duration of the %se eating event%s
. Up until here the Czech verb forms by themselves do no indicate that the passage is habitual, and apart from the adverbs `again maybe' %the clauses containing 
these three verb forms could also be used in the description of single events. Only the last form, frequentative \textit{čeká-va-la}% (derived from a simple IPFV)
, explicitly signals habituality%, since we are dealing with a frequentative \textit{čeká-va-la} (derived from a simple IPFV)%, derived from the simple IPFV \textit{čekala} `waited'
. 

Things are different in Russian when the chain starts: the translator switched to two present tense forms in the vivid-exemplifying use (PFV `eat', IPFV `see'), and this tense switch is accompanied by the addition of the habituality marker \textit{byvalo} `it used to be', which is absent from the Czech original. It is commonly assumed that the switch to the vivid-exemplifying use %of the present tense 
has to be accompanied by expressions like \textit{byvalo} %for the switch to be possible 
\citep[see, e.g.][]{zaliznjaksmelev97}. The vivid-exemplifying use of the PFV present is obviously a stylistic device; the translator could also have stayed in the past tense, in which case, however, the PFV would not have been possible. For the last verb form the translator switched back to the past tense and translated the frequentative Czech verb for `waited' with a simple IPFV verb `waited', as Russian cannot derive a frequentative from this kind of verb (\textit{*žda-va-la}); however, they added the adverb `usually' to render the habitual nature of the Czech verb form, even though there is no such adverb in the Czech original.

Let us see what these examples illustrate in general, as argued in more detail in \citet{gehrke22}. In Russian, the finite verb forms are explicitly marked for two things: a) event sequencing and/or event completion (on a par with what would be the case in the description of single events, as we saw in \S \ref{sec:chains}), and b) habituality, by additionally imperfectivising the verb forms, which is the main difference between chains of single vs. habitual events in this language. In Czech, on the other hand, only few verb forms make aspect use in habitual contexts different from single events: the indeterminate motion verb \textit{jezdily} in \REF{ex:jezdily} %sextcz}, this is possibly the indeterminate motion verb \textit{jezdily%chodilo}% (if it does not signal distributivity alone) and an SI with the same suffix we find with frequentative verb
and the frequentative verb \textit{čekávala} in \REF{ex:cekavala}. Both types of verbs are common means in Czech to signal habituality (see \citealt{eckert91} for indeterminate motion verbs and \citealt{Filip.Carlson1997} for frequentatives). These verb forms appear once in a passage, to mark it as habitual; other verb forms in the same passage display the same kind of aspect use as with single events: ``completed'' and/or quick events are described by PFVs, stative events or events of some duration are described by IPFVs. Thus, in habitual contexts, Czech uses aspect more or less the same as in the description of single %, episodic 
events (see \citealt{eckert84} for the same conclusion).\footnote{For example, only a small fraction (about 7\%) of the past tense forms in Czech habitual contexts in \citeposst{gehrke02} corpus analysis explicitly mark habituality: Out of about 500 past tense forms, these were 16 frequentatives% (vs. only 3 occurrences of \textit{byvat'} `be.\textsc{freq}' in Russian)
, 3 indeterminate verbs of motion, 3 SI verbs of motion \citep[which sometimes also mark habituality, according to][]{eckert91}, 9 SIs that (at least formally) with a suffix that is formally identical to frequentatives (\textit{-va}), and 3 SIs with other suffixes (all exclusively in the Czech translation of Sergej Dovlatov's \textit{Zapiski nadziratelja}).%\footnote{The latter three forms occurred exclusively in the translation of Dovlatov, which, upon further inspection, was perceived by native speakers consulted as not the best translation to begin with, so there could have been some transfer errors from the Russian original as well \citep[for further discussion see][]{gehrke22}.} 
}

This of course does not mean that Czech does not use the IPFV in habitual contexts. Rather, habituality is not directly marked on the verb form in most of the cases, apart from a few specialised IPFVs %, or possibly also transfer errors from Russian
once in a longer passage. In \citet{gehrke22}, it is therefore argued that the use of the IPFV in Czech habitual contexts can be explained by the same reasoning that explains its use in the description of single events: for atelic states and activities, as well as for accomplishments with a focus on their duration. Furthermore, nothing is said about the verb forms in Czech being the only options, it might very well be that some PFVs could be replaced by IPFVs, this would have to be checked on a case-by-case basis. It is just that as a generalisation on the data taken into account, the conclusion is that whenever there is no need to use an IPFV (for atelic events or for events of a certain duration) Czech simply uses the PFV; this is also in line with the intuitions reported by Czech native speakers that investigated differences between Czech and Russian, such as \citet{eckert84} and \citet{stunova93}. And this is exactly where Russian differs: in Russian one has to use the IPFV, the PFV would be ungrammatical in habitual contexts. 
%In sum, aspect use in habitual contexts in Czech is quite similar to its use in the description of single events; only sporadically we get (mostly) specialised IPFV verb forms to signal that we are in a habitual context, so that habituality is in most cases only contextually induced. In contrast, in Russian each finite verb form has to be IPFV in habitual contexts (with the exception of the stylistically marked vivid-exemplifying use of the PFV present). In many cases this involves adding an imperfectivising suffix, so that the overall discourse strategies one finds in chains of single events, outlined in the previous section, remain intact. 
Let us then turn to contexts that involve the historical (or narrative) present. 

\subsection{Historical present}
\label{sec:HP}

The historical present is a stylistic device in narrative texts to describe events using present tense forms (even if the events happened in the past), to make them appear more vividly, as unfolding before one's eyes (see also \citealt{anandtoo} for recent formal discussion). From a theoretical and cross-linguistic point of view, it is commonly assumed that with a regular present tense semantics, the R(eference Time) coincides with the S(peech Time) (\textsc{now}). Since  \textsc{now} is a point and not an interval, but since the truth-conditions for ``completed'' events require intervals to be evaluated, there is a common idea that (``true'') present tense and PFV are semantically incompatible. In both Czech and Russian, for example, present tense morphology on PFV verbs is (often) interpreted as reference to events in the future, rather than in the present. Hence, one could expect that in the historical present, which also seems to describe events as if evolving before our eyes, only the IPFV is used. On the other hand, if the present tense in the historical present is semantically not a ``true'' present tense, %or more drastically, if Aspect and Tense are two completely different systems, 
we would not have this expectation. While I refrain from giving a semantic account of the historical present, we seem to observe two different situations in Czech and Russian. In particular, Russian exclusively uses the IPFV in historical present contexts (except for the vivid-exemplifying use of the PFV; recall \REF{ex:naglprim}), whereas we find both aspects in Czech \citep{krizkova55,%bondarko58,
bondarko59,petruxina83,stunova93,%stunova94,
Dickey2000,fortuin2015typology}.  

\iffalse 

Let me illustrate with two examples from works by Karel Čapek, with translations into Russian, which are discussed in \citet{stunova93} (my English glosses and translations).

\ea\label{ex:histpres}
\ea 
\gll Tedy ten sedlář sedí po obědě v kruhu své rodiny, a najednou mu někdo \textit{tluče} na okno: `Ježíšmarjá, sousede, vždyt' vám hoří střecha nad hlavou!' Ten sedlář \textit{vyletí}, a na mou duši, krov je v plamenech. To se ví, děti se \textit{dají} do křiku, žena s pláčem \textit{vynáší} hodiny (...)\\
so this saddler sits.\textsc{ipfv} after lunch in circle \textsc{possrefl.gen} family.\textsc{gen} and suddenly he.\textsc{dat} someone knocks.\textsc{ipfv} on window {Jesus Mary} neighbour.\textsc{voc} always you.\textsc{dat} burns.\textsc{ipfv} roof over head this saddler out.flies.\textsc{pfv} and on my soul roof is in flames that \textsc{refl} knows.\textsc{ipfv} children \textsc{refl} give.\textsc{pfv.3pl} to scream woman/wife with cry outcarries.\textsc{si} clock {} \\
\glt `So, this saddler sits together with his family after lunch, and suddenly someone knocks on his window: ``Jesus, neighbour, your roof is on fire!" The saddler runs out, and, I swear, the roof is in flames. Obviously, the children start screaming, the wife carries out the clock, crying.' \hfill (Czech)
\ex
\gll Tak vot, ėtot šornik sidit sebe posle obeda so svoim semejstvom i vdrug kto-to \textit{stučit} v okno. -- Sosed, u vas kryša gorit! Šornik \textit{vybegaet} na ulicu -- i verno, kryša u nego vsja v ogne. Nu, konečno, deti \textit{revut}, žena s plačem \textit{vynosit} stennye časy.\\
so there this saddler sits.\textsc{ipfv} self after lunch with \textsc{possrefl} family and suddenly someone knocks.\textsc{ipfv} in window {} neighbour at you roof burns.\textsc{ipfv} saddler outruns.\textsc{si} on street {} and truly roof at him all in fire now of.course children roar.\textsc{ipfv} women/wife with cry outcarries.\textsc{si} wall.\textsc{adj} clock \\
\glt `So, this saddler sits down after lunch with his family and suddenly someone knocks on the window. ``Neighbour, your roof is burning!'' The saddler runs onto the street, and really, his whole roof is on fire. Now, of course, the children roar, the wife carries out the wall clock crying.' \hfill (Russian)
\z
\z

\noindent The whole passage of the narration in \REF{ex:histpres} \citep[from][185]{stunova93} is in the present tense so it arguably involves the use of the historical present. The first part (sitting at the table) serves to set the scene for the following sequence of events that is taking place (the relevant verb forms are italicised): someone knocking on the window and saying something to the family, the saddler running out of the house, the children starting to scream or roaring, and the wife carrying out the wall clock. In this sequence of events, Russian exclusively uses IPFVs (PFVs would not be possible), but the second and third events are described by PFVs in the Czech original. \citeauthor{stunova93} argues that replacing these two forms with IPFVs would focus too much in a literal process of flying out in the first case, and would express hesitation or distributiveness in the case of the children starting to scream, which would not fit the context. She does not comment on the first event of knocking on the window (which could be one or several knocks) or the last event of carrying out, which is described by a SI in both languages. I suspect that in the latter case there might be an additional effect of the use of this SI in Czech, namely that the wife is dragging out the wall clock (or at least that there is more of a focus on the process or duration of this event), whereas in Russian the IPFV is used because it is required in the historical present, and whether or not this was a longer event of dragging would remain implicit in this context.  
\fi

Let me illustrate with the example in \REF{ex:krak} from \citet[][187; my English glosses and translations, context abbreviated]{stunova93}, from Karel Čapek's novel \textit{Krakatit}. %For further, similar examples, see \citet{stunova93}. 

\eanoraggedright\label{ex:krak}
Context: `Indeed, nothing compares to the beauty of a summer morning, but Prokop looks down to the ground, smiles as far as he is able to do so, and wanders through small gates to the river.'
\ea
\gll %... vskutku, nic se nevyrovná kráse letního jutra, ale Prokop se dívá do země, usmívá se, pokud to vůbec dovede, a putuje samými závorkami až k řece. 
Tam \textit{objeví} -- ale u druhého břehu -- poupata leknínů; tu zhrdaje vším nebezpečím se \textit{svlékne}, \textit{vrhne} se do hustého slizu zátoky, \textit{pořeže} si nohy o nějakou zákeřnou osřici a \textit{vrací} se s náručí leknínů.\\
%{} indeed nothing \textsc{refl} compares beauty.\textsc{dat} summer.\textsc{adj.gen} morning.\textsc{gen} but Prokop \textsc{refl} looks.\textsc{ipfv} to ground smiles.\textsc{si} \textsc{refl} if that {at all} can.\textsc{ipfv} and wanders.\textsc{ipfv} same.\textsc{instr.pl} brackets.\textsc{instr} up to river 
there appear.\textsc{pfv.prs.3pl} {} but at other bank {} buds.\textsc{nom} water.lilies.\textsc{gen} here disregarding all danger \textsc{refl} undresses.\textsc{pfv} throws.\textsc{pfv} \textsc{refl} in thick slime bay.\textsc{gen} cuts.\textsc{pfv} \textsc{refl} legs about some insidious sedge and returns.\textsc{ipfv} \textsc{refl} with armful water.lilies.\textsc{gen} \\
\glt `%Indeed, nothing compares to the beauty of a summer morning, but Prokop looks down to the ground, smiles, if he can at all, and wanders through gates to the river. 
There appear -- but on the other river bank -- water lily buds; disregarding all dangers, he undresses, throws himself into the thick slime of the bay, cuts his legs on some insidious sedge and returns with an armful of water lilies.' \hfill (Czech)
\ex 
\gll %(...) {v samom dele,} ničto ne svranitsja krasotoju s letnim utrom, no Prokop ne otryvaet glaz ot zemli, ulybaetsja v meru svoego umen'ja i čerez množestvo kalitok dobiraetsja do reki. 
I tam, tol'ko u protivopoložnogo berega, \textit{obnaruživaet} butony kuvšinok; prenebregaja vsemi opasnostjami, on \textit{snimaet} plat'je i \textit{brosaetsja} v gustuju sliz' zavodi, \textit{ranit} nogi o kakie-to kovarnye ostrye list'ja, no \textit{vozvraščaetsja} s oxapkoj cvetov.\\
%{} indeed nothing \textsc{neg} compares.\textsc{ipfv} beauty.\textsc{instr}  with summer.\textsc{adj.instr} morning.\textsc{instr} but Prokop \textsc{neg} away-tears.\textsc{si} eye from ground smiles.\textsc{si} in measure \textsc{possrefl.gen} ability.\textsc{gen} and through many gates reaches.\textsc{si} to river 
and there only at opposite bank finds.\textsc{si} buds.\textsc{acc} water.lilies.\textsc{gen} disregard.\textsc{si.ap} all dangers he off.takes.\textsc{si} clothes and throws.\textsc{si.refl} in thick slime bay.\textsc{gen} injures.\textsc{ipfv} legs about some insidious sharp leaves but returns.\textsc{si} with  armful flowers.\textsc{gen}\\
\glt `%Indeed, nothing compares in beauty to a summer morning, but Prokop des not take his eyes off the ground, smiles to the degree of his ability and through many gates reaches the river. 
And there, just on the opposite river bank, he finds water lily buds; disregarding all dangers, he takes off his clothes and throws himself into the thick slime of the bay, injures his legs on some insidious sharp leaves, but returns with an armful of flowers.' \hfill (Russian)
\z
\z

\noindent This passage begins with setting the scene of the protagonist Prokop walking to the river, the Czech and Russian wordings I left out; both languages use the IPFV present here, arguably for different reasons though: In Czech we are dealing with states, activities or the dwelling on the process of an accomplishment, in Russian the historical present requires IPFVs. This scene-setting is followed by a sequence of events: water lilies becoming visible, Prokop undressing, throwing himself into the river, cutting or injuring himself, and returning. In Russian, we find only IPFV present tense forms, but apart from the return from the river, all Czech present tense forms are PFV. \citeauthor{stunova93} states that discovering flowers is a momentaneous event, so an IPFV would not be suitable; describing the undressing with an IPFV, she argues, would suggest that it took a long time, IPFV in-throwing would have the effect of an eye-witness report, and IPFV cutting would suggest intentionality, all not desired effects in this context. She furthermore states that the event of %s of walking to and 
returning from the river is described by an IPFV in Czech because we get the impression that Prokop is on his way, rather than arriving (the same is stated for the first event of wandering to the river). 

In sum, aspect use in Czech historical present contexts does not differ much from other contexts we have discussed so far: IPFVs appear with states, activities and accomplishments of a longer duration, PFVs appear in the other cases. In contrast, Russian has to use the IPFV throughout, and so the overall picture is quite similar to aspect use in habitual contexts, as addressed in the previous section. Finally, let us turn to factual contexts. 

\subsection{Factual contexts}
\label{sec:OFCZRU}

As mentioned in \sectref{sec:OF}, the existential IPFV appears in contexts in which the precise time at which the event occurred is not relevant or not known, or it might have happened more than once. \citet{paduceva96}, for instance, argues for Russian that potential repeatability (what she calls \textit{kratnost'}, as outlined in the beginning of \sectref{sec:hab}) is a requirement for the existential IPFV. Recall from \sectref{sec:OF} that with scopally non-specific temporal expressions like \textit{kogda-nibud'} `ever' (see also \sectref{sec:defspec}), Russian has to use the IPFV. This is different in Czech, as can be seen in \REF{ex:kogda} \citep[adapted from][10]{klimek22} \citep[see also discussion in][]{Dickey2000,gehrke02,gehrke22,Mueller-Reichau2018a}.\footnote{Some of the examples from \citet{klimek22} discussed in this section are adjusted qua glosses, translations, and scientific transliterations for Russian.}$^,$\footnote{\label{fn:kogdadoubt}A Czech reviewer observes that he finds the ``universal \textit{-koli(v)}-words'' bad with the PFV \textit{ztratil} in \REF{ex:kogda} and would use the existential \textit{někdy} instead or adjust the sentence. I keep the Czech judgments reported in \citet{klimek22} here, but I will come back to this point in \sectref{sec:definiteness}.}

\ea\label{ex:kogda}
\ea 
\gll \minsp{\{} Ztratil / \minsp{\textsuperscript{??}} ztrácel\} jsi kdykoliv klíče?\\
    {} lost.\textsc{pfv} {} {} lost.\textsc{ipfv} \textsc{aux.2sg} ever keys\\
\glt `Have you ever lost keys?' \hfill (Czech)
\ex
\gll Ty kogda-libo \minsp{\{*} poterjal / terjal\} ključi?\\
    you ever {} lost.\textsc{pfv} {} lost.\textsc{ipfv} keys\\
\glt    `Have you ever lost keys?' \hfill (Russian)
\z 
\z 

\noindent In this example, Czech has a preference for the PFV aspect, and one could argue that this is so because losing keys is a punctual event that does not involve a process or a duration (or even intentionality); Russian, on the other hand, has to use the IPFV because of the non-specific temporal adverb \textit{kogda-libo} `ever'. 

\citet{Dickey2000}, who does not differentiate between different subtypes of factual IPFVs, argues that factual IPFVs are incompatible with achievement predicates in Czech but not in Russian, and the examples he provides in \REF{ex:ofach} \citep[from][99 \& 101; my glosses and translations]{Dickey2000} illustrate the existential IPFV.\footnote{A Czech reviewer notes that \REF{ex:ofacha} should be modified to reflect the information structure: Since \textit{toho} `this.\textsc{gen}' brings about givenness, \textit{z toho stromu} `from this tree' should be placed before \textit{jednou} `once'. This does not affect the judgments about (I)PFVs though.} 

\ea\label{ex:ofach}
\ea 
\gll Jako dítě jsem jednou \minsp{\{} spadl / \minsp{*} padal\} z toho stromu.\\
    as child \textsc{aux.1sg} once {} fell.\textsc{pfv} {} {} fell.\textsc{ipfv} from this tree\\
\glt `As a child, I once fell from this tree.' \hfill (Czech)\label{ex:ofacha}
\ex 
\gll Ja pomnju, v detstve odnaždy ja \minsp{\{} upal / padal\} s \.etogo dereva.\\
    I remember in childhood once I {} fell.\textsc{pfv} {} fell.\textsc{ipfv} from this tree\\
\glt `I remember, in my childhood I once fell from this tree.' \hfill (Russian)
\z 
\z 

\noindent \citet{Mueller-Reichau2018a}, in turn, argues that it is not about achievements, but about necessarily unique events; he provides examples parallel to those in \REF{ex:ofach}, only with an accomplishment VP, `fell our single tree', in which Czech has to use the PFV as well.

\citet{klimek22} obtained data from Czech, Polish, and Russian, using questionnaires in which the informants had to fill in the missing verbs in both existential and presuppositional factual contexts. Existential contexts are further divided into neutral and resultative ones, presuppositional ones into strongly and weakly resultative ones and also whether focus is on the initiator or on the result state of the event. Independently of the context, all presuppositional examples involving verbs of creation (e.g. `cook', `build') were treated as strongly resultative, because reference to created objects is argued to involve the presupposition of a result state; in contrast, weakly resultatives are assumed to involve other ``accomplishment'' predicates (e.g., `iron' and `wash'). 

\REF{ex:re} \citep[adapted from][16; relevant verb forms italicised]{klimek22} is an example that is argued to involve a resultative existential.\footnote{A Czech reviewer notes that there is a discrepancy in \REF{ex:rea} between singular \textit{kytka} `flower' and the plural glosses `flowers'. I kept the example as it is in the cited paper.}

\ea\label{ex:re}
\ea
\gll Vidím, že kytka na římse zvadla. Určitě jsi je dnes \textit{zalíval}?\\
    see.\textsc{ipfv.1sg} that flowers on window.sill wilted.\textsc{pfv} surely are.\textsc{ipfv.2sg} them today watered.\textsc{si}\\
\glt `I see that the flowers on the window sill have wilted. Are you sure that you watered them today?' \hfill (Czech)\label{ex:rea}
\ex 
\gll Ja vižu, čto cvety na podokonnike zasoxli. Ty uveren, čto \textit{polival} ix segodnja?\\
    I see.\textsc{ipfv.1sg} that flowers on window.sill wilted.\textsc{pfv} you certain that watered.\textsc{si} them today\\
\glt `I see that the flowers on the window sill have wilted. Are you certain that you watered them today?' \hfill (Russian)
\z 
\z 

\noindent Both languages can use the IPFV in this context, and \citeauthor{klimek22} takes this as an argument that both languages make use of existential IPFVs. On the other hand, it is not obvious to me that this example necessarily involves resultativity; it could also be a question about whether the process of watering took place, in which case the use of IPFVs in Czech could be motivated by the process reading of IPFVs. I will come back to this in \sectref{sec:OFgehrke}.

The example in \REF{ex:srpfr} \citep[adapted from][17]{klimek22} is argued to illustrate strongly resultative presuppositional IPFVs, with focus on the result.

\ea\label{ex:srpfr}
\ea
\gll Ta placka je skvělá, Maryšo. Z jakých ingrediencí jsi ji \textit{pekla}?\\
    this.\textsc{f} cake.\textsc{f} is delicious Maryša.\textsc{voc} out.of which ingredients \textsc{aux.2sg} her baked.\textsc{ipfv}\\
\glt `Your cake is delicious, Maryša, What ingredients did you bake it with?' \hfill (Czech)
\ex 
\gll Kakoj vkusnyj pirog, Marysia! Iz čego ty jego \textit{pekla}?\\
    how delicious cake Marysia out.of what you it baked.\textsc{ipfv}\\
\glt `What a delicious cake, Marysia! What did you bake it with?' \\\hfill (Russian)
\z 
\z 

\noindent This example involves a presuppositional IPFV, because the baking event is presupposed by the object that came into existence through this event, the cake, which the previous context talked about. \citeauthor{klimek22} argues that focus is on the result, because the question is concerned with properties of the result of the baking event (the cake). Alternatively, it could be argued that a question about the ingredients of a baking event is more concerned with the process of baking, rather than with the result, so that there is no focus on the result. 

Finally, \REF{ex:wrpfi} \citep[adapted from][19%; I added the question mark in \REF{ex:wrpfia}
]{klimek22} is presented as an example for a weak resultative presuppositional IPFV, with focus on the initiator. 

\ea\label{ex:wrpfi}
\ea
\gll Vidím, že tvé kolo je nakonec funkční. Právě také hledám odborníky. Můžeš mi říct, kdo ti ho \textit{opravoval}?\\
    see.\textsc{ipfv.1sg} that your bike is finally working just also look.for.\textsc{ipfv.1sg} specialists can.\textsc{2sg} me tell who you.\textsc{dat} it.\textsc{acc} repaired.\textsc{si}\\
\glt `I see that your bike is finally working. I am also looking for specialists. Can you tell me who repaired it for you?' \hfill (Czech)\label{ex:wrpfia}
\ex
\gll Ja vižu, čto tebe nakonec-to počinili velosiped. Ja tože išču mastera. Podskaži, kto tebe jego \textit{činil}?\\ 
    I see.\textsc{ipfv.1sg} that you.\textsc{dat} finally repaired.\textsc{pfv.3pl} bike I also look.for.\textsc{ipfv.1sg} specialist tell.\textsc{pfv.imp.sg} who you.\textsc{dat} it.\textsc{acc} repaired.\textsc{ipfv}\\
\glt `I see that they finally repaired your bike. I am also looking for a specialist. Tell me, who repaired it for you?' \hfill (Russian)
\z 
\z 

\noindent Again, both languages can employ the IPFV, and in Russian the presupposed event is even referred to by the same PFV lexeme in the first sentence of the context (\textit{počinili} `repaired.\textsc{pfv}'), whereas in Czech the repairing event can be accommodated in a context that is preceded by a question about a bike repair specialist. Since we are not dealing with a verb of creation but with a verb that describes an event (of repairing) that the object in question (the bike) can in principle undergo multiple times, \citeauthor{klimek22} labels this a weakly resultative context; it is furthermore argued to involve focus on the initiator, because the question is about the person that initiated the event.  

Based on the data that she obtained, \citet{klimek22} provides a statistical analysis that yields the following significant results. In all factual contexts, Russian employs more IPFVs than Polish and Czech. In neutral existential contexts all three languages use significantly more IPFVs than in the other contexts, whereas in resultative existential contexts more PFVs were used than in neutral ones. Furthermore, only Czech and Polish can also use the PFV in neutral existential contexts with a temporal adverb `ever', as we already saw in \REF{ex:kogda}. There were no statistically significant differences between weakly and strongly resultative presuppositional contexts, nor between focus on initiator or result, for either of the languages. I will therefore not come back to these details in the discussion of \citeposst{klimek22} account in \sectref{sec:klimek}.

In the three examples in \REF{ex:re}--\REF{ex:wrpfi}, both Czech and Russian use the IPFV, so \citet{klimek22} assumes that both languages use the IPFV here precisely because we are dealing with (different types of) factual contexts. On the other hand, all her examples involve verbs of creation (a subtype of incremental theme verbs) or predicates that she groups under the label ``accomplishment''. However, it is not clear that incremental theme verbs are accomplishments to begin with and not just activities \citep[see, e.g.,][]{kennedyincr}, or what \citet{rappaportlevin10} would label manner verbs (in contrast to result verbs), and also some of the other verbal predicates she employed could be argued to involve activities, rather than accomplishments.\footnote{The translations of the complete list of the verbs used in Czech and Russian are as follows: `eat', `mow', `renovate', `drill', `clean', `water', `take out', `bring', `milk', `vacuum', `build', `sculpt', `write', `paint', `sew', `embroider', `bake', `sign', `comb', `wash', `cook', `iron', `cut'.} So it might very well be that Czech uses the IPFV in some of these examples for the same reason it uses IPFVs in other contexts: to focus on the process of a durative event. I will come back to this point in \sectref{sec:OFgehrke}%, but for now I will set the empirical picture aside and turn to the few formal semantic accounts or more formally framed thoughts on the differences between Czech and Russian in aspect use
. Let us then turn to a new empirical domain, for which cross-Slavic differences in aspect use have not been well-described yet: passives. 

\section{Passives}
\label{sec:passives}

This section addresses differences in aspect use between Czech and Russian with different types of passives. Passives are a typical syntactic topic, with many subtopics which are orthogonal to aspect use, so in this chapter I concentrate only on those subtopics that might stand in direct relation to the semantics of aspect.\footnote{As we will see below, a passive reading in Slavic languages is found either with past passive participles (PPPs) or with reflexively marked verb forms. The literature on these is usually not (primarily) concerned with aspect. For further information on passives, not related to aspect, see \citet{veselovskakarlik, karlik17, karlik20, medovawiland, cahamedova20} for Czech PPPs, \citet{Schoorlemmer1995,Paslawska.Stechow2003,borik13,borik14} for Russian PPPs, %\citet{polancec15} for Croatian PPPs,\citet{medovawiland,%bondarozwa18a, bondarozwa18b} for Polish PPPs, 
\citet{babby09} for Russian more generally, \citet{medovadiss} for Czech reflexives, and \citet{fehrmann+10} for Slavic reflexives.}
 In this context, one might even wonder why aspect would play a role at all, given the standard assumption that passives and actives do not differ truth-conditionally. Thus, one might expect aspect to function the same in passives as in actives. It turns out that this is true for Czech, but not for Russian. In the following, I first address general cross-linguistic assumptions about passives to then turn to a comparison between Czech and Russian.

\subsection{Different types of passives, cross-linguistically}
\label{sec:passtypes}

From a morphological point of view, different constructions have been labeled passives or taken to express a verbal passive meaning/syntax (on which see below). For Slavic languages, two morphological passive types are relevant. 
The first type is the \textsc{participial passive}, which is a periphrastic verb form made up of an auxiliary in combination with a past passive participle (PPP), as it is found in, for instance, Germanic and Romance languages. The second type is the \textsc{reflexive passive}, not found in Germanic languages, but in, e.g., Romance: a reflexively marked verb form that is interpreted on a par with a passive construction (among other readings reflexive forms can have, such as reflexive, reciprocal, inchoative/anticausative, impersonal, middle; see, e.g., \citealt{babby09,medovadiss,fehrmann+10} for discussion of Slavic). Examples from Czech for both types of passives are given in \REF{ex:czpasstypes} (from \citealt{karlik17}, my glosses and translation).

\ea \label{ex:czpasstypes}
\ea 
\gll Píše se stížnost. \\
	writes.\textsc{ipfv} \textsc{refl} complaint.\textsc{nom} \\
\ex	
\gll Je psána stížnost.\\
	is write.\textsc{ipfv.ppp} complaint \\
\z	
`A complaint is (being) written.' \hfill (Czech)
\z 

\noindent Syntactically and semantically it is common for participial passives to distinguish between verbal and adjectival passives, or between eventive and stative passives.\footnote{In many works these two dichotomies are used synonymously, even though they are not straightforwardly synonymous (there are stative verbs and possibly also eventive adjectives). Since adjectival passives involve adjectivisation of a PPP, as argued below, but reflexive verb forms are finite verb forms that cannot be adjectivised, this distinction has, to my knowledge, not been addressed and is probably not relevant for reflexive passives.} Within the adjectival or stative type, a further distinction is assumed to hold between target state and resultant state participles \citep[following][]{kratzer00}, based for example on modifiability by `still', which is only possible with the former. \citeauthor{kratzer00} argues that \textsc{target state participles} are derived from category-neutral stems that make available both an event and a target state argument ($\sim$~accomplishments and achievements), either by lexical or phrasal stativisation. \textsc{Resultant state participles}, in turn, are proposed to additionally involve a perfect operator (\citeauthor{kratzer00} labels this the ``job-done'' reading) and to necessarily be the result of phrasal stativisation. Another influential distinction among adjectival/stative PPPs is made by \citet{embick04} who differentiates between \textsc{resultatives} ($\sim$ \citeauthor{kratzer00}'s phrasal target and resultant state participles) and \textsc{statives} ($\sim$ \citeauthor{kratzer00}'s lexical target state participles).%\footnote{For further information about adjectival/stative participles in various languages, see \citet{rapp96,anagno03, %alexanagno08, 
%maienborn09, %maienborn11, 
%meltzer11, doron12, mcintyreadjpass, alexiadou+lingua, alexiadou+15, gehrkenllt}.}  

In this chapter, I assume that a \textsc{verbal passive} is a canonical passive that involves an operation on the argument structure of a verbal predicate, suppressing the external argument, which optionally surfaces in a `by'-phrase, and promoting the internal argument to subject position. The input requirement to a canonical verbal passive are therefore verbs with external and internal arguments. In contrast, I take an \textsc{adjectival passive} to be a copular construction in which a copula combines with an adjectival or adjectivised PPP. There is crosslinguistic variation as to whether adjectival passives allow for all kinds of `by'-phrases or whether there are restrictions on their availability (e.g. German, English vs. Greek; see \citealt{alexiadou+lingua}). For example, in German, in which a verbal PPP combines with the auxiliary \textit{werden} `become' and an adjectival PPP with the copula \textit{sein} `be', the verbal passive allows for all kinds of `by'-phrases, but the adjectival passive can only appear with non-referential `by'-phrases that access an event kind but not an event token \REF{ex:germadjverb} \citep[following][]{gehrkenllt}.

\ea\label{ex:germadjverb} (German)
\ea
\gll Das Bild wurde \minsp{\{} von Marta / von einem Kind\} gemalt.\\
	the picture became {} by Marta {} by a child paint.\textsc{ppp}\\
\glt	`The picture was/has been painted by \{Marta / a child\}.' \hfill \textsc{verbal PPP}
\ex
\gll Das Bild war \minsp{\{\#} von Marta / von einem Kind\} gemalt.\\
	the picture was {} by Marta {} by a child paint.\textsc{ppp}\\
\glt	$\sim$ `The picture was painted in a childish manner.' \hfill \textsc{adjectival PPP}	
\z 
\z

\noindent Additional tests to distinguish verbal/eventive from adjectival/stative PPPs involve the possibility of spatiotemporal modification of the underlying event (only with verbal/eventive PPPs), or the compatibility with adjectival morphology (if at all, only with adjectival/stative PPPs). I will come back to some of these tests in the discussion of Czech and Russian PPPs below. 

For languages like Czech and Russian it is additionally important that PPPs in predicative position can appear in two forms, i.e. with long or short agreement morphology (in the following: \textsc{long vs. short form}), e.g. Czech \textit{otevřena/otevře\-ná} / Russian \textit{otkryta/otkrytaja} `open(ed).\textsc{(nom.)fem.sg}'. (Predicative) long forms are commonly assumed to be adjectival \citep[e.g.][]{veselovskakarlik, borik13}; I will set them aside here and focus only on predicative short form PPPs.

\subsection{Russian passives}
\label{sec:russpassives}

Russian \textsc{participial passives} (i.e. predicative short form PPPs in combination with `be') have been argued to exclusively be adjectival by \citet{Paslawska.Stechow2003} (in analogy to their Greek counterparts; cf. \citealt{anagno03}), or to be ambiguous between a verbal and adjectival passive reading \citep[]{Schoorlemmer1995,borik13,Borik.Gehrke2018}, and by form alone it is impossible to distinguish between the two. Based on the following evidence, I side with the latter view in this chapter: Any type of `by'-phrase, which is an NP in instrumental case, is available, and it is possible to locate the underlying event in space and time, e.g. by temporal adverbials. For example in \REF{ex:ambig} \citep[adapted from][122]{borik13}, both the underlying event and the (result) state can be targeted by temporal modifiers, whereas with adjectival/stative passives, only the state itself should be available for such modification.

\ea\label{ex:ambig}
\gll Vorota \minsp{(} byli) otkryty storožem rovno v 6 utra na 2 časa.\\
gates {} were open.\textsc{pfv.ppp} watchman.\textsc{instr} exactly in 6 morning on 2 hours\\
\glt `The/A gate was opened by a/the watchman exactly at 6 in the morning for two hours.'
\z 

\noindent The auxiliary \textit{byt'} `be' that combines with the PPP appears with past (\textit{byl(a/o/i)}) or future (\textit{budu} etc.) tense marking, or it is zero (phonologically empty) in the present tense.\footnote{Throughout I gloss passive auxiliaries as forms of `be', whereas I use \textsc{aux} to gloss the auxiliary that appears in Czech past tense forms (see \sectref{sec:CZPPP} on the Czech counterparts).} Unlike the copula \textit{byt'}, which can appear with frequentative marking (e.g. \textit{by-va-t'}), the passive auxiliary in Russian does not allow it. 

There is a wideheld assumption that Russian PPPs can only be derived from PFV verbs \citep[e.g.][]{Paslawska.Stechow2003}, but there are counterexamples to this claim. Based on a corpus study, \citet{Borik.Gehrke2018} conclude that compositional IPFV PPPs with a predictable meaning exist, with both verbal and adjectival passive readings, but that they are quantitatively few and there are further restrictions.\footnote{Similar examples are mentioned in \citet{Schoorlemmer1995}, who nevertheless concludes that ``there is a finite set of imperfective verbs that occurs in syntactic passives in Russian. This set differs per speaker and reduces to zero for many.'' \citep[][226]{Schoorlemmer1995}.} In particular, in contemporary Russian, IPFV PPPs are only derived from simple but not from secondary IPFVs \REF{ex:storozppp} \citep[from][54]{Borik.Gehrke2018} and they cannot express a process reading \citep[see also][]{knjazev07}. 

\ea 
\label{ex:storozppp}
\gll Storož \minsp{\{} otkryval / otkryl\} vorota.   \\
	watchman.\textsc{nom} {} opened.\textsc{si} {} opened.\textsc{pfv} gates.\textsc{acc}\\
\glt	`A/The watchman was opening/opened a/the gate.'
\ea[]
{\gll Vorota byli otkryty storožem. \\
	gates.\textsc{nom} were open.\textsc{pfv}.\textsc{ppp} watchman.\textsc{instr} \\
\glt	`A/The gate was opened by a/the watchman.'}
\ex[*] 
{\gll 	Vorota byli otkryvany storožem.\\
 	  gates.\textsc{nom} were open.\textsc{si}.\textsc{ppp} watchman.\textsc{instr}\\ 
\glt	Intended: `A/The gate was opened by a/the watchman.'}
\z 
\z 

\noindent \citeauthor{Borik.Gehrke2018} hypothesise that IPFV PPPs can only have factual readings; this is arguably due to the strictly resultative meaning associated with Russian PPPs more generally (see also \citealt{Schoorlemmer1995}, who dubs this the ``Perfect Effect''). Relevant examples are given in \REF{pisano} \citep[from][58 \& 60]{Borik.Gehrke2018}.

\ea\label{pisano}
\ea
\gll Bylo pito, bylo edeno, byli slezy prolity. \label{pito}\\
	was drink.\textsc{ipfv.ppp.3sg.n} was eat.\textsc{ipfv.ppp.3sg.n} were tears pour.\textsc{pfv.ppp.pl}\\
\glt	`(Things) were drunk, (things) were eaten, tears were shed.' \\{} \hfill \textsc{existential}
\ex
\gll Pisano ėto bylo Dostoevskim v 1871 godu [\dots]\\ 
	write.\textsc{ipfv.ppp.3sg.n} that was Dostoevskij.\textsc{instr} in 1871 year \\
\glt	`That was written by Dostoevskij in 1871.'\label{pisanodost} \hfill \textsc{presuppositional}
\z 
\z 

\noindent In a recent corpus study on IPFV PPPs in Russian and Polish, \citet{wiemer+23} confirm the generalisations concerning the absence of a process reading and of SI PPPs for Russian. Taking into account also diachronic data, they show that SI PPPs and other IPFV readings cannot be found anymore in Russian texts from the 19th century onwards, so this is a rather recent development.\footnote{In contrast, \citet{wiemer+23} show that contemporary Polish still has SI PPPs and all IPFV readings are possible. We will see in \sectref{emp} that the same holds for Czech.}

As for Russian \textsc{reflexive passives}, it has been argued by \citet{fehrmann+10} that `by'-phrases are generally available with these. However, there is, again, the received view that there is an aspectual restriction, and this time it is commonly assumed that only IPFVs can derive reflexive passives. PFV counterexamples are discussed in the literature \citep[e.g.][]{Schoorlemmer1995,fehrmann+10}, but to my knowledge no systematic investigation into this issue has been conducted. The relevant reflexive counterparts to \REF{ex:storozppp}, including the availability of a `by'-phrase, are given in \REF{ex:storozrefl} \citep[from][54]{Borik.Gehrke2018}.

% \ea 
% \label{ex:storozrefl}
% \ea	
% \gll Vorota otkryvalis' storožem.\\
% 	gates.\textsc{nom} opened.\textsc{si}.\textsc{refl} watchman.\textsc{instr}\\
% \glt 	`The/A gate was (being) opened by a/the watchman.'\label{Olegom}
% \ex
% \gll \minsp{*} Vorota otkrylis' storožem.\\
% 	  {} gates.\textsc{nom} opened.\textsc{pfv}.\textsc{refl} watchman.\textsc{instr}\\ 
% 	\glt  
% \z 
% \z 

\ea 
\label{ex:storozrefl}
\ea[]{\gll Vorota otkryvalis' storožem.\\
	gates.\textsc{nom} opened.\textsc{si}.\textsc{refl} watchman.\textsc{instr}\\
\glt 	`The/A gate was (being) opened by a/the watchman.'\label{Olegom}}
\ex[*]{\gll Vorota otkrylis' storožem.\\
	  gates.\textsc{nom} opened.\textsc{pfv}.\textsc{refl} watchman.\textsc{instr}\\ 
	\glt  Intended: `The/A gate was (being) opened by a/the watchman.'}
\z 
\z 

\noindent Thus, we can draw the interim conclusion for Russian hat the functional IPFV-PFV distinction is defect with PPPs, and possibly also with reflexive passives. 

\subsection{Czech passives}

To my knowledge, there is nothing in the literature on Czech (verbal) passives that would suggest that the semantic role aspect plays is different than in actives. The literature on Czech reflexive and participial passives has examples with both aspects \citep[e.g.][]{veselovskakarlik,medovadiss,fehrmann+10,karlik17,karlik20,cahamedova20}. This could be due to the implicit assumption (as also suggested by Denisa Lenertová, Radek Šimík, p.c.) that (I)PFV works exactly the same with passives as it does with actives. This seems to be corroborated by my data investigation, to which I turn shortly. First, however, some general words about Czech passives. 

\subsubsection{Participial passives}
\label{sec:CZPPP}

Czech participial passives also combine a form of `be' with a PPP, so again, by form alone it seems impossible to distinguish a verbal from an adjectival passive. To my knowledge, there is no research on the formal semantics of Czech PPPs or on the types of `by'-phrases we might find, and the literature is primarily concerned with morphosyntactic issues. 

\citet{veselovskakarlik} propose that there are three derivational options for Czech PPPs: a) lexical derivation of an adjective, b) syntactic derivation of an adjective (in adjectival/stative passives), c) postsyntactic derivation of an adjective at PF (in verbal/eventive passives); arguments for adjectivisation come from the adjectival agreement on the PPP. The first two are argued to always appear in the long form and to be interpreted (at LF) as stative/adjectival (expressed by the feature [STATE]); the latter in turn is assumed to always involve a short form and to be interpreted as eventive/verbal (expressed by the feature [ACTIVITY]). %They argue that IPFV PPPs are always post-syntactic (expressing a verbal/eventive passive), but that PFV PPPs can be syntactic (in adjectival/stative passives) or post-syntactic. 
Evidence for the verbal vs. adjectival nature of the relevant PPPs comes from standard diagnostics, such as the availability of `by'-phrases, spatiotemporal and manner modifiers with verbal/eventive PPPs as opposed to adjectival `un'-prefixation and the (albeit limited) gradability of stative/adjectival PPPs.

The passive auxiliary is proposed to be inserted in \textit{v} and to be distinct both from the past auxiliary (which is inserted in T) and from copular/existential `be' (inserted in V). The passive auxiliary (like copular/existential `be' but unlike the past auxiliary) can combine with frequentative marking, which is argued to be the only functional Aspect marking in Czech ((I)PFV are treated as features on V); this marking can additionally appear on the PPP \REF{ex:PPPfreqa} \citep[adapted from][174]{veselovskakarlik}, but more often not on the PPP alone \REF{ex:PPPfreqb} \citep[adapted from][188]{veselovskakarlik}.\footnote{With past auxiliaries it is only the \textit{l}-participle that can bear frequentative marking (see op. cited for relevant examples). Recall from \sectref{sec:russpassives} that in Russian the passive auxiliary cannot bear frequentative marking, unlike its Czech counterpart, but like Czech past auxiliaries.}\textsuperscript{,}\footnote{In the examples from \citet{veselovskakarlik}, I keep the translations provided by the authors, but it is not entirely clear whether their use of progressive marking (or absence thereof) is supposed to reflect a difference in the aspectual or passive readings of the Czech examples, or whether the choice is not meaningful at all. For example, the progressive could be intended to either mark an imperfective reading (in any tense), or to mark an eventive/verbal (as opposed to stative/adjectival) passive reading (in the present tense). To avoid any bias caused by a potential English translation, I do not provide any translations in examples from Czech sources, e.g. those taken from \citet{karlik17}.} 

\ea \label{ex:PPPfreqa}
\ea 
\gll Já bývám \minsp{\{} chválen / chváliván\}. \\
 I am.\textsc{freq} {} praise.\textsc{ipfv.ppp} {} praise.\textsc{ipfv.freq.ppp} \\
\glt `I am being praised (repeatedly).'
\ex 
\gll Já jsem \minsp{\{} chválen / chváliván\}. \\
 I am {} praise.\textsc{ipfv.ppp} {} praise.\textsc{ipfv.freq.ppp} \\
\glt `I am praised (repeatedly).'
\z
\z

\ea \label{ex:PPPfreqb}
\ea[*] 
{\gll Diktát byl psáván každý pátek.  \\
dication.\textsc{nom} was written.\textsc{ipfv.freq.ppp} every Friday\\
\glt Intended: `Dictations were being written every Friday.'}
\ex[*] 
{\gll Ten test byl opisováván docela pravidelně.\\
this.\textsc{nom} test.\textsc{nom} was copy.\textsc{si.freq.ppp} rather regularly \\
\glt Intended: `This test was being copied rather regularly.'}
\z
\z 

\noindent Furthermore, the passive auxiliary can appear in past (\textit{byl(a/o/i/y)}), present (\textit{jsem} etc.), or future tense (\textit{budu} etc.).\footnote{To be precise, \citet{veselovskakarlik} argue that the future examples involve the future auxiliary in T in combination with an infinitival zero passive auxiliary in \textit{v}.}

\citet{veselovskakarlik} argue that Aspect is interpretable only on short form PPPs, as in, e.g., \REF{ex:czpppasp} \citep[adapted from][224]{veselovskakarlik}, and their discussion quite generally suggests that we are dealing with the typical aspectual readings we also find in the active.

\ea \label{ex:czpppasp}
\gll Pět vojaků bylo \minsp{\{} (z)raněno / (z)ra\v{n}ováno\}.\\
    five soldiers.\textsc{gen} was.\textsc{3sg.n} {} wound.\textsc{pfv.ppp.3sg.n} {} wound.\textsc{si.ppp.3sg.n}\\
\glt `Five soldiers were \{wounded / repeatedly wounded\}.'
\z

\noindent In the discussion of PFV short form PPPs in predicative position, \citet{veselovskakarlik} state that many Czech speakers view this form as archaic \citep[see also][96, for ``stative"-interpreted short form PPPs]{Biskup2019}. They furthermore argue that there is also variation in the tense options on the auxiliary, and that in general both the auxiliary and the PPP contribute to the overall temporal-aspectual interpretation of the sentence. In particular, they claim that past and future auxiliaries in combination with PFV short form PPPs tend to be interpreted eventively, while the present tense auxiliary is dispreferred with the short form \REF{ex:tenseppp} \citep[adapted from][225]{veselovskakarlik} (a long form is used instead; see op. cited for examples).\footnote{We will see shortly that PFV short form PPPs with a present tense auxiliary are discussed in the literature elsewhere, so it is not clear that the claim reported here is shared by everyone.}

\ea\label{ex:tenseppp}
\gll Okno \minsp{\{\textsuperscript{??}} je / bylo / bude\} otevřeno policií.\\
window.\textsc{nom} {} is {} was {} will.be open.\textsc{ipfv.ppp} police.\textsc{instr}\\
\glt `The window \{is/was/will be\} opened by the police.'
\z 

\noindent Moreover, the PFV short form expresses resultativity, with the result preceding or following the moment of speech, ``without inferring anything about the present state or action'' \REF{ex:pppresultative} \citep[adapted from][225]{veselovskakarlik}.

\ea\label{ex:pppresultative}
\ea
\gll Hala \minsp{\{} byla / bude\} uklizena špatně.\\
hall.\textsc{nom} {} was {} will.be clean.\textsc{pfv.ppp} badly\\
\glt `The hall \{was / will be\} cleaned badly.'
\ex
\gll Hala je uklizena špatně.\\
hall.\textsc{nom} is clean.\textsc{pfv.ppp} badly\\
\glt `The hall has (already) been cleaned badly.' (not: `is being cleaned')\label{ex:pppresultativeb}
\z 
\z 

\noindent In sum, according to \citet{veselovskakarlik}, all short form PPPs are unambiguously verbal/eventive (and adjectival only postsyntactically, at PF), with PFV PPPs in combination with a present tense auxiliary always being interpreted resultatively. The discussion of short form PPPs in other papers suggests that there might still be an ambiguity between adjectival/stative and verbal/eventive passive readings, at least for PFV PPPs. For example, \citet{cahamedova20} argue for the syntactic three-way syncretism of Czech PPPs in \REF{ex:nano}. 

\ea\label{ex:nano}
\ea Eventive passive: [Init [Proc Res]]
\ex `Intermediate' passive (R-state):  [Proc Res]
\ex Purely stative passive (T-state):  [Res]
\z 
\z 

\noindent Their account is spelled out in a nanosyntactic framework that incorporates \citeposst{kratzer00} distinction between resultant-state and target state participles (R-states and T-states in \REF{ex:nano}), and \citeposst{Ramchand2008} syntactic decomposition of events into maximally (an) Init(iator Phrase), (a) Proc(ess Phrase), and (a) Res(ult phrase). %To distinguish between these different types, \citeauthor{cahamedova20} employ standard tests for the eventive/verbal vs. stative/adjectival status (the terms are used interchangeably), such as (in)compatibility with adjectival morphology (e.g. adjectival `un'-prefixation, as well as comparative (degree) morphology; in the latter case all examples involve long form PPPs only though), temporal modification of the state vs. the event, (in)compatibility with `by'-phrases, modification by `recently'. 
For example, \REF{ex:pokoj} \citep[adapted from][119]{cahamedova20} is argued to be ambiguous between an adjectival and a verbal passive reading. 

\ea\label{ex:pokoj}
\gll Pokoj byl včera uklizen.\\
	room.\textsc{nom} was yesterday tidy.\textsc{pfv.ppp}\\
\glt	`The room was tidy/tidied yesterday.'\\
	(i) adjectival passive (the state of the room being tidy held yesterday) \\
	(ii) verbal passive (the event took place yesterday)
\z 

\noindent To argue for the distinction in \REF{ex:nano}, they employ similar tests as \citet{veselovskakarlik} for short vs. long form PPPs, only this time \citet{cahamedova20} contrast short form PPPs. For example, with adjectival PPPs (their ``purely stative'' PPPs) negation directly combines with the PPP \REF{ex:neukl}, whereas with verbal (eventive) PPPs sentential negation combines with the auxiliary `be' \REF{ex:nebyl} \citep[both examples adapted from][120]{cahamedova20}.

\ea\label{ex:ne}
\ea\label{ex:neukl}
\gll Pokoj byl včera neuklizen.\\
	room.\textsc{nom} was yesterday \textsc{neg}.tidy.\textsc{pfv.ppp}\\
\glt	`Yesterday, the room was untidy.'
\ex\label{ex:nebyl}
\gll Pokoj nebyl včera uklizen.\\
	room.\textsc{nom}  \textsc{neg}.was yesterday tidy.\textsc{pfv.ppp}\\
\glt	`The room was not tidied yesterday.'
\z 
\z 

\noindent Thus, I conclude that Czech (I)PFV short forms regularly appear in eventive/ver\-bal passives, and the limited discussion of the role of aspect in the literature suggests that we get regular (I)PFV readings. At least PFV short form PPPs can also get a stative/adjectival passive reading, so we will have to see whether this is restricted to PFVs and possibly also to the tense on the auxiliary. 

\subsubsection{Reflexive passives}

To my knowledge, it is not addressed in the literature whether the role of aspect in Czech reflexive passives is the same as in actives, so the discussion in this section will be brief and will only concern `by'-phrases. Unlike what we find with Czech participial passives, it has been argued that Czech reflexive passives do not allow for `by'-phrases \citep[e.g.][]{fehrmann+10,karlik17}. \citet{karlik17}, for example provides the contrast in \REF{ex:karlikby} (my glosses).


\ea \label{ex:karlikby}
\ea
\gll Škola je stavěna \minsp{(} zedníky).\\
	school.\textsc{nom} is build.\textsc{ipfv.ppp} {} mason.\textsc{instr.pl}\\
\ex
\gll Škola se stavi \minsp{(*} zedníky).\\
	school.\textsc{nom} \textsc{refl} builds.\textsc{ipfv} {} mason.\textsc{instr.pl}\\
\z 
\z 

\noindent From their cross-Slavic investigation of reflexive constructions \citet{fehrmann+10} conclude that there is variation in this respect: Belarusian, Bulgarian, Russian, Ukrainian, and Upper Sorbian allow for `by'-phrases with reflexive passives, whereas BCMS, Czech, Polish, Slovak, and Slovenian do not. For Czech and Russian this is illustrated in \REF{ex:fehrmby} \citep[adapted from][210f.]{fehrmann+10}.\footnote{The referential status of the `by'-phrases is not discussed at all; while the Czech example in \REF{ex:fehrmby} most likely has a referential/definite `by'-phrase, the Russian one could be interpreted non-referentially (or generically), in which case it might not be the best example for a `by'-phrase as a diagnostics for a verbal passive reading.}

\ea \label{ex:fehrmby}
\ea
\gll Šaty se právě šijí \minsp{(*} babičkou).\\
	dress.\textsc{nom.pl} \textsc{refl} right.now sew.\textsc{ipfv.3pl} {} grandmother.\textsc{instr}\\
\glt	`The dress is being made right now.' (`by'-phrase impossible) \hfill (Czech)
\ex
\gll Dom stroitsja \minsp{(} plotnikami).\\
	house.\textsc{nom} builds.\textsc{ipfv.refl} {} carpenters.\textsc{instr}\\
\glt	`The house is being built (by carpenters).' \hfill (Russian)
\z 
\z 

\noindent \citet{fehrmann+10} note further differences with reflexive constructions between Slavic languages in general, which for reasons of space I cannot go into, and this leads to their proposing two different types of reflexive morphemes in Slavic languages, which are in complementary distribution (some languages employ just one, others employ either but in different reflexive contexts): the argument-blocking one converts any argument into a semantically unbound variable, which allows further specification by a `by'-phrase if the external argument of a transitive verb is involved; the argument-binding one always binds the highest argument, which then obligatorily gets interpreted as arbitrary human so that further specification by a `by'-phrase is not possible. 

Even though reflexive passives in Czech are commonly assumed to be passives (and probably verbal passives), their incompatibility with `by'-phrases is suspicious (at least from a Germanic point of view) and might call into question the claim that we are really dealing with a ``true'' verbal passive and not with some kind of impersonal construction. On the other hand, the two native speaker consultants I worked with for the empirical investigation presented in the following section both note that reflexive passives are much more frequent than participial passives, in particular in colloquial use. %\footnote{Note also that \citet[][96]{Biskup2019} assumes that short form predicative PPPs in Czech are ``bookish or archaic''.} 
So for now I will assume that this is a passive construction, but this question deserves further investigation. 

\subsection{Empirical investigation: Aspect in Czech passives}
\label{emp}

In order to get a clear idea of whether the role of aspect in Czech passives is the same as in actives, I took examples discussed in \citet{karlik17}, as well as \REF{ex:ne} from \citet{cahamedova20}, and consulted two native speakers, asking them the following questions:\footnote{Thanks to Petra Charvátová and Denisa Lenertová; both are Moravian speakers (as is Karlík), so potential disagreement in judgments is probably not due to dialectal differences.}
Are reflexive and participial passives interchangeable? Are IPFV and PFV interchangeable? What happens when we switch to a different tense? Which (I)PFV readings does one get (process, general, habitual; completed, one-time, perfect/job done)? Since \citet{karlik17} is written in Czech, I added my own glosses but not translations; these are derivable from the speakers' judgments reported for the examples discussed here. 

Let us start with the judgments I obtained for \REF{ex:ne}%, the only example I took from \citet{cahamedova20}%, repeated in \REF{ex:nerep}
. 
%
%\ea\label{ex:nerep}
%\ea\label{ex:neuklrep}
%\gll Pokoj byl včera ne-uklizen.\\
%	room.\textsc{nom} was yesterday \textsc{neg}-tidy.\textsc{pfv.ppp}\\
%\glt	`Yesterday, the room was untidy.'
%\ex\label{ex:nebylrep}
%\gll Pokoj ne-byl včera uklizen.\\
%	room.\textsc{nom}  \textsc{neg}-was yesterday tidy.\textsc{pfv.ppp}\\
%\glt	`The room was not tidied yesterday.'
%\z 
%\z 
%
%\noindent 
The native speakers agreed that we are dealing with an adjective in \REF{ex:neukl}, and one added that it would be incompatible with a `by'-phrase. This is the same pattern we find in, e.g., German, for which it is commonly assumed that `un'-prefixation is incompatible with phrasal adjectivisation, which in turn is necessary for a `by'-phrase to be possible \citep[see, e.g.,][]{rapp96,gehrkeonli}. Furthermore, according to my consultants, the PFV \textit{uklizen} can be replaced by the IPFV \textit{uklízen} (a SI) in \REF{ex:nebyl}, but not in \REF{ex:neukl}. This could indicate that adjectival/stative short form PPPs, at least those that are not derived phrasally, are incompatible with a SI input or maybe even with an IPFV input in general, but this would have to be tested further.\footnote{The fact that the SI PPP cannot be the complement to `seem' \REF{ex:seemukl} \citep[adapted from][99]{Biskup2019} could also point in this direction. 

% \ea \label{ex:seemukl}
% \gll \minsp{*} Pokoj se zdá být uklízen.\\
% {} room.\textsc{nom} \textsc{refl} seems be.\textsc{inf} clean.\textsc{si.ppp}\\
% \glt Intended: `The room seems to be cleaned.'
% \z 

\ea[*]{\label{ex:seemukl}
\gll Pokoj se zdá být uklízen.\\
room.\textsc{nom} \textsc{refl} seems be.\textsc{inf} clean.\textsc{si.ppp}\\
\glt Intended: `The room seems to be cleaned.'}
\z 

} The meaning difference we get for \REF{ex:nebyl} if we change the aspect is that the tidying up was not started or finished (PFV), or that it was not started (IPFV), and this is what we would expect with a regular (I)PFV semantics.

Let us then turn to \citeposst{karlik17} examples and start with \REF{ex:chval}, which involves the IPFV activity predicate `praise'. 

\ea\label{ex:chval}
\gll Žák je chválen \minsp{(} učitelem).\\
	student.\textsc{nom} is praise.\textsc{ipfv.ppp} {} teacher.\textsc{instr}\\
\z 

\noindent My consultants reported that with the IPFV PPP in \REF{ex:chval} we get an ongoing or a regular reading; a perfect reading (job done), which would indicate that we are dealing with an adjectival PPP, was not available for this example. The IPFV PPP can also be exchanged by PFV \textit{pochválen} to give rise to a regular reading, or to a ``resultative perfect'' reading (Petr Biskup, p.c.; see also discussion around \REF{ex:pppresultativeb}). A habitual reading was stated to become available if the auxiliary is changed to frequentative \textit{bývat}, and this is possible with both IPFV and PFV PPPs (\textit{bývá (po)chválen}). %Examples with a frequentatively marked auxiliary can also be found in \citet{karlik17}, for instance \REF{ex:byva}.

%\ea\label{ex:byva}
%\gll Petr bývá často kritizován \minsp{(} svými přáteli).\\
%	Petr.\textsc{nom} is.\textsc{freq} often criticise.\textsc{ipfv.ppp} {} \textsc{possrefl.instr.pl} friend.\textsc{instr.pl}\\
%	\z 

%\noindent Corresponding examples in Russian are not available: While it is possible to add the frequentative marker to `be' (presumably as a copula), this is not possible for passive auxiliary `be' in Russian.\footnote{\citet{veselovska08} argues for three different types of `be' in Czech and shows that while the passive auxiliary, copula, and presentative `be' can be frequentatively marked, the perfect and subjunctive auxiliary `be' cannot.}  

Furthermore, without the `by'-phrase a reflexive passive is also possible in \REF{ex:chval} (even if the reflexive interpretation is more prominent), in both aspects and both present and past tense (\textit{se (po)chválí/-lil}). With the IPFV present we get an ongoing or general reading, with the PFV present either a future orientation or a generic reading (as part of a daily routine), with the IPFV past a process or general reading, and with the PFV past a one-time reading. Thus, based on the judgments concerning reflexive passives with the activity predicate in \REF{ex:chval}, (I)PFVs come with the same readings we also expect from active verb forms (as we have seen in \sectref{sec:crossslav}). With the participial passive in combination with present tense `be' we also seem to get typical IPFV meanings with the IPFV PPP (ongoing, regular), but a somewhat less typical picture with PFV PPPs (regular or ``perfect''). As we will see, this is different in all the other examples, so this might be due to the combination of an activity predicate and present tense `be'. 

Let us then turn to \REF{ex:delegaty}, which arguably involves an achievement predicate (or at least an accomplishment predicate), PFV `decide'. 

\ea\label{ex:delegaty}
\gll O tom bylo rozhodnuto \minsp{(} delegáty) včera.\\
	about that was decide.\textsc{pfv.ppp} {} delegates.\textsc{instr} yesterday\\
\z 

\noindent My consultants stated that such an example is odd with \textit{často} `often' (instead of `yesterday'), but with the frequentatively marked auxiliary \textit{bývalo} in combination with \textit{často} and the PFV PPP it is fully acceptable. Exchanging the PFV PPP with the IPFV \textit{rozhodováno} (a SI) has the effect that there was some deliberation but that it was not finished; with \textit{často} also a habitual reading is possible%, with or without \textit{bývalo}, and this could be due to the fact that this is a SI (recall discussion in \sectref{sec:hab})
. Finally, without the `by'-phrase the reflexive passive is acceptable in this context as well: with a PFV present we get a future orientation, for example in combination with the temporal adverbial \textit{zítra} `tomorrow' (\textit{se rozhodne zítra}), and with the PFV past a one-time interpretation (\textit{se rozhodlo \{včera / \textsuperscript{??}často\}});\footnote{As pointed out by Petr Biskup (p.c.), the reflexive PFV \textit{rozhodlo se} is also possible in the iterative interpretation, it is just that \textit{často} needs to have some element in its nuclear scope, which is why it is out in the example above.} with the IPFV past (\textit{se rozhodovalo}), we get either a process or a habitual interpretation (depending on the context and possibly additional adverbials). I conclude that we get the same (I)PFV readings we would get in active contexts.  

Let us move on to \REF{ex:pisese}, parts of which already appeared in \REF{ex:czpasstypes} and which involves an IPFV verb of creation.

\ea\label{ex:pisese}
\ea \label{ex:pisesea}
\gll Píše se stížnost. / Právě se píše stížnost.\\
	writes.\textsc{ipfv} \textsc{refl} complaint.\textsc{nom} {} just \textsc{refl} writes.\textsc{ipfv} complaint.\textsc{nom} \\
\ex\label{ex:piseseb}
\gll Je psána stížnost. / právě psaná stížnost\\
	is write.\textsc{ipfv.ppp} complaint {} just write.\textsc{ipfv.ppp.lf} complaint\\
	\z 
 \z 
 
\noindent Both examples allow for either a regular (without `just') or an ongoing (with or without `just') interpretation, again typical IPFV readings. Changing \REF{ex:piseseb} to PFV \textit{je napsána} leads to a finished-reading only (in the sense that the job is done). \citeposst{karlik17} second example in \REF{ex:piseseb} involves an attributive (and therefore long form) PPP, but my consultants agreed that a predicative short form PPP is possible as well (\textit{Právě je psána stížnost.} `right-now is written.\textsc{ipfv} complaint'); this time also the IPFV PPP gives rise to a job done-reading. As noted in \sectref{sec:passtypes} this is a typical adjectival passive reading, so here it is interesting that we get this kind of reading with the PPP in either aspect (where the IPFV is a simple one), even though in previous examples (e.g. \REF{ex:chval}) we did not get it with the (secondary) IPFV. \citet{karlik20} argues that the adverb `just' facilitates an adjectival/stative passive reading; in addition, maybe also the difference in word order might play a role: In \REF{ex:chval} we have a canonical subject-initial order, but in \REF{ex:pisese} the syntactic subject appears sentence-finally/postverbally. Whether word order plays a role for the kinds of readings we get with participial passives needs to be explored further.\footnote{In contrast, Petr Biskup (p.c.) cannot get a ``stative'' (job-done) reading with the IPFV PPP (\textit{Právě je psána stížnost.}) and states that it is only eventive for him, independently of the word order or the context.} Finally, it was reported that frequentative `be' in combination with a PFV PPP (\textit{bývá napsána}) was slightly odd but with the right context one could get the reading that events of this type are regularly finished.

Let us return to \REF{ex:karlikby}, repeated in \REF{ex:karlikbyrep}, which involves an IPFV verb of creation.

\ea\label{ex:karlikbyrep}
\ea
\gll Škola je stavěna \minsp{(} zedníky).\\
	school.\textsc{nom} is build.\textsc{ipfv.ppp} {} mason.\textsc{instr.pl}\\
\ex
\gll Škola se staví \minsp{(*} zedníky).\\
	school.\textsc{nom} \textsc{refl} builds.\textsc{ipfv} {} mason.\textsc{instr.pl}\\
\z 
\z 

\noindent Here, both IPFV passives get the ongoing reading as the first reading. With the right context, for example if the subject `school' appears in the plural, a habitual reading is also possible. With the PFV reflexive passive in the present tense we get reference to a result in the future. With the PFV participial passive and present tense `be', we get the job done-reading; so it might additionally be important to have present tense `be' for this adjectival passive reading to arise. The frequentative auxiliary in combination with the PFV PPP (\textit{bývá postavena}) is acceptable with a plural subject (`school buildings are regularly finished'). Finally, in PFV past, both reflexive and participial passives get the ``completed'' reading, and in the IPFV past, both passives get an ongoing reading. With plural subjects, a habitual reading is possible with both passives in the past and both aspects, so just as we saw in \sectref{sec:hab}, habituality seems to depend more on the context than on aspect, unless the auxiliary is marked for it. The unavailability of the `by'-phrase with the reflexive passive was confirmed for both aspects and all tenses. 
	
\citet{karlik17} notes that there are a few counterexamples to the claim that reflexive passives cannot combine with `by'-phrases, such as \REF{ex:reflbyok}. %Karlík: ``existují ale nemnohé protipříklady'' 

\ea\label{ex:reflbyok}
\gll V USA se prezident volí všemi občany.\\
	in USA \textsc{refl} president.\textsc{nom} votes all.\textsc{instr.pl} citizen.\textsc{instr.pl}\\
\glt	`In the US, the president is elected by all citizens.'
\z 

\noindent Only one of my consultants accepted this example, but stated that it did not feel equivalent to the PPP version with the instrumental NP; so maybe in \REF{ex:reflbyok} we do not have a true `by'-phrase but just the means by which an election is done; the consultant furthermore speculated that it might also be necessary to have `all' here and that it is stated like a rule. The same speaker observed that the IPFV past reflexive passive would also be ``kind of ok'' but not the PFV one. %(with the `by'-phrase). %With the IPFV PPP, we get the statement of a rule, with the PFV one that the election happened once.

In sum, Czech aspect fulfils its typical functions in both types of passives, as it does in actives% (unlike what we find in Russian)
: With the PFV we get a completed or finished reading with both types of passives. With the IPFV, we get an ongoing or regular interpretation. In the right context both (simple) IPFV and PFV passives can be interpreted as habitual, especially the participial one with the frequentatively marked auxiliary. As for adjectival passives it is possible that lexical target state participles might only be derived from PFVs (recall the discussion of \REF{ex:ne}) but that with the right context (which could be related to tense, adverbs, word order, or other factors) both IPFV and PFV participles can be interpreted at least as resultant state participles (giving rise to the job done-reading); this would have to be investigated further but would support the claims in \citet{cahamedova20} and go against the proposal by \citet{veselovskakarlik}, according to which all short form PPPs are eventive/verbal (in syntax proper). Finally, reflexive passives most likely do not allow for `by'-phrases. 

\subsection{Summary: Czech vs. Russian passives}

This section has shown that Czech and Russian differ in another domain, which has not received a lot of attention in the literature on cross-Slavic aspectual variation, namely passives. In Czech, both participial and reflexive passives are derived from all kinds of (I)PFVs (including SIs) and express the same (I)PFV readings we find with actives. PFV (short form) PPPs can also get a stative/adjectival reading. In contrast, Russian reflexive passives are restricted to IPFVs, whereas PPPs are regularly derived from PFVs, and there are severe restrictions on IPFVs: there are no SI PPPs, and we do not get the full range of IPFV meanings. In particular, there is no process reading, but maybe only factual readings. 

Thus, the category of aspect is fully functional in Czech passives, but defective in Russian passives. An interesting parallel can be drawn to nominalisations in Czech \textit{-ní/-tí} / Russian \textit{-nie/-tie}. These morphologically share the \textit{-n/-t}-component with PPPs, and they can also be regularly derived from both (I)PFVs in Czech, with regular (I)PFV readings, but only from one or the other in Russian (though this time also from SIs), with no predictable aspectual meaning \citep[see, e.g.,][]{Dickey2000,biskup23adv}. Data like these seem to suggest that Russian aspect is not fully functional in non-finite contexts, and this question is further explored in \citet{gehrkeadv23}. A formal-semantic account of cross-Slavic variation in passives, as well as in nominalisations, has to be left for future research.

In the following, I address existing formal-semantic accounts of cross-Slavic aspectual differences in other domains. These primarily focus on factual contexts, but (implicitly or explicitly) assume their accounts to be more general. %Let us then turn to the final section of this paper. I will first highlight an overarching difference between Czech and Russian in the interplay between Aspect and finiteness that emerges from some of the empirical differences outlined in this paper, to then return to the notion of temporal definiteness.  

\section{Accounting for the differences}
\label{sec:formal}

\noindent In this section, I address formal semantic accounts of cross-Slavic differences in aspect use. First, however, I discuss \citeposst{Dickey2000} cognitive semantic proposal that PFVs require temporal definiteness in Russian but not in Czech, an idea that is further explored more formally in \citet{Mueller-Reichau2018a,mueller23} (\sectref{sec:mueller}) and \citet{klimek22} (\sectref{sec:klimek}). \citeposst{gronndiss} influential account of Russian factual IPFVs, which was outlined in \sectref{sec:gronn}, is taken as a point of departure for two proposals dealing with differences between Czech and Russian, among others \citep[][]{alvestad13,Mueller-Reichau2018a,mueller23} (\sectref{sec:alvestad}--\sectref{sec:mueller}). Finally, in \sectref{sec:OFgehrke}, I call into question both the widespread assumption that Czech has existential IPFVs, as well as the analysis of factual IPFVs as ``fake'' %, rather than ``true'' 
IPFVs.   

\subsection{\citet{Dickey2000,dickey15}} 
\label{sec:dickey}

Based on data from the literature and native consultants, \citet{Dickey2000} describes differences in the use of (I)PFVs between ten Slavic languages in the following contexts%, some of which were already addressed in \sectref{sec:empirical}
: habituals, factuals, historical present, running instructions and commentaries, coincidence (performatives), sequences of events and ingressivity, verbal nouns. Some of these have already been presented in Table \ref{unterschiede} and described in detail in \sectref{sec:chains}--\sectref{sec:OFCZRU}. %; in table \ref{unterschiedeII}, I summarise the findings for the other contexts. 
%
%\begin{table}
	%\caption{Further aspectual differences between Czech and Russian described in \citet{Dickey2000}}
%	\label{unterschiedeII}
%\setlength{\tabcolsep}{8.2mm}
%		\begin{tabular}{l|ll}
%	  \lsptoprule
%			&\textbf{Czech} 	&\textbf{Russian}  \\ \hline 
				%Chains of single events (past)	&IPFV, PFV &(almost excl.) PFV\\
				%Iterativity, habituality (past, present)	&IPFV, PFV&(almost excl.) IPFV 	
				%\\
				%Historical Present	&IPFV, PFV&(almost excl.) IPFV	 \\
% 			Running instructions \& commentaries	&IPFV, PFV&(almost excl.) IPFV	 \\
%            Coincidence &IPFV/PFV &(almost excl.) IPFV	\\
%            Verbal nouns &aspectual pairs &aspectually neutral\\
%      \lspbottomrule
	%\end{tabular}
%\end{table}
%
%The first two contexts involve present tense forms in descriptions of situations that unfold while they are uttered or described. These involve, for example, recipes (running instructions), commentaries of sports games, and performatives (e.g. `I hereby declare ...'). In all these contexts, Czech can use PFVs quite freely, while Russian has to use the IPFV in most cases. This pattern is comparable to what we find in the historical present (recall \sectref{sec:HP}).  %which show a similar pattern as with the historical present. I will only discuss performatives here. While Czech can use the PFV for performatives with a range of verbs, also those that are not verbs of saying \REF{ex:perfcz}, Russian only allows it with a small subset of verbs of saying (e.g. `request', `admit', `note'), but not with others (e.g. `suggest', `swear'), and generally not with verbs that are not verbs of saying \REF{ex:perfru}. 
%
\iffalse 

\ea 
\ea \\
\gll\\
\glt `'\hfill (Czech)\label{ex:perfcz}
\ex A sejčas vam \minsp{\{*} podarju / darju \} {\.e}tot samovar na pamjat'.\\
\gll and now you.\textsc{dat} {} give.\textsc{1sg.pres.pfv} {} give.\textsc{1sg.pres.ipfv} samovar on memory\\
\glt `And now I present to you as this samovar as a keepsake.'\hfill (Russian)\label{ex:perfru}
\z 
\z
\fi 
%
%The third context involves deverbal nominalisations. While Czech has true aspectual pairs also in the domain of nominalisations, with the aspects contributing essentially the same meanings as with finite verb forms (e.g. IPFV ongoing vs. PFV ``completed'' event), Russian does not but derives them from one or the other aspectual partner, with no predictable aspectual meanings (see also \citealt{prochazkovama} for Czech, \citealt{Schoorlemmer1995,pazelskaya12} for Russian, and \citealt{biskup23adv,gehrkeadv23} for a comparison of the two); \citeauthor{Dickey2000} therefore calls the Russian forms ``aspectually neutral''. Some relevant examples are given in \REF{ex:nominals} \citep[see][ch. 9]{Dickey2000}. [NOT SURE ALL THIS IS NEEDED -- I MIGHT HAVE MISUNDERSTOOD PETR'S REVIEW]
%
%\ea e.g. `realise, execute' $>$ `realisation, execution'\label{ex:nominals}
%\ea pfv. provést / ipf. provádět  $>$ pfv. provedení / ipf. provádění			\hfill (Czech)
%\ex pfv. osuščestvit’ / ipfv. osuščestvljat’ $>$ osuščestvlenie/*osuščestvljanie		%\hfill (Russian)
%\z
%\z 
%
Given these differences, \citet{Dickey2000} argues that in the Slavic aspectual system of the western type, with its prototype Czech, the meaning of the PFV is totality, and the meaning of the IPFV quantitative temporal indefiniteness. In contrast, in the system of the eastern type, with its prototype Russian, the semantics of the PFV is proposed to be temporal definiteness and of the IPFV qualitative temporal indefiniteness. Following \citet{leinonen82}, \textsc{temporal definiteness} holds when a situation is ``uniquely locatable in a context, contiguous in time to qualitatively different states of affairs'' \citep[][19f.]{Dickey2000}, \textsc{quantitative temporal indefiniteness} involves ``assignability of a situation to several points in time'' \citep[][107]{Dickey2000}, and \textsc{qualitative temporal indefiniteness} ``the non-assignment of a situation to a single, unique point in time'' \citep[][108]{Dickey2000}. \citeauthor{Dickey2000} furthermore argues that Serbo-Croatian and Polish belong to transitional zones, tending towards the western or eastern type, respectively. 

\citeposst{Dickey2000} typology is further refined in \citet{dickey15}% \citep[see also discussion in][]{dickey18, fortuin2015typology, fortuin2015typology18}
, who characterises the distinctions between the western and eastern type as in Table \ref{tab:xslavEASTWEST}. 

\begin{table}
\caption{``Slavic East-West Division'' \citep[adapted from][]{dickey15}}
\label{tab:xslavEASTWEST}
 \begin{tabular}{lll}
 	  \lsptoprule
			&\textsc{WEST}: Czech, Slovak, &\textsc{EAST}: Russian,   \\ 
   			&Sorbian, Slovenian & Ukrainian, Belarusian  \\ \midrule [0.25ex] 
	Functional scope of PFV & maximal & minimal \\[0.25ex]
    IPFV general-factual & minimal usage & maximal usage \\[0.25ex]
    Productive delimitative \textit{po-} & no & yes\\[0.25ex]
    Productive distributive \textit{po-} & yes & no\\[0.25ex]
    Préverbe vide  & \textit{s-/z-} & \textit{po-}\\
        \lspbottomrule
 \end{tabular}
\end{table}

\citet{dickey15} argues that the western PFV has maximal functional scope because it is also used in present and past tense habitual contexts, the historical present, and running instructions, among others, which are all contexts which require the IPFV in the eastern type (recall also \sectref{sec:hab} \& \sectref{sec:HP}). Furthermore, \citeauthor{dickey15} claims that the western type makes minimal, the eastern type maximal use of factual IPFVs (recall \sectref{sec:OFCZRU}). Further distinctions between the two types are assumed to be in the choice of what \citeauthor{dickey15} labels ``préverbes vides'', a term for prefixes with a purely perfectivising function. In particular, he focuses on \textit{po-}, which is argued to be partially grammaticalised as such a perfectivising prefix in the eastern type, but to retain its surface-contact meaning in the western type, where its distributive use is also productive.\footnote{Relevant Russian examples are, for instance, delimitative \textit{pomolčal} `\textsc{po}-was-silent.\textsc{pfv}' in \REF{ex:po} and distributive \textit{po\-stu\-ča\-la} `\textsc{po}-knocked.\textsc{pfv}' in \REF{ex:sum}. %Differences between Czech and Russian \textit{po-} are also described in more detail in \citet{gehrke02,gehrke22}.
} The western type, in turn, is argued to employ \textit{s-/z-} as partially grammaticalised perfectivising prefixes \citep[on variation with \textit{po-} and \textit{s-/z-} see also][]{dickeyhutch, dickey05, dickey11}.\footnote{While \citet{Dickey2000} treats Bulgarian and Macedonian as belonging to the eastern type, he observes some deviations from this type in \citet{dickey15} in that the functional scope of PFV is argued to be narrower and the use of factual IPFVs greater, but not maximal.} 

\iffalse 
\citeposst{dickey15} transitional zones are illustrated in Table \ref{tab:xslavTRANS}. 

\begin{table}
\caption{``Slavic East-West Division'' II: Transitional \citep[adapted from][]{dickey15}}
\label{tab:xslavTRANS}
 \begin{tabular}{l|ll}
 	  \lsptoprule
			&\textbf{TENDING WEST}: &\textbf{TENDING EAST}:   \\ 
   			& BCMS & Polish (standard)   \\ \hline 
	Functional scope of PFV & greater  & relatively greater \\
    IPFV general-factual & lower usage & [not indicated]  \\
    Productive delimitative \textit{po-} & no & yes \\
    Productive distributive \textit{po-} & yes & yes \\
    Préverbe vide     & -- & \textit{s-/z-}  \\
	      \lspbottomrule
 \end{tabular}
\end{table}

\fi

\citet{dickey15} correlates the cross-Slavic aspectual variation with the diachronic developments and different contact situations of the languages involved: Whereas the western type had maximal contact with German, the eastern type had varying levels of Finno-Ugric language contact. The transitional zones, in turn, had either contact with German and East Slavic (Polish), or at various levels with German, Romance, and some also with Turkish (BCMS). 

The notion of temporal definiteness is not formally defined in Dickey's works, and this raises several questions: Are we dealing with definiteness of the event (\citeauthor{Dickey2000} speaks of ``situations''), or with definiteness of a particular time (e.g. the event time, or the reference time)? What kind of definiteness is involved: familiarity (anaphoricity), uniqueness (situational or world knowledge)? Dickey speaks of uniqueness, but other formal notions could be explored. For example, it is also possible that we are dealing with specificity, or with determinacy, which in the nominal domain just means that the nominal denotes an individual of type $e$ \citep[see][]{coppockbeaver15}. All these are questions that one would have to explore if one wanted to formalise Dickey's notion of temporal definiteness. The few works aiming at formalising this, to which I turn now, have so far only scratched the surface and exploring various avenues in this respect can be fruitful. %Let us then turn to formal-semantic accounts of some of the differences between particularly Czech and Russian,.

\subsection{\citet{alvestad13}}
\label{sec:alvestad}

\citet{alvestad13} investigates the use of (I)PFVs in imperatives in twelve Slavic languages, using the ParaSol corpus \citep[][]{parasol} \citep[see also][]{waldenfels12%,alvestad14
}. %\footnote{Note that Serbian and Croatian count as two separate languages in this corpus.} 
The cross-Slavic variation in aspect use she observes in imperatives is illustrated in \REF{ex:alvimp} \citep[see][312]{alvestad13}.

\eanoraggedright\label{ex:alvimp}
Choice of IPFV in imperatives \\
Russian (60\%) $>$ Belarusian (59\%) $>$ Ukrainian (58\%) $>$ Bulgarian (48\%) $>$ Polish (47\%) $>$ Serbian, Croatian (45\%) $>$ Macedonian (44\%) $>$ Upper Sorbian (43\%) $>$ Slovak (33\%) $>$ Czech (31\%) $>$ Slovenian (29\%) 
\z

\noindent She argues that the use of IPFVs in imperatives is a kind of general-factual use, and she builds on \citeposst{gronndiss} proposal for Russian factual IPFVs, as well as \citeposst{gronnstechow10} research programme (which is further developed in \citealt{gronnstechow16}, as outlined in \sectref{sec:gronn}). Her general proposal is summarised in Table \ref{tab:alv} \citep[cf.][229]{alvestad13}.

% \begin{table}
% 	\caption{Slavic Tense and Aspect \citep[see][229]{alvestad13}}
% 	\label{tab:alv}
% 		\begin{tabular}{ll}
% 	  \lsptoprule
% 		\textsc{unique-def. T \& indef. Asp} 	&\textsc{indef. T \& indef. Asp}  \\ 
%         Prediction: &Prediction:\\
%         Morphological PFV & Aspect neutralisation, \\
%         &existential fake IPFV\\\hline
% 		\textsc{def. T \& def. Asp} 	&\textsc{indef. T \& def. Asp}  \\ 
%         Prediction: &Prediction:\\
%         Aspect neutralisation, & Aspect neutralisation, \\
%         presuppositional fake IPFV & presuppositional fake IPFV\\        
% 	      \lspbottomrule
% 	\end{tabular}
% \end{table}

\begin{table}
	\caption{Slavic Tense and Aspect \citep[see][229]{alvestad13}}
	\label{tab:alv}
		\begin{tabularx}{.85\textwidth}{XX}
	  \lsptoprule
		\textsc{unique-def. T \& indef. Asp} 	&\textsc{indef. T \& indef. Asp}\\[0.25ex] 
        ~~~Prediction: &~~~Prediction:\\
        ~~~Morphological PFV &~~~Aspect neutralisation, \\
        &~~~existential fake IPFV\\[2ex]
		\textsc{def. T \& def. Asp} 	&\textsc{indef. T \& def. Asp}\\[0.25ex] 
        ~~~Prediction: &~~~Prediction:\\
        ~~~Aspect neutralisation, &~~~Aspect neutralisation, \\
        ~~~presuppositional fake IPFV &~~~presuppositional fake IPFV\\
	      \lspbottomrule
	\end{tabularx}
\end{table}


In particular, \citet{alvestad13} argues that the PFV is always indefinite, in the sense that it introduces a new eventive discourse referent. IPFV, on the other hand, is argued to be compatible with both (in)definite events. Similarly, Tense can involve a definite or indefinite reference time (new vs. an\-a\-phor\-ic or unique). In Table \ref{tab:alv}, ``unique-def'' involves uniqueness, whereas the other instances of ``def'' involve anaphoric definiteness.

Inspired by \citeposst{gronnstechow10} research programme, she revises their Aspect Neutralisation Rule in \REF{ex:aspneugst} as follows:

\eanoraggedright \textsc{Aspect Neutralisation Rule} (revised) \citep[see][230]{alvestad13}\label{ex:aspneu}
\eanoraggedright When a semantically perfective aspect is definite/anaphoric, it is morphologically neutralised to IPFV. This holds irrespective of whether the tense is indefinite or definite. When this rule is adhered to, we see an instance of the presuppositional type fake IPFV.
\ex When a semantically perfective aspect is indefinite AND the tense is indefinite%, (the extended future in the case of imperatives)
, the aspect is morphologically neutralised to IPFV. When this rule is adhered to, we see an instance of the existential type fake IPFV.
\z
\z 

\noindent Her account of aspect use differences is primarily pragmatic as she argues that ``Slavic languages adhere to the Aspect Neutralization Rule to varying degrees'' and views ``Russian as the most `law-abiding' language'' \citep[][312]{alvestad13}. 

The proposal by \citet{alvestad13}, as well as the one by \citet{gronn15,gronnstechow16}, raise several questions. How is the system with covert (in)definite operators restricted, i.e. when do we get one and when the other?  Why is the reference time with presuppositional IPFVs necessarily definite, as argued by \citet{gronnstechow16}, whereas it can be either for \citet{alvestad13}? The reference time could also be treated as indefinite with each clause introducing a new reference time into the discourse, as argued by, for instance, \citet{gehrketrueipf} (I will come back to this in \sectref{sec:OFgehrke}). Furthermore, why do PFVs always involve indefinite events \citep{alvestad13}?  In \citet{Mueller-Reichau2018a}, which I address in the following section, PFV events are always definite (unique). Why is only Russian ``law-abiding'' \citep{alvestad13}, and what exactly triggers the choice of one or the other aspect in languages with more optionality, such as Czech? Finally, two different notions of definiteness play a role in \citeposst{alvestad13} account: uniqueness of Tense, i.e. the reference time (requiring the PFV, in case Aspect, i.e. the event, is indefinite), as opposed to anaphoricity of Tense/reference time and Aspect/the event (leading to the use of a presuppositional IPFV, no matter whether Tense/the reference time is definite or indefinite); these are quite different notions of definiteness and it would be good to figure out when and why we use one or the other. 

\subsection{\citet{Mueller-Reichau2018a,mueller23}}
\label{sec:mueller}

\citet{Mueller-Reichau2018a} provides a semantic account of cross-Slavic differences in existential factual contexts, in which Polish and Czech use the PFV but Russian the IPFV, mainly discussing data from the literature and consulting speakers for additional data points. Building on \citet{gronndiss}, he treats the IPFV as denoting an underspecified relation between E and R in all three languages \REF{ex:muelleripfv}.

\ea\label{ex:muelleripfv} 
\sib{IPFV} = $\lambda P \lambda t \exists e [P(e) \wedge e \bigcirc t]$
\z 

\noindent He proposes that the differences between the three languages lie in the semantics of the PFV: In all three, PFVs involve event uniqueness, but in Russian there is an additional requirement of target state validity, see \REF{ex:muellerpfv}.

\ea\label{ex:muellerpfv}
\ea \sib{PFV\textsubscript{POL/CZ}} = $\lambda P \lambda t \exists e [P(e) \wedge e \subseteq t \wedge \neg \exists e' [P(e') \wedge e' \not = e]]$
\ex \sib{PFV\textsubscript{RU}} = $\lambda P \lambda t \exists e [P(e) \wedge e \subseteq t \wedge \neg \exists e' [P(e') \wedge e' \not = e] \wedge f\textsubscript{\textsc{end}}(t) \subseteq f\textsubscript{\textsc{target}}(e)]$
\z 
\z 

\noindent The denotation of the PFV requires that the event time is part of or equal to the reference time ($e \subseteq t$) in all three languages, which, as we saw, is a cross-linguistically common account of the semantics of PFVs. However, there is an additional requirement for the event to be unique, which is expressed in the second conjunct for both PFV operators in \REF{ex:muellerpfv} (there is no event $e'$, of which the property holds as well and which is not identical to $e$). For the Russian PFV, \citeauthor{Mueller-Reichau2018a} adds another conjunct ($f\textsubscript{\textsc{end}}(t) \subseteq f\textsubscript{\textsc{target}}(e)$), which is meant to capture the intuition of target state validity, as it ``requires the reference time to end when the target state is in force'' \citep[][305]{Mueller-Reichau2018a}.\footnote{\citet{Mueller-Reichau2018a} argues, following \citet{mittwoch08}, that target state validity by itself already requires uniqueness and completion so that the first two conjuncts in the definition of the Russian PFV could be omitted. For ease of comparison I leave them here though.} \citeauthor{Mueller-Reichau2018a} argues that the requirement of target state validity prevents the Russian PFV to occur in any existential context (including contexts in which the result state has been reversed), whereas its absence allows Czech and Polish to use the PFV in these contexts. 

While the differences between Czech and Russian (and possibly other languages) might very well be captured by a difference in the conditions on the PFV, while maintaining the same semantics for the IPFV, there is an empirical problem with the proposal that the PFV in both languages requires event uniqueness. Czech quite regularly employs PFVs in contexts in which the events are not unique, namely in iterative and habitual contexts, as we saw in \sectref{sec:hab}. Furthermore, the notion of target state validity should only apply to predicates that already come with a target state \citep[in the sense of, e.g.,][]{kratzer00}, which are predicates that involve changes of states (accomplishments and achievements). But what about predicates without target states, which in Russian can be prefixed by, e.g., delimitative \textit{po-} or ingressive \textit{za-} to become PFV (recall discussion in \sectref{sec:chains})? %I believe \citeposst{Mueller-Reichau2018a} account is also problematic for Czech existential contexts, as hinted at in \sectref{sec:OF}, as it can be argued that the use of Czech IPFVs in existential contexts can be derived from its general use also in other contexts (despite event uniqueness). I will come back to this line of argumentation in \sectref{sec:OFgehrke}. 

\citeposst{Mueller-Reichau2018a} proposal is further refined in \citet{mueller23}, where he adds other contexts as well as colloquial Upper Sorbian (CUS) data discussed in the literature. Like Czech and unlike Russian, CUS can use the PFV in iterative and habitual contexts (e.g. with `often'), as illustrated in \REF{ex:cushab} \citep[attributed to][]{breu00}.

\ea\label{ex:cushab}
\gll Wón je husto jenož jednu knihu předał.\\
	he \textsc{aux} often only one book sold.\textsc{pfv}\\
\glt 	`He often sold only one book.' \hfill (Colloquial Upper Sorbian)
\z 

\noindent Unlike both Czech and Russian, CUS can also, rather surprisingly, use the PFV to describe an ongoing event \REF{ex:cusproc} \citep[attributed to][]{breu00}.


\ea\label{ex:cusproc}	
\gll Jurij jo rune jen text šełožił, hdyž sym ja nutř šišoł.\\
Jurij \textsc{aux} now a text translated.\textsc{pfv} when \textsc{aux} I in came.\textsc{pfv}\\
\glt	`When I came in, Jurij was translating a text.' \hfill (Colloquial Upper Sorbian)
\z 

\noindent To account for the differences in aspect use between the three languages, \citet{mueller23} replaces his previous definition of event uniqueness by determinacy. In particular, he builds on the notion of a path, as used in, e.g, \citet{Krifka1998,zwartsprepasp,Gehrke2008}, with the relevant notion of a determinate event predicate, which involves unidimensionality, directedness, and boundedness %, informally characterised as in \REF{ex:path} 
\citep[for the formalisation of these notions see][]{mueller23}.

\iffalse 
\ea\label{ex:path}	
An event predicate $P$ is determinate iff it is unidimensional, directed and bounded, whereby:
\ea $P$ is unidimensional iff the events in its denotation set share a path structure such that all elements of this path structure are parts of a common path which likewise belongs to the part structure.
\ex $P$ is directed iff the events in its denotation set share a path structure such that there are no two non-overlapping elements in this path structure that occupy the same space.
\ex $P$ is bounded iff the events in its denotation set share a path structure which includes an element that cannot be concatenated by another element of the same path structure such that the resulting path likewise belongs to this path structure.
\z 
\z 
\fi 

%\noindent 
\citet{mueller23} employs the same weak semantics for IPFVs in all three languages that we saw in \REF{ex:muelleripfv}. What is new now and what directly translates the idea that the IPFV is the unmarked member of the aspectual opposition, is that this semantics is also part of the semantics of the marked PFV operators, which adds further requirements in conjuncts \REF{ex:muellerpfvnew}. 

\ea\label{ex:muellerpfvnew} 
\ea\label{ex:muellerpfvnewcus} 
\sib{PFV\textsubscript{CUS}} = $\lambda P \lambda t \exists e.P(e) \wedge \cnst{det}(P) \wedge t \bigcirc \tau(e)$
\ex\label{ex:muellerpfvnewcz} 
\sib{PFV\textsubscript{CZ}} = $\lambda P \lambda t \exists e.P(e) \wedge \cnst{det}(P) \wedge t \bigcirc \tau(e) \wedge f\textsubscript{\textsc{end}}(\tau(e)) \subseteq t$
\ex\label{ex:muellerpfvnewru}
\sib{PFV\textsubscript{RU}} = $\lambda P \lambda t \exists e.P(e) \wedge \cnst{det}(P) \wedge t \bigcirc \tau(e) \wedge f\textsubscript{\textsc{end}}(\tau(e)) \subseteq t \wedge f\textsubscript{\textsc{end}}(\tau(e)) \subseteq f\textsubscript{\textsc{target}}(e)$
\z
\z 

\noindent In all three languages, the event predicate has to additionally involve a determinate path ($\cnst{det}(P)$), in Czech and Russian the end of the temporal trace has to additionally be included in the reference time, and only in Russian, target state validity is also necessary, building on \citet{Mueller-Reichau2018a}.  

While this proposal solves the issue of event uniqueness for Czech (and presumably also CUS), by replacing it with determinacy, it is now less clear why Russian PFVs require single events. In addition, the question about target state validity remains, and further questions arise. One would be how to integrate Russian ingressives (recall the discussion of the prefix \textit{za-} in \sectref{sec:chains}), which are PFV but with which it is not the end but the beginning of the temporal trace that is picked up by the ingressive prefix. Furthermore, the only difference between PFV and IPFV in CUS is the nature of the event path, which is argued to have to be directed, unidimensional, and bounded with PFVs. However, the idea of directed, unidimensional, bounded paths playing a role for the interpretation of events is commonly thought of as characterising the predicates as telic \citep[see, e.g.,][]{Krifka1998}, which in turn belongs to the level of inner, rather than outer aspect. Is it the case, then, that the CUS (I)PFV distinction is an inner-aspectual one, as also indicated by the discussion in \citet{breu00} and \citet{scholze08}, who argue informally that CUS aspect is based on ``terminativity'' (recall fn. \ref{fn:innerouter})? %This might also not be true, since we still get the imperfective paradox etc. 

\subsection{\citet{klimek22}}
\label{sec:klimek}

In her account of the differences in aspect use in factual contexts between Czech, Polish, and Russian, \citet{klimek22} builds on \citet{Ramchand2004,ramchandlingua}, in assuming that the ``First phase syntax'' ($\sim$ \textit{v}P/VP; see also \sectref{sec:CZPPP}) introduces an event variable, and grammatical aspect introduces a time variable, which is crucially an instant, rather than an interval (as otherwise commonly assumed); the event variable and the temporal variable are related by the temporal trace function $\tau(e)$. In Ramchand's proposal, Russian PFV events introduce a definite reference time R, which is a  ``a single unique moment'' \citep[][]{Ramchand2004} or ``a specific moment'' \citep[][]{ramchandlingua}. With accomplishments (Ramchand's procP/resP-syntax), R must be within both subevents (process and result state), and since it is an instant it has to be the transition itself. Russian IPFV events, in contrast, are argued to introduce an indefinite reference time, i.e. an arbitrary moment within the temporal trace of the event; with accomplishment predicates this would be an arbitrary time within the process. Finally, Tense is argued to bind the time variable and to relate it to S. 

\citet{klimek22} proposes more generally that R can be (in)definite either with respect to the temporal trace of the event (at the micro-level), or with respect to the utterance time (at the macro-level), building on earlier work \citep[][]{klimek20}. %, she argues that with complex events (i.e. accomplishments), the placement of the temporal variable with respect to the temporal trace of an event depends on whether the focus is more on the initiation, process, or result. With focus on the result, focus is on the transition (the change of state or location), so that we get definiteness with respect to the temporal trace of the event. With focus on the initiation or process of the event, we get an arbitrary point and therefore indefiniteness with respect to the event's temporal trace. 
Since with presuppositional IPFVs, the result is assumed to be presupposed, \citet{klimek20} argues that the resultee is part of the conversation and event completion is inferred. %[that doesn't really make sense -- in the imperfective paradox the resultee is also present, so we might as well have a process reading]
In the case of existential IPFVs that combine with `once' or `ever', for instance, she argues that we are dealing with indefiniteness with respect to the utterance time. More generally, the constellation in existential factual contexts leads to an aspectual competition, in which speakers can choose to go for the PFV because there is definiteness with respect to the temporal trace (at AspP), or for the IPFV because there is indefiniteness with respect to the utterance time (at TP). \citet{klimek20} assumes that both options are in principle available, because the verb form competes for lexical insertion at the CP level. %She observes variation within Polish in this respect: Whereas in Western Polish there is a preference to choose definiteness with respect to the temporal trace/AspP and thus a higher use of the PFV in existential contexts, in Eastern Polish there is a preference to express definiteness with respect to the utterance time/TP and a higher use of IPFVs is observed.

%With the additional data from Czech and Russian as well as from different types of presuppositional IPFV contexts, 
\citet{klimek22} argues that temporal (in)definiteness at the macro-level should be understood in terms of temporal specificity, but throughout the paper she uses these terms interchangeably. She proposes that %In the case of %neutral 
existential contexts involve %, she proposes that we are dealing with 
temporal indefiniteness at the macro-level and underspecification for definiteness at the micro-level. This leads to the following cross-Slavic variation. In Western Polish and Czech, we find a preference for definiteness with respect to the temporal trace/AspP with accomplishments, and this is obligatory with achievements. This is argued to be so because achievements are instantaneous and the time variable can only be located at the time instant at which the change of state happens. In Eastern Polish and Russian, on the other hand, there is a preference for definiteness with respect to the utterance time/TP, and in some cases this is obligatory in Russian (e.g. with `ever', recall discussion in \sectref{sec:OF}). It is argued that the Russian (but not the Polish or Czech) PFV has to be anchored to a specific temporal location on the timeline. With presuppositional contexts, there is a relation to the earlier discourse and we are dealing with ``pragmatic specificity''. Therefore the PFV is also possible in Russian, but more so in Czech and Polish, and it is again a matter of speakers' choice which aspect is used.\footnote{This proposal is reminiscent of \citeauthor{petruxina00}'s (\citeyear{petruxina00}: 59f.) idea that aspect choice is conditioned by two, sometimes conflicting factors: the ``objective'' nature of the event (i.e. completed or not) vs. the ``subjective'' assessment of the temporal contour of the event by the speaker, which gets resolved differently in, e.g., Czech and Russian.} 

\citeposst{klimek22} account raises a number of questions, some of which have already been raised for the other accounts. If it is a matter of speakers' choice, what regulates the choice? In addition, if it is a matter of speakers' choice, it seems to be a pragmatic account, but at various points it is stated that in some cases the use of a particular aspect is obligatory. Shouldn't obligatoriness be reflected in the semantics of (I)PFV? It is also not clear whether we are dealing with definiteness (and then uniqueness or anaphoricity) or with specificity? Finally, if it is true that with presuppositional IPFVs, the result is presupposed, the resultee is part of the conversation, and event completion is inferred, doesn't that come close to accomplishments under a process reading, in which, granted, event completion is \textit{not} inferred, but at least it is conceptually given and it is equally not in focus (as \citeauthor{klimek22} argues for presuppositional IPFVs without focus on the result)? Considerations like these support the idea presented in \citet{gehrke22} that presuppositional IPFVs are but a subcase of process/durative IPFVs. I address this proposal in the following section, in which I also call into question whether Czech has existential IPFVs to begin with.

\subsection{\citet{gehrke22,gehrketrueipf}}
\label{sec:OFgehrke}

Based on the corpus study in \citet{gehrke02}, I concluded in \citet{gehrke22} that aspect use in the description of single vs. repeated (e.g. habitual) events in Czech is almost identical, except for a few verb forms that are specialised to mark entire passages as habitual (recall discussion in \sectref{sec:hab}). Therefore the Czech IPFVs that are used in these contexts are not motivated by the fact that we are dealing with non-single events, but appear for other reasons, e.g. focus on the process, duration or state. This is different in Russian, where many verb forms in habitual contexts are additionally imperfectivised, so that the IPFV in such contexts is primarily motivated by the fact that we are not dealing with a single event. 

If, in addition, \citet{paduceva96} is correct in assuming that the use of Russian IPFVs in existential contexts arises because we are dealing with potential repeatability, then the IPFV appears for the same reason it appears in habitual contexts: we are dealing with (potentially) plural events. This reasoning leads to the hypothesis in \citet{gehrke02,gehrke22} that the existential IPFV is but a special case of the use of IPFVs in the description of non-single events. Coupled with the observation above that Czech does not use the IPFV to signal habituality (instead it is motivated by other IPFV readings), this leads to the assumption that Czech lacks existenial IPFVs, contra the received view. Or, in other words, the use of Czech IPFVs in existential contexts is not motivated by the existential context. For example, in the discussion of \REF{ex:re} in \sectref{sec:OFCZRU}, which \citet{klimek22} assumes to involve an existential IPFV also in Czech, I argued that the Czech IPFV could instead be motivated by the process use of IPFVs, as we are dealing with the question whether the process of watering has taken place. Given that many of her Czech examples involve incremental theme verbs, which I take to be activity predicates by themselves, following \citet{rappaportlevin10}, \citet{kennedyincr}, it is empirically difficult to distinguish the two views. We would need to find clear resultative examples instead, but I do not see such Czech IPFV examples discussed in the literature.\footnote{Petr Biskup (p.c.) suggests that in the following exchange between a waiter (W) and a customer (C), the IPFV in the customer's reply is existential:

\ea\label{ex:waiter}
\begin{xlist}
\exi{W:}{\gll Mám tady nějaká piva.\\
have.\textsc{1sg} here some beers\\
\glt `I have some beers here (that still need to be paid).'}
\exi{C:}{\gll To pivo jsem platil.\\
\textsc{dem} beer.\textsc{acc} \textsc{aux.1sg} paid.\textsc{ipfv}\\
\glt `I've paid this beer already.'}
\end{xlist}
\z 

\noindent However, I am not convinced that this is an existential use (`this beer was involved in a paying event'), since it could also be analysed as presuppositional (`with respect to this beer, the paying event (retrievable from the discourse) happened'). We would have to analyse the prosody and information structure to resolve this issue.}

Things are different with the presuppositional IPFV. Under the analysis of presuppositional IPFVs in \citet{gehrketrueipf}, their use leads to a zooming in on the reference time of a previously introduced or contextually retrievable event; thus, the IPFV in such contexts involves a partitive (a ``true") IPFV semantics, contra \citet{gronn15}. This analysis makes presuppositional IPFVs just a special case of the use of IPFVs for ongoing events, as hypothesised in \citet{gehrke22}. Since both Czech and Russian use the IPFV to describe ongoing events, both languages are expected to use the IPFV in presuppositional factual contexts. 

Let me illustrate the account of Russian presuppositional IPFVs as ``true'' imperfectives with the example in \REF{ex:zaplatili} \citep[adapted from][]{gehrketrueipf}. 

\ea \label{ex:zaplatili}
\gll Zaplatili. \textit{Plačeny} byli naličnymi šest' tysjač rublej.\\
	paid.\textsc{pfv.3pl} pay.\textsc{ipfv.ppp} were in.cash six thousand Rubles\\
\glt	`They paid. It was paid 6.000 Rubles in cash.'
\z 

\noindent In this example, the PFV verb form \textit{zaplatili} `(they) paid' in the first sentence introduces a ``completed'' paying event. The presuppositional IPFV PPP \textit{plačeny} `paid' in the second sentence links back to this already introduced event. The marked word order and the most natural way to read this example indicate a marked information structure, a hallmark of the presuppositional IPFV: The paying event appears in the beginning of the (second) sentence and is backgrounded, focus lies on the sentence-final subject and (probably also on) the modifier (`6.000 Rubles (in cash)'). The DRT analysis \citep[][]{Kamp1993} of this example proposed in \citet{gehrketrueipf} is given in \REF{ex:zaplatilib}.\footnote{A linear notation for discourse representation structures (DRSs) is employed, where discourse referents are written on the left-hand side, before $|$ (in a traditional DRS they appear at the top of the DRS), and the conditions on these discourse referents are listed to the right of $|$, separated by commas (which in a different notation can be translated as conjuncts).}

% \ea	\label{ex:zaplatilib}
% $[e_1, e_2, t_1, t_2, n, x\,|\,\textbf{pay}(e_1), \tau(e_1) \subset t_1, t_1 < n, \textbf{pay}(e_2), \textsc{theme}(e_2, x),$\\ $\textbf{6.000R}(x), \textbf{in-cash}(e_2), e_2 = e_1, t_2 \subset \tau(e_2), t_2 < n]$
% \z 

\ea	\label{ex:zaplatilib}
$[e_1, e_2, t_1, t_2, n, x\,|\,\textsc{pay}(e_1), \tau(e_1) \subset t_1, t_1 < n, \textsc{pay}(e_2), \cnst{theme}(e_2, x),$\\\ $\textsc{6.000R}(x), \textsc{in cash}(e_2), e_2 = e_1, t_2 \subset \tau(e_2), t_2 < n]$
\z 

\noindent Under this analysis, the first sentence introduces a new eventive discourse referent $e_1$, and its run time, $\tau(e_1)$, is included in the reference time $t_1$ (PFV semantics), which is before $n$(ow) (past tense semantics). The presuppositional IPFV in the second sentence introduces a new paying event $e_2$, which -- due to the information structural cues -- is anaphorically linked to $e_1$, i.e. $e_2 = e_1$. The new information in focus is about $e_2$, and since $e_2$ is identical to $e_1$ it is also about $e_1$: the theme of $e_2$ is `6.000 Rubles' and this was paid `in cash' (an event modifier). The IPFV semantics specifies that there is a new reference time, $t_2$, that is included in the run time of $e_2$, $\tau(e_2)$; past tense indicates that $t_2$ is before $n$. 

How does this analysis still capture the intuition that the paying event is completed, if the presuppositional IPFV is analysed as involving IPFV semantics? I argue that event completion information is already given in the first sentence about $e_1$ (its run time is included in the first reference time $t_1$). Since $e_2$ equals $e_1$, the actual event of paying remains completed. Furthermore, the second reference time, $t_2$, is included in the run time of $e_2$, and therefore it is also included in the run time of $e_1$ (since $e_2$ is identical to $e_1$). By transitivity, $t_2$ must also be included in the first reference time, $t_1$. The effect of the presuppositional IPFV, then, is that it is used to zoom in on a narrower reference time within the bigger reference time introduced in the first sentence; the link between the two reference times $t_1$ and $t_2$ is only indirect, via the events involved, but it can still be made. The assertion that the sentence with the presuppositional IPFV makes, then, is only for part of the bigger reference time (and therefore also only for part of the actual event), and this is what is captured by the IPFV semantics. %In other words, the event description provided by the presuppositional IPFV elaborates on the first event% \citep[on elaboration as a discourse relation between events, see, e.g.][]{lascaridesasher}. 

To conclude, I argued that presuppositional and existential IPFVs are just a subcase of the canonical IPFV readings (ongoing process or potentially plural events), as suggested in \citet{gehrke22}. This proposal goes directly against \citeposst{gronn15} notion of ``fake'' IPFVs and furthermore makes the rather dubious Aspect Neutralisation Rules employed by \citet{gronnstechow10,gronnstechow16} and \citet{alvestad13} obsolete; this is a welcome result particularly for Czech, in which the optionality of the rule is rather problematic. Whether or not Czech has an IPFV use that is motivated by the existential factual context alone needs to be further explored, but judging from the data discussed in the literature there are no clear cases that would suggest this. However, both languages have presuppositional IPFVs, under the assumption that this is just a special case of the process reading of IPFVs. The empirical question that has to be further explored in this case is whether achievements (or necessarily unique events, as argued by \citealt{Mueller-Reichau2018a}) can appear in presuppositional IPFVs, in both languages. To my knowledge, achievements have been explicitly addressed in the literature only for existential but not for presuppositional contexts. 

The following and final section takes stock and zooms out again, focusing on the notion of temporal definiteness that played a role in all the accounts of cross-Slavic variation in aspect use discussed in this chapter. I outline a general research programme that draws parallels between the nominal and verbal domain with respect to definiteness.

\section{Taking stock: Temporal definiteness}
\label{sec:definiteness}

Let us take stock and review how different authors employ the notion of definiteness or specificity in the verbal and sentential domain to capture the semantics of aspect and cross-Slavic differences in aspect use. As already mentioned in \sectref{sec:defspec}, various definitions of definiteness and specificity play a role, as well as definite/specific events and times. \citeposst{Dickey2000} account uses the notion of temporal definiteness in the sense of uniqueness: a situation is uniquely locatable in a context (Russian PFV), it can be assigned to several points in time (Czech IPFV), or it is not assigned to a single, unique point in time (Russian IPFV). From a formal-semantic point of view, this could be translated as either the event time or (maybe more likely) the reference time being (non-)unique. 

\citet{klimek22}, in turn, proposes that PFVs in the North Slavic languages she investigates involve a definite or specific reference time point \citep[as proposed for Russian by][]{ramchandlingua}. She does not make precise what theory of definiteness she employs (uniqueness, familiarity, or other), and the notions ``definite'' and ``specific'' are used interchangeably; from the general discussion, I assume that familiarity might be closer to what she has in mind. To account for the cross-Slavic aspect variation, she proposes that in Czech and Western Polish, the PFV is used when the reference time is definite/specific with respect to the temporal trace of the event (which is what \citeauthor{ramchandlingua} assumed for Russian), whereas in Russian and Eastern Polish, the IPFV is used when the reference time is indefinite/non-specific with respect to the speech time. The former is a standard aspect relation (E-R), whereas the R-S relation is commonly taken to express tense \citep[e.g.][]{Reichenbach47,Klein1994}. This account, then, would suggest that Russian and Eastern Polish aspect is not a standard aspect but closer to a tense.\footnote{In \citet{gehrkeadv23}, I explore what it could mean that Russian aspect is more like a tense, and how the interaction of aspect with finiteness plays a crucial role in Russian, but not in Czech. A non-standard aspectual semantics in terms of the S-R relation, albeit different from the one outlined here, has also been proposed by \citet{Borik2006} for Russian.} 

\citet{Mueller-Reichau2018a,mueller23} also assumes definiteness for PFVs and locates the point of variation in additional conditions for PFVs that vary across the Slavic languages he is interested in. In particular, \citet{Mueller-Reichau2018a} argues that PFVs in Czech, Polish, and Russian involve reference to unique events; only Russian PFVs additionally require target state validity. \citet{mueller23}, in turn, does not employ uniqueness anymore but assumes that PFVs in Czech, Sorbian, Russian involve determinate event paths, so this again appears to be a switch from uniqueness to specificity. Additional requirements on Czech and Russian PFVs that are not directly related to definiteness are assumed to account for the cross-Slavic variation.

The accounts discussed so far do not agree on the kind of definiteness theory we need, and whether it is the event itself that is definite/specific or the reference time. This distinction is explicitly addressed in \citeposst{gronnstechow16} research programme, which I had outlined in \sectref{sec:gronn}. Drawing parallels to articleless languages in the way one could capture definiteness in the nominal domain in the absence of articles, they propose that also aspects and tenses do not directly correspond to (in)definite events and reference times,\footnote{This point is also already made in \citet{Klein1995}.} but only express a relational semantics. Whether we are dealing with an indefinite or definite event or reference time is determined by covert (in)definite operators, which are proposed to be dynamic generalised quantifiers that either introduce a new event or time (indefinite) or are anaphoric to a previous event or time (maximally presupposing given information). The exact mechanisms that lead to the insertion of a definite or indefinite operator are not fully worked out, but as a general research programme this looks interesting, since it opens up a whole area of parallels one can draw to the nominal domain. I will come back to this in a bit.

Embedded within this research programme is \citeposst{gronn15} account of Russian aspect. He proposes that existential IPFVs have the semantics of a ``fake'' IPFV (essentially a PFV semantics), which arises due to an indefinite event and an indefinite reference time; therefore, the reference time is too large for the PFV semantics to be informative and the ``fake'' IPFV is inserted instead. Presuppositional IPFVs, in turn, are proposed to involve a definite event and a definite reference time; the ``fake'' IPFV comes into play because narrative progression, commonly associated with the PFV (which involves an indefinite event), is to be avoided. \citet{gronnstechow16} explicitly state that whenever a ``semantically PFV aspect" (i.e. E is included in R) is definite/anaphoric, it gets morphologically neutralised to IPFV (recall their Aspect Neutralisation Rule from \sectref{sec:gronn}). 

Intuitively it makes sense to assume that with existential IPFVs both the event and the reference time are indefinite: they assert the existence of an event (recall the paraphrase `There is an event of this type'), and it is either not known when the event happened, suggesting that the reference time is indefinite and non-specific,\footnote{Recall from \sectref{sec:OF} that Russian scopally non-specific temporal adverbs of the \textit{-nibud'}-series, such as \textit{kogda-nibud'} even require the existential IPFV.} or the event could have happened more than once, suggesting that the reference time is non-unique and this could be captured as an anti-uniqueness implicature (recall discussion in \sectref{sec:defspec}). Intuitively, it also makes sense to think of the event with presuppositional IPFVs as definite: it is anaphorically linked to a previously introduced or contextually accommodatable discourse referent. It is less clear to me though why the reference time has to be definite as well%, especially if we assume that in discourse the reference time moves \citep[see also discussion in][]{Altshuler2012}
; I will come back to this below. 

\citeposst{alvestad13} account is also embedded in the kind of research programme outlined above, but she makes assumptions for Slavic that are sometimes different from \citeposst{gronn15}. Just like \citeauthor{gronn15}, she takes the PFV to involve an indefinite event; additionally she specifies that the reference time has to be unique%, and this is possibly similar to what \citet{ramchandlingua} and \citet{klimek22} have in mind
. While there is some discussion about the difference between uniqueness and anaphoricity, it is not made explicit enough why \citeauthor{alvestad13} uses uniqueness only for tense with PFVs, but anaphoricity everywhere else. At first sight, it might make intuitive sense: we have a unique (i.e. only one) reference time. However, this proposal faces the same issue outlined in the discussion of \citet{Mueller-Reichau2018a} in \sectref{sec:mueller}: If this is a general proposal for all Slavic languages, why would Czech ever use PFVs in habitual contexts (recall discussion in \sectref{sec:hab})? Also like \citeauthor{gronn15}, \citeauthor{alvestad13} assumes that with existential IPFVs both the event and the reference time are indefinite. Where her account differs is with respect to presuppositional IPFVs: while the event is definite-anaphoric as well, she now assumes that the reference time can be either definite-anaphoric or indefinite.

\citeposst{gehrketrueipf} proposal for presuppositional IPFVs outlined in \REF{ex:zaplatilib} in the previous section does not explicitly discuss definite events and times, but the account itself makes clear that both start out as indefinite: each (finite) verb form introduces a new event and a new reference time. Due to the information-structural cues, the newly introduced event is anaphorically linked to a previously introduced event ($e_2 = e_1$), and this brings about the impression that we are dealing with a definite. The reference time stays indefinite, but because of the anaphoric link between the two events, it indirectly has to be part of the previously introduced event's reference time. Furthermore, there is no ``fake'' IPFV \citep[][]{gehrke22,gehrketrueipf}: Presuppositional IPFVs involve an IPFV semantics (R is included in E) and a certain information-structurally marked discourse, whereas  existential IPFVs involve potential repeatability, and plural events in Russian (but not in Czech) require the IPFV. Obviously, this proposals calls for close exploration of the role of discourse in a dynamic framework \citep[see also][]{Altshuler2012}, but we can dispense with any type of aspect neutralisation rule (as those in \sectref{sec:gronn} and \sectref{sec:alvestad})%; this is a welcome result particularly for Czech, in which the optionality of the rule is rather problematic
.

Let us go back to the point of departure for \citet{gronnstechow16}, then, namely the idea that we can draw parallels to the way (in)definite interpretations arise in the absence of articles in articleless languages. While there are accounts that employ covert (in)definite operators or covert type shifts of the relevant kind also for these languages \citep[e.g.][]{chierchia98}, there are recent proposals that question whether bare nominals in articleless languages should ever (semantically) be analysed as definites, even if we get the impression that contextually they are definite. For example, \citet{heim11} proposes that in articleless languages all bare nominals (in argument position) are existential, i.e. indefinite. They are also acceptable in definite contexts because these languages lack a definite article that could block them (i.e. there is no competition with definites and thus no anti-uniqueness implicature; recall discussion in \sectref{sec:defspec}). 

\citeposst{heim11} line of reasoning is picked up in recent work on the semantics of bare nominals in Russian \citep{simikdemian20, %borik+20, 
seresborik21} (see also discussion in \citetv{chapters/02_borik}). \citeauthor{seresborik21} show that Russian bare nominals can appear fairly freely in both definite and indefinite contexts, and other factors, such as topicality and the overall context only indirectly contribute to the impression that we get a definite interpretation. As \citet[][340]{seresborik21} put it: ``The perceived definiteness in Russian is analysed as a pragmatic effect (not as a result of a covert type-shift), which has the following sources: ontological uniqueness, topicality, and familiarity/anaphoricity." \citeauthor{simikdemian20}, in turn, show experimentally that factors that have been described to correlate with a definite (in the sense of unique/maximal) interpretation of Russian bare nominals (word order, prosody, number) are rather weak, and that the nominals in such contexts do not behave like German definites. They conclude that their data rather support \citeposst{heim11} proposal that Russian bare nominals are always indefinite.  

So what if we end up with the following, which could serve as a research programme for future investigations. Aspects, tenses, VPs, NPs are predicates (following \citealt{gronnstechow16} for events/times and \citealt{coppockbeaver15} for nominals); additional information about them (e.g. provided by adverbials) are added via predicate modification. At the relevant syntactic positions we get existential closure over the respective variables. In the absence of overt determiners, all events and times are indefinite, just like bare nominals \citep[following][]{heim11,simikdemian20,seresborik21}; the impression of a definite interpretation is only due to the context (including information structural cues) but there is no iota (or similar) shift. Note, again, that the idea that events and times are always indefinite is already implicit in \citeposst{gehrketrueipf} account of presuppositional IPFVs. 

If we want to fully exploit parallels to definiteness in the nominal domain when building a theory of aspect (or tense), we will have to exploit the full spectrum: different types of both definites (unique, familiar, anaphoric, weak vs. strong, kind reference, etc.) and indefinites, including a discussion of specificity, free choice, ignorance implicatures, and so on. Furthermore, we need to take into account the contribution of adverbials that further specify the event and the reference time, information structure and context more generally. For example, the different Russian indefinite series briefly mentioned in \sectref{sec:defspec} (e.g. \textit{kto-nibud'/to} `someone')  also exist in the temporal domain (e.g. \textit{kogda-nibud'/to} `sometime'), which provides additional information about the reference time (epistemically or scopally (non-)specific, and similar). In the nominal domain, \citet{geist10a} shows that the \textit{-nibud'}-series is always scopally non-specific and needs to be licensed, e.g. by a quantifier or a modal operator (see also \citetv{chapters/06_geist}). Similarly, \citet{pereltsvaignibud} claims that the \textit{-nibud'}-series needs to covary within the scope of an operator or quantifier. In the verbal domain, we saw in \sectref{sec:OF} that \textit{kogda-nibud'} requires the IPFV; we are certainly not dealing with a single time/event, in fact, the reference time is non-referential in this case. Does that mean that there is covariation with a silent generic or other operator?

Czech (as well as other Slavic languages) has comparable indefinite series, though the forms and functions differ. For example, \citet{aguilar+10} observe that the \textit{koli}-series primarily expresses free choice \citep[see also][]{aloni22}, the \textit{si}-series epistemic indefiniteness; \textit{ně}-prefixed \textit{wh}-words are treated as plain indefinites, e.g. \textit{někdo} (based on `who') `someone', \textit{někdy} (based on `when') `sometime'. Recall from the discussion of \REF{ex:kogda} in \sectref{sec:OFCZRU}, fn. \ref{fn:kogdadoubt}, that a Czech native speaker doubted the preference of the Czech PFV in the context of the ``universal \textit{-koli(v)}-words'' (the example contained \textit{kdykoliv} `ever') and would use the ``existential'' \textit{někdy} instead. So if the reference time is a free choice indefinite, this leads to a preference for the Czech IPFV (for this speaker, but maybe not for the informants of \citealt{klimek22}), but when it is a plain existential indefinite, the PFV is the first choice (in that particular example)?

Furthermore, \citet[][25]{klimek22} discusses the Russian corpus example in \REF{ex:libo} as potentially problematic for both \citeposst{Mueller-Reichau2018a} account of PFVs as involving unique events, as well as \citeposst{gehrketrueipf} account of existential IPFVs being motivated by potential repeatability: 

\ea\label{ex:libo}
\gll Ėto byla vešč' lučšaja iz vsex veščej, kotorye ja kogda-libo sozdal.\\
this was thing best of all things which I when-\textsc{libo} created.\textsc{pfv}\\
\glt `It was the best thing of all the things I had ever created.' \hfill (Russian)
\z 

\noindent At first sight, this indeed seems to be a problem for both accounts. According to Daria Seres (p.c.), \REF{ex:libo} is acceptable, but we cannot replace \textit{kogda-libo} by \textit{kogda-nibud'}, the temporal adverb that is usually found in discussions of Russian existential IPFVs. So we have to explore what the difference is between \textit{kogda-libo} and \textit{kogda-nibud'}, both of which are regularly translated as `ever'. More generally, according to \citeposst{Haspelmath1997} semantic-map approach to indefinite pronouns, the \textit{-nibud'}-series occurs in the contexts irrealis, question, and conditional antecedent; these are a proper subset of the functions of \textit{-libo}, which additionally appears in indirect negation and comparatives. In \REF{ex:libo}, we are dealing with a comparative context, so whatever the semantics here is makes it different enough from the contexts which require the Russian IPFV, and the PFV can be used. 

Setting aside stylistic differences between \textit{-nibud'} and \textit{-libo}, which are also commonly noted in the literature, we see for example that the native speakers that \citet{ward77} consulted could always replace \textit{-nibud'} by \textit{-libo}, but not always the other way around% \citep[see also][]{veyrenc64}
. \citeauthor{ward77} proposes that \textit{-libo} presupposes or implies the existence of a set: ``There exists or can exist a set of $x$'s but it is not asserted that there exists a particular member of that set such that that member can or does participate in the event" \citep[][465]{ward77}. In contrast, \textit{-nibud'} ``leaves the existence of a set unmarked'' \citep[][467]{ward77}. So it seems that the existence presupposition of \textit{-libo} is compatible with the Russian PFV but its absence is not. These and similar considerations, which so far only scratch the surface, would have to be explored further to get at the full picture.
 
Finally, let us return to \citeposst{Dickey2000} proposal that the Russian PFV expresses temporal definiteness, whereas the Czech PFV expresses totality. Traditionally, it is common to view the semantics of PFV in (many) Slavic languages in terms of totality; for example \citet[][]{Filip1999} argues that the ``Slavic'' PFV maps events of any type to total events, to represent them ``as integrated wholes (i.e, in their totality, as single indivisible wholes)'' \citep[][184]{Filip1999}. In \citet{Filip2008}, she argues, again for Slavic in general, that PFVs involve a maximalisation operator on events, and a similar idea is formalised in \citet{altshuler14} who proposes that PFV operators cross-linguistically require ``a maximal stage of an event in the extension of the VP that it combines with'' \citep[][762]{altshuler14}. So what if this general proposal for PFVs cross-linguistically does not hold for Russian but only for Czech? And what if, as some of the accounts discussed here, the Russian PFV is truly definite, comparable to definite articles in the nominal domain? More generally, it is possible that certain tenses and aspects (in some languages) come with uniqueness (or other) presuppositions, leading to competition with other aspects/tenses, similar to the competition that \citet{heim11} and others assume for the nominal domain (see e.g. \citealt{zhaodiss} for recent work on the competition between perfect and past in various languages). This is where various approaches discussed in this chapter could sneak back in and be refined accordingly.

\section*{Acknowledgments}
\largerpage

Thanks to Petra Charvátová, Denisa Lenertová, Barbora Schnelle, and Radek Šimík for help with the Czech data, to Petr Biskup for invaluable feedback on the first draft of this chapter, to two anonymous reviewers, and to Hans Robert Mehlig for wonderful discussion of all things related to aspect over the years.


\section*{Abbreviations}

\begin{tabularx}{.5\textwidth}{@{}lQ}
\textsc{1}&first person\\
\textsc{2}&second person\\
\textsc{3}&third person\\
\textsc{acc}&accusative\\
\textsc{adj}&adjective\\
\textsc{adv}&adverb\\
\textsc{ap}&adverbial participle\\
\textsc{aux}&auxiliary\\
CUS&Colloquial Upper Sorbian\\
\textsc{dat}&dative\\
\textsc{dem}&demonstrative\\
E&event time\\
\textsc{fem}&feminine\\
\textsc{freq}&frequentative\\
\textsc{gen}&genitive\\
\textsc{imp}&imperative\\
\textsc{indet}&indterminate\\
\textsc{inf}&infinitive\\
\textsc{instr}&instrumental\\
\textsc{ipfv}&imperfective\\
\end{tabularx}%
\begin{tabularx}{.5\textwidth}{lQ@{}}
\textsc{lf}&long form\\
\textsc{m}&masculine\\
\textsc{n}&neuter\\
\textsc{neg}&negation\\
\textsc{nom}&nominative\\
\textsc{pap}&past active participle\\
\textsc{pl}&plural\\
\textsc{pfv}&perfective\\
\textsc{po}&delimitative \textit{po-}\\
\textsc{possrefl}&possessive reflexive\\
\textsc{ppp}&past passive participle\\
\textsc{prs}&present tense\\
R&reference time\\
\textsc{refl}&reflexive\\
\textsc{sg}&singular\\
\textsc{si}&secondary IPFV\\
S&speech time\\
SOE&sequence of events\\
\textsc{voc}&vocative\\
\textsc{za}&ingressive \textit{za-}\\
%&\\ % this dummy row achieves correct vertical alignment of both tables
\end{tabularx}

\printbibliography[heading=subbibliography,notkeyword=this]

\end{document}
